% Generated mechanically from the frozen derham.tex target.
% Do not edit this generated file; edit the bound locale source instead.

% OK, start here
%

\setcounter{chapter}{49}
\chapter{德拉姆上同调}



\phantomsection
\label{derham-section-phantom}


\section{引言}
\label{derham-section-introduction}

\noindent
本章先讨论概形态射的德拉姆复形，最后证明：当基域的特征为零时，
德拉姆上同调定义了一种韦伊上同调理论。




\section{德拉姆复形}
\label{derham-section-de-rham-complex}

\noindent
设 $p : X \to S$ 是概形态射。存在复形
$$
\Omega^\bullet_{X/S} =
\mathcal{O}_{X/S} \to \Omega^1_{X/S} \to \Omega^2_{X/S} \to \ldots
$$
它是 $p^{-1}\mathcal{O}_S$-模复形，其中
$\Omega^i_{X/S} = \wedge^i(\Omega_{X/S})$
置于次数 $i$，其微分由局部截面上的法则
$\text{d}(g_0 \text{d}g_1 \wedge \ldots \wedge \text{d}g_p) =
\text{d}g_0 \wedge \text{d}g_1 \wedge \ldots \wedge \text{d}g_p$
决定。参见《模》，第 \ref{modules-section-de-rham-complex} 节。

\medskip\noindent
给定概形的交换图
$$
\xymatrix{
X' \ar[r]_f \ar[d] & X \ar[d] \\
S' \ar[r] & S
}
$$
则存在复形的典范映射
$f^{-1}\Omega_{X/S}^\bullet \to \Omega^\bullet_{X'/S'}$ 以及
$\Omega_{X/S}^\bullet \to f_*\Omega^\bullet_{X'/S'}$。
参见《模》，第 \ref{modules-section-de-rham-complex} 节。
将其线性化后，对每个 $p$ 都得到线性映射
$f^*\Omega^p_{X/S} \to \Omega^p_{X'/S'}$。

\medskip\noindent
特别地，若 $f : Y \to X$ 是基概形 $S$ 上的概形态射，则有复形映射
$$
\Omega^\bullet_{X/S} \longrightarrow f_*\Omega^\bullet_{Y/S}
$$
将其线性化可知，对每个 $p \geq 0$ 都得到典范映射
$$
\Omega^p_{X/S} \otimes_{\mathcal{O}_X} f_*\mathcal{O}_Y
\longrightarrow
f_*\Omega^p_{Y/S}
$$

\begin{lemma}
\label{derham-lemma-base-change-de-rham}
设
$$
\xymatrix{
X' \ar[r]_f \ar[d] & X \ar[d] \\
S' \ar[r] & S
}
$$
为概形的笛卡儿图。则上述映射诱导同构
$f^*\Omega^p_{X/S} \to \Omega^p_{X'/S'}$。
\end{lemma}

\begin{proof}
将《态射》，引理 \ref{morphisms-lemma-base-change-differentials}
与外幂构造和基变换相容这一事实结合即可。
\end{proof}

\begin{lemma}
\label{derham-lemma-etale}
考虑概形的交换图
$$
\xymatrix{
X' \ar[r]_f \ar[d] & X \ar[d] \\
S' \ar[r] & S
}
$$
若 $X' \to X$ 与 $S' \to S$ 均为平展态射，则上述映射诱导同构
$f^*\Omega^p_{X/S} \to \Omega^p_{X'/S'}$。
\end{lemma}

\begin{proof}
我们有 $\Omega_{S'/S} = 0$ 和 $\Omega_{X'/X} = 0$；例如参见
《态射》，引理 \ref{morphisms-lemma-etale-at-point}。于是，由《态射》中的引理
\ref{morphisms-lemma-triangle-differentials} 与
\ref{morphisms-lemma-triangle-differentials-smooth}
所给出的短正合列可知
$\Omega_{X'/S'} = \Omega_{X'/S} = f^*\Omega_{X/S}$。
取外幂即得结论。
\end{proof}






\section{德拉姆上同调}
\label{derham-section-de-rham-cohomology}

\noindent
设 $p : X \to S$ 是概形态射。我们把
{\it $X$ 在 $S$ 上的德拉姆上同调}定义为上同调群
$$
H^i_{dR}(X/S) = H^i(R\Gamma(X, \Omega^\bullet_{X/S}))
$$
由于 $\Omega^\bullet_{X/S}$ 是 $p^{-1}\mathcal{O}_S$-模复形，
这些上同调群自然是 $H^0(S, \mathcal{O}_S)$-模。

\medskip\noindent
给定概形的交换图
$$
\xymatrix{
X' \ar[r]_f \ar[d] & X \ar[d] \\
S' \ar[r] & S
}
$$
利用第 \ref{derham-section-de-rham-complex} 节中的典范映射，我们得到拉回映射
$$
f^* :
R\Gamma(X, \Omega^\bullet_{X/S})
\longrightarrow
R\Gamma(X', \Omega^\bullet_{X'/S'})
$$
以及
$$
f^* : H^i_{dR}(X/S) \longrightarrow H^i_{dR}(X'/S')
$$
这些拉回满足显然的复合法则。
特别地，若固定基概形 $S$，则德拉姆上同调是 $S$ 上概形范畴的反变函子。

\begin{lemma}
\label{derham-lemma-de-rham-affine}
设 $X \to S$ 是由环同态 $R \to A$ 给出的仿射概形态射。则
$R\Gamma(X, \Omega^\bullet_{X/S}) = \Omega^\bullet_{A/R}$ 在 $D(R)$ 中成立，且
$H^i_{dR}(X/S) = H^i(\Omega^\bullet_{A/R})$。
\end{lemma}

\begin{proof}
这由《概形上同调》，引理
\ref{coherent-lemma-quasi-coherent-affine-cohomology-zero}
和勒雷无环引理
（《导出范畴》，引理 \ref{derived-lemma-leray-acyclicity}）推出。
\end{proof}

\begin{lemma}
\label{derham-lemma-quasi-coherence-relative}
设 $p : X \to S$ 是概形态射。若 $p$ 拟紧且拟分离，则
$Rp_*\Omega^\bullet_{X/S}$ 是 $D_\QCoh(\mathcal{O}_S)$ 中的对象。
\end{lemma}

\begin{proof}
存在一个第一叶为 $E_1^{a, b} = R^bp_*\Omega^a_{X/S}$、
收敛到 $Rp_*\Omega^\bullet_{X/S}$ 的上同调的谱序列
（参见《导出范畴》，引理 \ref{derived-lemma-two-ss-complex-functor}）。
因此，由《同调》，引理 \ref{homology-lemma-first-quadrant-ss}，
只需证明 $R^bp_*\Omega^a_{X/S}$ 拟凝聚。
这由《概形上同调》，引理
\ref{coherent-lemma-quasi-coherence-higher-direct-images} 推出。
\end{proof}

\begin{lemma}
\label{derham-lemma-coherence-relative}
设 $p : X \to S$ 是真概形态射，且 $S$ 局部诺特。则
$Rp_*\Omega^\bullet_{X/S}$ 是 $D_{\textit{Coh}}(\mathcal{O}_S)$ 中的对象。
\end{lemma}

\begin{proof}
在此情形下，由《态射》，引理 \ref{morphisms-lemma-finite-type-differentials}，
模 $\Omega^i_{X/S}$ 是凝聚的。因此，可以逐字采用引理
\ref{derham-lemma-quasi-coherence-relative} 证明中的论证，并使用《概形上同调》，命题
\ref{coherent-proposition-proper-pushforward-coherent}。
\end{proof}

\begin{lemma}
\label{derham-lemma-finite-de-Rham}
设 $A$ 是诺特环，并设 $X$ 是 $S = \Spec(A)$ 上的真概形。
则 $H^i_{dR}(X/S)$ 是有限 $A$-模（对所有 $i$）。
\end{lemma}

\begin{proof}
这是引理 \ref{derham-lemma-coherence-relative} 的特例。
\end{proof}

\begin{lemma}
\label{derham-lemma-proper-smooth-de-Rham}
设 $f : X \to S$ 是真光滑概形态射。则
$Rf_*\Omega^p_{X/S}$（$p \geq 0$）和 $Rf_*\Omega^\bullet_{X/S}$
都是 $D(\mathcal{O}_S)$ 中的完美对象，且其构造与任意基变换相容。
\end{lemma}

\begin{proof}
由于 $f$ 光滑，模 $\Omega^p_{X/S}$ 是有限局部自由
$\mathcal{O}_X$-模；参见《态射》，引理
\ref{morphisms-lemma-smooth-omega-finite-locally-free}。
由引理 \ref{derham-lemma-base-change-de-rham}，其构造与任意基变换相容。因此，
$Rf_*\Omega^p_{X/S}$ 是 $D(\mathcal{O}_S)$ 中的完美对象，
且其构造与任意基变换相容；参见《概形的导出范畴》，引理
\ref{perfect-lemma-flat-proper-perfect-direct-image-general}。
这证明了本引理的第一个断言。

\medskip\noindent
为了证明 $Rf_*\Omega^\bullet_{X/S}$ 在 $S$ 上完美，可以在 $S$ 上局部地论证。
因此可设 $S$ 拟紧。这意味着可以假设 $\Omega^n_{X/S}$ 在 $n$ 充分大时为零。
对每个 $p \geq 0$，我们断言
$Rf_*\sigma_{\geq p}\Omega^\bullet_{X/S}$ 是
$D(\mathcal{O}_S)$ 中的完美对象，且其构造与任意基变换相容。
由上述论证可知，当 $p \gg 0$ 时断言成立。假设该断言对 $p$ 成立，
并考虑可区别三角
$$
\sigma_{\geq p}\Omega^\bullet_{X/S} \to
\sigma_{\geq p - 1}\Omega^\bullet_{X/S} \to
\Omega^{p - 1}_{X/S}[-(p - 1)] \to
(\sigma_{\geq p}\Omega^\bullet_{X/S})[1]
$$
它位于 $D(f^{-1}\mathcal{O}_S)$ 中。
应用正合函子 $Rf_*$，在 $D(\mathcal{O}_S)$ 中得到一个可区别三角。
由于完美性具有三取二性质
（《上同调》，引理 \ref{cohomology-lemma-two-out-of-three-perfect}），
可得 $Rf_*\sigma_{\geq p - 1}\Omega^\bullet_{X/S}$ 是
$D(\mathcal{O}_S)$ 中的完美对象。其构造与任意基变换相容的断言同理。
\end{proof}





\section{杯积}
\label{derham-section-cup-product}

\noindent
考虑映射
$\Omega^p_{X/S} \times \Omega^q_{X/S} \to \Omega^{p + q}_{X/S}$
其定义为 $(\omega , \eta) \longmapsto \omega \wedge \eta$。
利用第 \ref{derham-section-de-rham-complex} 节给出的 $\text{d}$ 的公式，
以及 $\text{d} : \mathcal{O}_X \to \Omega_{X/S}$ 的莱布尼茨法则，可知
$\text{d}(\omega \wedge \eta) = \text{d}(\omega) \wedge \eta +
(-1)^{\deg(\omega)} \omega \wedge \text{d}(\eta)$。这说明
$\wedge$ 定义了态射
\begin{equation}
\label{derham-equation-wedge}
\wedge :
\text{Tot}(
\Omega^\bullet_{X/S} \otimes_{p^{-1}\mathcal{O}_S} \Omega^\bullet_{X/S})
\longrightarrow
\Omega^\bullet_{X/S}
\end{equation}
它是 $p^{-1}\mathcal{O}_S$-模复形的态射。

\medskip\noindent
把《上同调》，第 \ref{cohomology-section-cup-product} 节中的杯积
与（\ref{derham-equation-wedge}）结合，得到一个
$H^0(S, \mathcal{O}_S)$-双线性杯积映射
$$
\cup : H^i_{dR}(X/S) \times H^j_{dR}(X/S) \longrightarrow H^{i + j}_{dR}(X/S)
$$
例如，若 $\omega \in \Gamma(X, \Omega^i_{X/S})$ 和
$\eta \in \Gamma(X, \Omega^j_{X/S})$ 都是闭的，则
$\omega$ 与 $\eta$ 的德拉姆上同调类之杯积，就是 $\omega \wedge \eta$ 的
德拉姆上同调类；参见《上同调》，第 \ref{cohomology-section-cup-product} 节中的讨论。

\medskip\noindent
给定概形的交换图
$$
\xymatrix{
X' \ar[r]_f \ar[d] & X \ar[d] \\
S' \ar[r] & S
}
$$
则拉回映射
$f^* : R\Gamma(X, \Omega^\bullet_{X/S}) \to R\Gamma(X', \Omega^\bullet_{X'/S'})$
以及
$f^* : H^i_{dR}(X/S) \longrightarrow H^i_{dR}(X'/S')$
与上述杯积相容。

\begin{lemma}
\label{derham-lemma-cup-product-graded-commutative}
设 $p : X \to S$ 是概形态射。
$H^*_{dR}(X/S)$ 上的杯积满足结合律和分次交换律。
\end{lemma}

\begin{proof}
这由《上同调》中的引理 \ref{cohomology-lemma-cup-product-associative} 与
\ref{cohomology-lemma-cup-product-commutative}
以及 $\wedge$ 满足结合律和分次交换律这一事实推出。
\end{proof}

\begin{remark}
\label{derham-remark-relative-cup-product}
设 $p : X \to S$ 是概形态射。可以把
$\Omega^\bullet_{X/S}$ 看作微分分次 $p^{-1}\mathcal{O}_S$-代数层；参见
《微分分次层》，定义 \ref{sdga-definition-dga}。
特别地，《微分分次层》，第 \ref{sdga-section-misc} 节中的讨论适用。
例如，这意味着对任意交换图
$$
\xymatrix{
X \ar[d]_p \ar[r]_f & Y \ar[d]^q \\
S \ar[r]^h & T
}
$$
都存在典范的相对杯积
$$
\mu :
Rf_*\Omega^\bullet_{X/S}
\otimes_{q^{-1}\mathcal{O}_T}^\mathbf{L}
Rf_*\Omega^\bullet_{X/S}
\longrightarrow
Rf_*\Omega^\bullet_{X/S}
$$
它位于 $D(Y, q^{-1}\mathcal{O}_T)$ 中，满足结合律，并且在上同调上给出上述杯积。
\end{remark}

\begin{remark}
\label{derham-remark-cup-product-as-a-map}
设 $f : X \to S$ 是概形态射，并设 $\xi \in H_{dR}^n(X/S)$。
根据《微分分次层》，第 \ref{sdga-section-misc} 节中的讨论，存在典范态射
$$
\xi' : \Omega^\bullet_{X/S} \to \Omega^\bullet_{X/S}[n]
$$
它位于 $D(f^{-1}\mathcal{O}_S)$ 中，并由下列性质中的 (1) 和 (2) 唯一刻画：
\begin{enumerate}
\item $\xi'$ 可提升为右微分分次 $\Omega^\bullet_{X/S}$-模的导出范畴中的态射，并且
\item $\xi'(1) = \xi$ 在
$H^0(X, \Omega^\bullet_{X/S}[n]) = H^n_{dR}(X/S)$ 中成立；
\item 态射 $\xi'$ 把 $\eta \in H^m_{dR}(X/S)$ 送到
$\xi \cup \eta$；这是 $H^{n + m}_{dR}(X/S)$ 中的元素；
\item $\xi'$ 的构造与到开集的限制相容：对开集 $U \subset X$，
限制 $\xi'|_U$ 是与像 $\xi|_U \in H^n_{dR}(U/S)$ 对应的态射；
\item 对备注 \ref{derham-remark-relative-cup-product} 中的任意图，得到交换图
$$
\xymatrix{
Rf_*\Omega^\bullet_{X/S}
\otimes_{q^{-1}\mathcal{O}_T}^\mathbf{L}
Rf_*\Omega^\bullet_{X/S} \ar[d]_{\xi' \otimes \text{id}}
\ar[r]_-\mu &
Rf_*\Omega^\bullet_{X/S} \ar[d]^{\xi'} \\
Rf_*\Omega^\bullet_{X/S}[n]
\otimes_{q^{-1}\mathcal{O}_T}^\mathbf{L}
Rf_*\Omega^\bullet_{X/S}
\ar[r]^-\mu &
Rf_*\Omega^\bullet_{X/S}[n]
}
$$
它位于 $D(Y, q^{-1}\mathcal{O}_T)$ 中。
\end{enumerate}
\end{remark}




\section{霍奇上同调}
\label{derham-section-hodge-cohomology}

\noindent
设 $p : X \to S$ 是概形态射。我们把
{\it $X$ 在 $S$ 上的霍奇上同调}定义为上同调群
$$
H^n_{Hodge}(X/S) = \bigoplus\nolimits_{n = p + q} H^q(X, \Omega^p_{X/S})
$$
并将其视为分次 $H^0(X, \mathcal{O}_X)$-模。把微分形式的楔积
与《上同调》，第 \ref{cohomology-section-cup-product} 节中的杯积结合，
便定义了一个 $H^0(X, \mathcal{O}_X)$-双线性杯积
$$
\cup :
H^i_{Hodge}(X/S) \times H^j_{Hodge}(X/S)
\longrightarrow
H^{i + j}_{Hodge}(X/S)
$$
当然，若 $\xi \in H^q(X, \Omega^p_{X/S})$ 且
$\xi' \in H^{q'}(X, \Omega^{p'}_{X/S})$，则 $\xi \cup \xi' \in
H^{q + q'}(X, \Omega^{p + p'}_{X/S})$。

\begin{lemma}
\label{derham-lemma-cup-product-hodge-graded-commutative}
设 $p : X \to S$ 是概形态射。
$H^*_{Hodge}(X/S)$ 上的杯积满足结合律和分次交换律。
\end{lemma}

\begin{proof}
其证明与引理 \ref{derham-lemma-cup-product-graded-commutative} 的证明完全相同。
\end{proof}

\noindent
给定概形的交换图
$$
\xymatrix{
X' \ar[r]_f \ar[d] & X \ar[d] \\
S' \ar[r] & S
}
$$
则有拉回映射
$f^* : H^i_{Hodge}(X/S) \longrightarrow H^i_{Hodge}(X'/S')$
它与分次以及上述杯积相容。







\section{两个谱序列}
\label{derham-section-hodge-to-de-rham}

\noindent
设 $p : X \to S$ 是概形态射。由于 $p^{-1}\mathcal{O}_S$-模在 $X$ 上的范畴
有足够多的内射对象，
$\Omega^\bullet_{X/S}$ 存在嘉当--艾伦伯格解消。
参见《导出范畴》，引理 \ref{derived-lemma-cartan-eilenberg}。
因此可以应用《导出范畴》，引理
\ref{derived-lemma-two-ss-complex-functor}，得到两个都收敛到 $X$ 在 $S$
上之德拉姆上同调的谱序列。

\medskip\noindent
第一个通常称为{\it 霍奇--德拉姆谱序列}。该谱序列的第一叶为
$$
E_1^{p, q} = H^q(X, \Omega^p_{X/S})
$$
这些就是 $X/S$ 的霍奇上同调群（名称由此而来）。这一叶上的微分
$d_1$ 由映射
$d_1^{p, q} : H^q(X, \Omega^p_{X/S}) \to H^q(X, \Omega^{p + 1}_{X/S})$
给出；这些映射由微分
$\text{d} : \Omega^p_{X/S} \to \Omega^{p + 1}_{X/S}$ 诱导。图示如下：
$$
\xymatrix{
H^2(X, \mathcal{O}_X) \ar[r] \ar@{-->}[rrd] \ar@{..>}[rrrdd] &
H^2(X, \Omega^1_{X/S}) \ar[r] \ar@{-->}[rrd] &
H^2(X, \Omega^2_{X/S}) \ar[r] &
H^2(X, \Omega^3_{X/S}) \\
H^1(X, \mathcal{O}_X) \ar[r] \ar@{-->}[rrd] &
H^1(X, \Omega^1_{X/S}) \ar[r] \ar@{-->}[rrd] &
H^1(X, \Omega^2_{X/S}) \ar[r] &
H^1(X, \Omega^3_{X/S}) \\
H^0(X, \mathcal{O}_X) \ar[r] &
H^0(X, \Omega^1_{X/S}) \ar[r] &
H^0(X, \Omega^2_{X/S}) \ar[r] &
H^0(X, \Omega^3_{X/S})
}
$$
其中用虚线箭头标出 $E_2$ 叶上微分的源与靶，并用点线箭头表示
$E_3$ 叶上的一个微分。考察次数 $0$，得到
$$
H^0_{dR}(X/S) =
\Ker(\text{d} : H^0(X, \mathcal{O}_X) \to H^0(X, \Omega^1_{X/S}))
$$
当然，从德拉姆复形的次数 $0$ 项起始为
$\mathcal{O}_X \to \Omega^1_{X/S}$ 这一事实也能立刻看出这一点。

\medskip\noindent
第二个谱序列通常称为{\it 共轭谱序列}。该谱序列的第二叶为
$$
E_2^{p, q} = H^p(X, H^q(\Omega^\bullet_{X/S})) = H^p(X, \mathcal{H}^q)
$$
其中 $\mathcal{H}^q = H^q(\Omega^\bullet_{X/S})$ 是德拉姆复形的第 $q$
个上同调层；这里的德拉姆复形是 $X/S$ 的德拉姆复形。
这一叶上的微分由 $E_2^{p, q} \to E_2^{p + 2, q - 1}$ 给出。
图示如下：
$$
\xymatrix{
H^0(X, \mathcal{H}^2) \ar[rrd] \ar@{..>}[rrrdd] &
H^1(X, \mathcal{H}^2) \ar[rrd] &
H^2(X, \mathcal{H}^2) &
H^3(X, \mathcal{H}^2) \\
H^0(X, \mathcal{H}^1) \ar[rrd] &
H^1(X, \mathcal{H}^1) \ar[rrd] &
H^2(X, \mathcal{H}^1) &
H^3(X, \mathcal{H}^1) \\
H^0(X, \mathcal{H}^0) &
H^1(X, \mathcal{H}^0) &
H^2(X, \mathcal{H}^0) &
H^3(X, \mathcal{H}^0)
}
$$
考察次数 $0$，得到
$$
H^0_{dR}(X/S) = H^0(X, \mathcal{H}^0)
$$
细想之下，这是显然的。在次数 $1$，得到正合列
$$
0 \to H^1(X, \mathcal{H}^0) \to H^1_{dR}(X/S) \to
H^0(X, \mathcal{H}^1) \to H^2(X, \mathcal{H}^0) \to H^2_{dR}(X/S)
$$
事实表明，若 $X \to S$ 光滑且 $S$ 处于特征 $p$，则层
$\mathcal{H}^q$ 可以借助某个微分形式层来计算，而共轭谱序列是一个很有价值的
工具（在此插入将来的参考文献）。



\section{霍奇滤过}
\label{derham-section-hodge-filtration}

\noindent
设 $X \to S$ 是概形态射。$H^n_{dR}(X/S)$ 上的霍奇滤过，是由
霍奇--德拉姆谱序列诱导的滤过（《同调》，定义
\ref{homology-definition-filtration-cohomology-filtered-complex}）。
为避免误解，下面将它明确地定义出来。

\begin{definition}
\label{derham-definition-hodge-filtration}
设 $X \to S$ 是概形态射。$H^n_{dR}(X/S)$ 上的{\it 霍奇滤过}，
是由下列各项组成的滤过：
$$
F^pH^n_{dR}(X/S) = \Im\left(H^n(X, \sigma_{\geq p}\Omega^\bullet_{X/S})
\longrightarrow H^n_{dR}(X/S)\right)
$$
其中 $\sigma_{\geq p}\Omega^\bullet_{X/S}$ 如《同调》，
第 \ref{homology-section-truncations} 节所述。
\end{definition}

\noindent
当然，$\sigma_{\geq p}\Omega^\bullet_{X/S}$ 是相对德拉姆复形的子复形，
于是得到相对德拉姆复形的滤过
$$
\Omega^\bullet_{X/S} = \sigma_{\geq 0}\Omega^\bullet_{X/S} \supset
\sigma_{\geq 1}\Omega^\bullet_{X/S} \supset
\sigma_{\geq 2}\Omega^\bullet_{X/S} \supset
\sigma_{\geq 3}\Omega^\bullet_{X/S} \supset \ldots
$$
其满足
$\text{gr}^p(\Omega^\bullet_{X/S}) = \Omega^p_{X/S}[-p]$。
把 $\Omega^\bullet_{X/S}$ 视为滤过的层复形，则《上同调》，引理
\ref{cohomology-lemma-spectral-sequence-filtered-object} 所构造的谱序列，
与第 \ref{derham-section-hodge-to-de-rham} 节中借助《上同调》，例
\ref{cohomology-example-spectral-sequence-bis} 构造的霍奇--德拉姆谱序列相同。
此外，楔积（\ref{derham-equation-wedge}）把
$\text{Tot}(\sigma_{\geq i}\Omega^\bullet_{X/S} \otimes_{p^{-1}\mathcal{O}_S}
\sigma_{\geq j}\Omega^\bullet_{X/S})$ 送入
$\sigma_{\geq i + j}\Omega^\bullet_{X/S}$。因此得到交换图
$$
\xymatrix{
H^n(X, \sigma_{\geq i}\Omega^\bullet_{X/S}))
\times 
H^m(X, \sigma_{\geq j}\Omega^\bullet_{X/S}))
\ar[r] \ar[d] &
H^{n + m}(X, \sigma_{\geq i + j}\Omega^\bullet_{X/S})) \ar[d] \\
H^n_{dR}(X/S) \times
H^m_{dR}(X/S)
\ar[r]^\cup &
H^{n + m}_{dR}(X/S)
}
$$
特别地，可以得到
$$
F^iH^n_{dR}(X/S) \cup F^jH^m_{dR}(X/S) \subset F^{i + j}H^{n + m}_{dR}(X/S)
$$






\section{K\"unneth 公式}
\label{derham-section-kunneth}

\noindent
德拉姆上同调的一个重要特征是它具有 K\"unneth 公式。

\medskip\noindent
设 $a : X \to S$ 和 $b : Y \to S$ 是具有相同靶的概形态射。
设 $p : X \times_S Y \to X$ 和 $q : X \times_S Y \to Y$ 为投影态射，并令
$f = a \circ p = b \circ q$。图示如下：
$$
\xymatrix{
& X \times_S Y \ar[ld]^p \ar[rd]_q \ar[dd]^f \\
X \ar[rd]_a & & Y \ar[ld]^b \\
& S
}
$$
在本节中，给定 $\mathcal{O}_X$-模 $\mathcal{F}$ 和
$\mathcal{O}_Y$-模 $\mathcal{G}$，记
$$
\mathcal{F} \boxtimes \mathcal{G} =
p^*\mathcal{F} \otimes_{\mathcal{O}_{X \times_S Y}} q^*\mathcal{G}
$$
拟凝聚模上的双函子
$(\mathcal{F}, \mathcal{G}) \mapsto \mathcal{F} \boxtimes \mathcal{G}$
可扩张为以拟凝聚模为对象、以 $S$ 上有限阶微分算子为态射的双函子；
参见《态射》，备注 \ref{morphisms-remark-base-change-differential-operators}。
由《模》，引理
\ref{modules-lemma-differentials-relative-de-rham-complex-order-1}，
德拉姆复形 $\Omega^\bullet_{X/S}$ 和 $\Omega^\bullet_{Y/S}$ 的微分是
$1$ 阶的、定义在 $S$ 上的微分算子。因此可以考虑复形
$$
\text{Tot}(\Omega^\bullet_{X/S} \boxtimes \Omega^\bullet_{Y/S})
$$
请参见《概形的导出范畴》，第
\ref{perfect-section-kunneth-complexes} 节中的讨论。

\begin{lemma}
\label{derham-lemma-de-rham-complex-product}
在上述情形下，存在典范同构
$$
\text{Tot}(\Omega^\bullet_{X/S} \boxtimes \Omega^\bullet_{Y/S})
\longrightarrow
\Omega^\bullet_{X \times_S Y/S}
$$
它是 $f^{-1}\mathcal{O}_S$-模复形的同构。
\end{lemma}

\begin{proof}
由《态射》，引理 \ref{morphisms-lemma-differential-product}，可知
$
\Omega_{X \times_S Y/S} = p^*\Omega_{X/S} \oplus q^*\Omega_{Y/S}
$
。取外幂得到
$$
\Omega^n_{X \times_S Y/S} =
\bigoplus\nolimits_{i + j = n}
p^*\Omega^i_{X/S} \otimes_{\mathcal{O}_{X \times_S Y}} q^*\Omega^j_{Y/S} =
\bigoplus\nolimits_{i + j = n}
\Omega^i_{X/S} \boxtimes \Omega^j_{Y/S}
$$
这是外幂的初等性质。这些等同给出了本引理箭头两侧复形各项之间的同构。
我们略去验证这些映射与微分相容的步骤。
\end{proof}

\noindent
置 $A = \Gamma(S, \mathcal{O}_S)$。把引理
\ref{derham-lemma-de-rham-complex-product} 的结论与《概形的导出范畴》中的映射
（方程 \ref{perfect-equation-de-rham-kunneth}）结合，得到杯积
$$
R\Gamma(X, \Omega^\bullet_{X/S})
\otimes_A^\mathbf{L}
R\Gamma(Y, \Omega^\bullet_{Y/S})
\longrightarrow
R\Gamma(X \times_S Y, \Omega^\bullet_{X \times_S Y/S})
$$
在上同调层面，利用《更多代数》，第
\ref{more-algebra-section-products-tor} 节中的讨论，得到典范映射
$$
H^i_{dR}(X/S) \otimes_A H^j_{dR}(Y/S)
\longrightarrow
H^{i + j}_{dR}(X \times_S Y/S),\quad
(\xi, \zeta) \longmapsto p^*\xi \cup q^*\zeta
$$
注意，上述构造确实是先作拉回，再取杯积。

\begin{lemma}
\label{derham-lemma-kunneth-de-rham}
假设 $X$ 和 $Y$ 在 $S = \Spec(A)$ 上光滑、拟紧，且对角态射仿射。则映射
$$
R\Gamma(X, \Omega^\bullet_{X/S})
\otimes_A^\mathbf{L}
R\Gamma(Y, \Omega^\bullet_{Y/S})
\longrightarrow
R\Gamma(X \times_S Y, \Omega^\bullet_{X \times_S Y/S})
$$
是 $D(A)$ 中的同构。
\end{lemma}

\begin{proof}
由《态射》，引理 \ref{morphisms-lemma-smooth-omega-finite-locally-free}，
层 $\Omega^n_{X/S}$ 和 $\Omega^m_{Y/S}$ 分别是有限局部自由
$\mathcal{O}_X$-模和 $\mathcal{O}_Y$-模。另一方面，$X$ 和 $Y$ 在 $S$ 上平坦
（《态射》，引理 \ref{morphisms-lemma-smooth-flat}），
因而 $\Omega^n_{X/S}$ 和 $\Omega^m_{Y/S}$ 在 $S$ 上平坦。
还要注意，$\Omega^\bullet_{X/S}$ 局部有界。因此，结论由引理
\ref{derham-lemma-de-rham-complex-product} 和《概形的导出范畴》，引理
\ref{perfect-lemma-kunneth-special} 推出。
\end{proof}

\noindent
杯积还有相对版本，即映射
$$
Ra_*\Omega^\bullet_{X/S}
\otimes_{\mathcal{O}_S}^\mathbf{L}
Rb_*\Omega^\bullet_{Y/S}
\longrightarrow
Rf_*\Omega^\bullet_{X \times_S Y/S}
$$
它位于 $D(\mathcal{O}_S)$ 中。该构造把引理
\ref{derham-lemma-de-rham-complex-product} 与《概形的导出范畴》中的映射
（方程 \ref{perfect-equation-relative-de-rham-kunneth}）结合起来。
由此构造可见，该映射由下图给出：
$$
\xymatrix{
Ra_*\Omega^\bullet_{X/S}
\otimes_{\mathcal{O}_S}^\mathbf{L}
Rb_*\Omega^\bullet_{Y/S}
\ar[d]^{\text{伴随单位}}  \\
Rf_*(p^{-1}\Omega^\bullet_{X/S})
\otimes_{\mathcal{O}_S}^\mathbf{L}
Rf_*(q^{-1}\Omega^\bullet_{Y/S}) \ar[r] \ar[d]^{\text{相对杯积}} &
Rf_*(\Omega^\bullet_{X \times_S Y/S})
\otimes_{\mathcal{O}_S}^\mathbf{L}
Rf_*(\Omega^\bullet_{X \times_S Y/S}) \ar[d]^{\text{相对杯积}} \\
Rf_*(p^{-1}\Omega^\bullet_{X/S}
\otimes_{f^{-1}\mathcal{O}_S}^\mathbf{L}
q^{-1}\Omega^\bullet_{Y/S})
\ar[d]^{\text{从导出函子到通常函子}} \ar[r] &
Rf_*(\Omega^\bullet_{X \times_S Y/S}
\otimes_{f^{-1}\mathcal{O}_S}^\mathbf{L}
\Omega^\bullet_{X \times_S Y/S})
\ar[d]^{\text{从导出函子到通常函子}} \\
Rf_*\text{Tot}(p^{-1}\Omega^\bullet_{X/S}
\otimes_{f^{-1}\mathcal{O}_S}
q^{-1}\Omega^\bullet_{Y/S}) \ar[r] \ar[d]^{\text{典范映射}} &
Rf_*\text{Tot}(\Omega^\bullet_{X \times_S Y/S}
\otimes_{f^{-1}\mathcal{O}_S}
\Omega^\bullet_{X \times_S Y/S})
\ar[d]^{\eta \otimes \omega \mapsto \eta \wedge \omega} \\
Rf_*\text{Tot}(\Omega^\bullet_{X/S} \boxtimes \Omega^\bullet_{Y/S})
\ar@{=}[r]
&
Rf_*\Omega^\bullet_{X \times_S Y/S}
}
$$
这里，第一个箭头使用伴随的单位 $\text{id} \to Rp_* p^{-1}$
和 $\text{id} \to Rq_* q^{-1}$，以及等同
$Rf_* p^{-1} = Ra_* Rp_* p^{-1}$ 和
$Rf_* q^{-1} = Rb_* Rq_* q^{-1}$。
第二个箭头是《上同调》，备注 \ref{cohomology-remark-cup-product} 中的相对杯积。
第三个箭头把复形的导出张量积送到复形张量积的全复形。
最后一个等号来自引理 \ref{derham-lemma-de-rham-complex-product}。
在整体截面上，该构造恢复了前面给出的构造。

\begin{lemma}
\label{derham-lemma-kunneth-de-rham-relative}
假设 $X \to S$ 和 $Y \to S$ 光滑且拟紧，并且态射
$X \to X \times_S X$ 和 $Y \to Y \times_S Y$ 仿射。则相对杯积
$$
Ra_*\Omega^\bullet_{X/S}
\otimes_{\mathcal{O}_S}^\mathbf{L}
Rb_*\Omega^\bullet_{Y/S}
\longrightarrow
Rf_*\Omega^\bullet_{X \times_S Y/S}
$$
是 $D(\mathcal{O}_S)$ 中的同构。
\end{lemma}

\begin{proof}
这是引理 \ref{derham-lemma-kunneth-de-rham} 的直接推论。
\end{proof}







\section{德拉姆上同调中的第一陈类}
\label{derham-section-first-chern-class}

\noindent
设 $X \to S$ 是概形态射。存在复形映射
$$
\text{d}\log : \mathcal{O}_X^*[-1] \longrightarrow \Omega^\bullet_{X/S}
$$
它把截面 $g \in \mathcal{O}_X^*(U)$ 送到截面
$\text{d}\log(g) = g^{-1}\text{d}g$，后者属于 $\Omega^1_{X/S}(U)$。
因此可以考虑映射
$$
\Pic(X) = H^1(X, \mathcal{O}_X^*) =
H^2(X, \mathcal{O}_X^*[-1]) \longrightarrow H^2_{dR}(X/S)
$$
其中第一个等号来自《上同调》，引理
\ref{cohomology-lemma-h1-invertible}。
可逆模 $\mathcal{L}$ 的同构类的像记为
$c^{dR}_1(\mathcal{L}) \in H^2_{dR}(X/S)$。

\medskip\noindent
还可以利用映射 $\text{d}\log : \mathcal{O}_X^* \to \Omega^1_{X/S}$
在霍奇上同调中定义陈类
$$
c_1^{Hodge} : \Pic(X) \longrightarrow H^1(X, \Omega^1_{X/S})
\subset H^2_{Hodge}(X/S)
$$
这些构造与拉回相容。

\begin{lemma}
\label{derham-lemma-pullback-c1}
给定概形的交换图
$$
\xymatrix{
X' \ar[r]_f \ar[d] & X \ar[d] \\
S' \ar[r] & S
}
$$
则下列图
$$
\xymatrix{
\Pic(X') \ar[d]_{c_1^{dR}} &
\Pic(X) \ar[d]^{c_1^{dR}} \ar[l]^{f^*} \\
H^2_{dR}(X'/S') &
H^2_{dR}(X/S) \ar[l]_{f^*}
}
\quad
\xymatrix{
\Pic(X') \ar[d]_{c_1^{Hodge}} &
\Pic(X) \ar[d]^{c_1^{Hodge}} \ar[l]^{f^*} \\
H^1(X', \Omega^1_{X'/S'}) &
H^1(X, \Omega^1_{X/S}) \ar[l]_{f^*}
}
$$
均交换。
\end{lemma}

\begin{proof}
略。
\end{proof}

\noindent
让我们在 {\v C}ech 上同调中“计算”元素 $c^{dR}_1(\mathcal{L})$
（{\v C}ech 微分的符号规定同《上同调》，第
\ref{cohomology-section-cech-cohomology-of-complexes} 节）。
具体地，选取开覆盖
$\mathcal{U} : X = \bigcup_{i \in I} U_i$，使得有平凡化截面
$s_i$，它是 $\mathcal{L}|_{U_i}$ 的截面；这对所有 $i$ 都成立。
在交集 $U_{i_0i_1} = U_{i_0} \cap U_{i_1}$ 上，存在可逆函数
$f_{i_0i_1}$，使得
$f_{i_0i_1} = s_{i_1}|_{U_{i_0i_1}} s_{i_0}|_{U_{i_0i_1}}^{-1}$\footnote{
$0$-上闭链 $\{a_{i_0}\}$ 的 {\v C}ech 微分为
$a_{i_1} - a_{i_0}$，它定义在 $U_{i_0i_1}$ 上。}。
当然有
$$
f_{i_1i_2}|_{U_{i_0i_1i_2}}
f_{i_0i_2}^{-1}|_{U_{i_0i_1i_2}}
f_{i_0i_1}|_{U_{i_0i_1i_2}} = 1
$$
$\mathcal{L}$ 在 $H^1(X, \mathcal{O}_X^*)$ 中的上同调类，
是上闭链 $\{f_{i_0i_1}\}$ 在
$\check{\mathcal{C}}^\bullet(\mathcal{U}, \mathcal{O}_X^*)$
中的 {\v C}ech 上同调类之像。因此，$c_1^{dR}(\mathcal{L})$ 是
与 {\v C}ech 上闭链 $\{\alpha_{i_0 \ldots i_p}\}$ 相伴的上同调类之像；
该上闭链位于
$\text{Tot}(\check{\mathcal{C}}^\bullet(\mathcal{U}, \Omega_{X/S}^\bullet))$
中、次数为 $2$，并由下列各项给出：
\begin{enumerate}
\item $\alpha_{i_0} = 0$ 属于 $\Omega^2_{X/S}(U_{i_0})$；
\item $\alpha_{i_0i_1} = f_{i_0i_1}^{-1}\text{d}f_{i_0i_1}$ 属于
$\Omega^1_{X/S}(U_{i_0i_1})$；并且
\item $\alpha_{i_0i_1i_2} = 0$ 属于 $\mathcal{O}_{X/S}(U_{i_0i_1i_2})$。
\end{enumerate}
假设有可逆模 $\mathcal{L}_k$（$k = 1, \ldots, a$），
每个模在所有 $U_i$（$i \in I$）上都已平凡化，并由此得到上述上闭链
$f_{k, i_0i_1}$ 和 $\alpha_k = \{\alpha_{k, i_0 \ldots i_p}\}$。
利用《上同调》，第
\ref{cohomology-section-cech-cohomology-of-complexes} 节中的法则，可以算出
$$
\beta = \alpha_1 \cup \alpha_2 \cup \ldots \cup \alpha_a
$$
它由上闭链 $\beta = \{\beta_{i_0 \ldots i_p}\}$ 给出；后者描述如下：
\begin{enumerate}
\item $\beta_{i_0 \ldots i_p} = 0$ 属于
$\Omega^{2a - p}_{X/S}(U_{i_0 \ldots i_p})$，除非 $p = a$；并且
\item $\beta_{i_0 \ldots i_a} = (-1)^{a(a - 1)/2}
\alpha_{1, i_0i_1} \wedge \alpha_{2, i_1 i_2} \wedge \ldots \wedge
\alpha_{a, i_{a - 1}i_a}$ 属于
$\Omega^a_{X/S}(U_{i_0 \ldots i_a})$。
\end{enumerate}
因此，这是表示
$c_1^{dR}(\mathcal{L}_1) \cup \ldots \cup c_1^{dR}(\mathcal{L}_a)$ 的上闭链。
当然，同一计算还表明：上闭链
$\{\beta_{i_0 \ldots i_a}\}$ 位于
$\check{\mathcal{C}}^a(\mathcal{U}, \Omega_{X/S}^a))$ 中，并表示上同调类
$c_1^{Hodge}(\mathcal{L}_1) \cup \ldots \cup c_1^{Hodge}(\mathcal{L}_a)$。

\begin{remark}
\label{derham-remark-truncations}
下面用更抽象的语言重新表述上述计算。
设 $p : X \to S$ 是概形态射，并设 $\mathcal{L}$ 是可逆
$\mathcal{O}_X$-模。若把 $\text{d}\log$ 看作映射
$$
\mathcal{O}_X^*[-1] \to \sigma_{\geq 1}\Omega^\bullet_{X/S}
$$
则如上利用 $\Pic(X) = H^1(X, \mathcal{O}_X^*)$，得到上同调类
$$
\gamma_1(\mathcal{L}) \in H^2(X, \sigma_{\geq 1}\Omega^\bullet_{X/S})
$$
$\gamma_1(\mathcal{L})$ 在映射
$\sigma_{\geq 1}\Omega^\bullet_{X/S} \to \Omega^\bullet_{X/S}$
下的像恢复 $c_1^{dR}(\mathcal{L})$。特别地，
$c_1^{dR}(\mathcal{L}) \in F^1H^2_{dR}(X/S)$, see
参见第 \ref{derham-section-hodge-filtration} 节。$\gamma_1(\mathcal{L})$ 在映射
$\sigma_{\geq 1}\Omega^\bullet_{X/S} \to \Omega^1_{X/S}[-1]$ 下的像
恢复 $c_1^{Hodge}(\mathcal{L})$。
取杯积（参见第 \ref{derham-section-hodge-filtration} 节），得到
$$
\xi = \gamma_1(\mathcal{L}_1) \cup \ldots \cup \gamma_1(\mathcal{L}_a) \in
H^{2a}(X, \sigma_{\geq a}\Omega^\bullet_{X/S})
$$
第 \ref{derham-section-hodge-filtration} 节中的交换图表明，$\xi$ 被映到
$c_1^{dR}(\mathcal{L}_1) \cup \ldots \cup c_1^{dR}(\mathcal{L}_a)$
，后者位于 $H^{2a}_{dR}(X/S)$ 中；这里所用的映射为
$\sigma_{\geq a}\Omega^\bullet_{X/S} \to \Omega^\bullet_{X/S}$。
还可知
$c_1^{dR}(\mathcal{L}_1) \cup \ldots \cup c_1^{dR}(\mathcal{L}_a)$
属于 $F^a H^{2a}_{dR}(X/S)$。类似地，映射
$\sigma_{\geq a}\Omega^\bullet_{X/S} \to \Omega^a_{X/S}[-a]$
把 $\xi$ 送到
$c_1^{Hodge}(\mathcal{L}_1) \cup \ldots \cup c_1^{Hodge}(\mathcal{L}_a)$
；后者位于 $H^a(X, \Omega^a_{X/S})$ 中。
\end{remark}

\begin{remark}
\label{derham-remark-log-forms}
设 $p : X \to S$ 是概形态射。对 $i > 0$，以
$\Omega^i_{X/S, log} \subset \Omega^i_{X/S}$ 表示由下述形式的局部截面
生成的阿贝尔子层：
$$
\text{d}\log(u_1) \wedge \ldots \wedge \text{d}\log(u_i)
$$
其中 $u_1, \ldots, u_n$ 是 $\mathcal{O}_X$ 的可逆局部截面。
当 $i = 0$ 时，子层 $\Omega^0_{X/S, log} \subset \mathcal{O}_X$
是 $\mathbf{Z} \to \mathcal{O}_X$ 的像。对每个 $i \geq 0$，都有复形映射
$$
\Omega^i_{X/S, log}[-i] \longrightarrow \Omega^\bullet_{X/S}
$$
因为对数形式的微分为零。此外，对数形式的楔积仍是对数形式，因而得到双线性映射
$$
\wedge :  \Omega^i_{X/S, log} \times
\Omega^j_{X/S, log} \longrightarrow \Omega^{i + j}_{X/S, log}
$$
它们与（\ref{derham-equation-wedge}）及上述映射相容。
设 $\mathcal{L}$ 是可逆 $\mathcal{O}_X$-模。利用阿贝尔层的映射
$\text{d}\log : \mathcal{O}_X^* \to \Omega^1_{X/S, log}$
和等同 $\Pic(X) = H^1(X, \mathcal{O}_X^*)$，得到典范上同调类
$$
\tilde \gamma_1(\mathcal{L}) \in H^1(X, \Omega^1_{X/S, log})
$$
这些类具有下列性质：
\begin{enumerate}
\item $\tilde \gamma_1(\mathcal{L})$ 在典范映射
$\Omega^1_{X/S, log}[-1] \to \sigma_{\geq 1}\Omega^\bullet_{X/S}$
下的像，是该映射把 $\tilde \gamma_1(\mathcal{L})$ 送到的类
$\gamma_1(\mathcal{L}) \in 
H^2(X, \sigma_{\geq 1}\Omega^\bullet_{X/S})$；
这就是备注 \ref{derham-remark-truncations} 中的类；
\item $\tilde \gamma_1(\mathcal{L})$ 在典范映射
$\Omega^1_{X/S, log}[-1] \to \Omega^\bullet_{X/S}$
下的像，是该映射把 $\tilde \gamma_1(\mathcal{L})$ 送到的
$c_1^{dR}(\mathcal{L})$，它位于 $H^2_{dR}(X/S)$ 中；
\item $\tilde \gamma_1(\mathcal{L})$ 在典范映射
$\Omega^1_{X/S, log} \to \Omega^1_{X/S}$
下的像，是该映射把 $\tilde \gamma_1(\mathcal{L})$ 送到的
$c_1^{Hodge}(\mathcal{L})$，它位于 $H^1(X, \Omega^1_{X/S})$ 中；
\item 这些类的构造与拉回相容；
\item 在此添加更多内容。
\end{enumerate}
\end{remark}





\section{线丛的德拉姆上同调}
\label{derham-section-line-bundle}

\noindent
线丛是向量丛的特例，而向量丛又是带有某些附加结构的锥。
为了恰当地讨论这些对象的德拉姆复形，先讨论分次环的德拉姆复形是合适的。

\begin{remark}[分次环的德拉姆复形]
\label{derham-remark-de-rham-complex-graded}
设 $G$ 是以加法记号书写、幺元为 $0$ 的阿贝尔幺半群。
设 $R \to A$ 是环同态，并假设 $A$ 带有分次
$A = \bigoplus_{g \in G} A_g$；它由 $R$-模给出，使得 $R$ 映入 $A_0$，且
$A_g \cdot A_{g'} \subset A_{g + g'}$。于是微分模带有分次
$$
\Omega_{A/R} = \bigoplus\nolimits_{g \in G} \Omega_{A/R, g}
$$
其中 $\Omega_{A/R, g}$ 是 $R$-子模；它包含于 $\Omega_{A/R}$，并
由 $a_0 \text{d}a_1$ 生成；这里 $a_i \in A_{g_i}$，且
$g = g_0 + g_1$。类似地，得到
$$
\Omega^p_{A/R} = \bigoplus\nolimits_{g \in G} \Omega^p_{A/R, g}
$$
其中 $\Omega^p_{A/R, g}$ 是 $R$-子模；它包含于 $\Omega^p_{A/R}$，并
由 $a_0 \text{d}a_1 \wedge \ldots \wedge \text{d}a_p$ 生成；
这里 $a_i \in A_{g_i}$，且 $g = g_0 + g_1 + \ldots + g_p$。
显然，微分保持分次，而楔积以显然的方式与分次相容。
\end{remark}

\noindent
设 $f : X \to S$ 是概形态射，并设 $\pi : C \to X$ 是一个锥；参见
《构造》，定义 \ref{constructions-definition-abstract-cone}。
回忆这意味着 $\pi$ 仿射，且有分次
$\pi_*\mathcal{O}_C = \bigoplus_{n \geq 0} \mathcal{A}_n$，其中
$\mathcal{A}_0 = \mathcal{O}_X$。
在仿射开集上应用备注 \ref{derham-remark-de-rham-complex-graded} 中的讨论，可得
\footnote{请原谅这里的记号！}
$$
\pi_*(\Omega^\bullet_{C/S}) =
\bigoplus\nolimits_{n \geq 0} \Omega^\bullet_{C/S, n}
$$
典范地分解为子复形的直和。此外，还有分解
$$
\Omega^\bullet_{X/S} \to \Omega^\bullet_{C/S, 0} \to
\pi_*(\Omega^\bullet_{C/S})
$$
并且知道 $\omega \wedge \eta \in \Omega^{p + q}_{C/S, n + m}$，
只要 $\omega \in \Omega^p_{C/S, n}$ 且
$\eta \in \Omega^q_{C/S, m}$。

\medskip\noindent
设 $f : X \to S$ 是概形态射，并设 $\pi : L \to X$ 是与可逆
$\mathcal{O}_X$-模 $\mathcal{L}$ 相伴的线丛。这意味着 $\pi$
是满足下式的唯一仿射态射：
$$
\pi_*\mathcal{O}_L = 
\bigoplus\nolimits_{n \geq 0} \mathcal{L}^{\otimes n}
$$
这里等式指 $\mathcal{O}_X$-代数的等式。因此 $L$ 是 $X$ 上的锥。
由上述讨论得到典范直和分解
$$
\pi_*(\Omega^\bullet_{L/S}) =
\bigoplus\nolimits_{n \geq 0} \Omega^\bullet_{L/S, n}
$$
它与楔积以及上述 $\pi_*\mathcal{O}_L$ 的分解相容，并使
$\Omega_{X/S}$ 映入部分 $\Omega_{L/S, 0}$；该部分的次数为 $0$。

\medskip\noindent
还有一种情形对我们有用。考虑补集\footnote{概形 $L^\star$ 是
$\mathbf{G}_m$-扭子；它位于 $X$ 上并与 $L$ 相伴。因此下面所得分次是
$\mathbf{Z}$-分次；比较《群胚》，例
\ref{groupoids-example-Gm-on-affine}，以及引理
\ref{groupoids-lemma-complete-reducibility-Gm} 和
\ref{groupoids-lemma-Gm-equivariant-module}。}
$L^\star \subset L$，它是线丛中零截面 $o : X \to L$ 的补，而该线丛为
$L$。局部计算表明，存在典范同构
$$
(L^\star \to X)_*\mathcal{O}_{L^\star} =
\bigoplus\nolimits_{n \in \mathbf{Z}} \mathcal{L}^{\otimes n}
$$
这是 $\mathcal{O}_X$-代数的同构。右端是 $\mathbf{Z}$-分次拟凝聚
$\mathcal{O}_X$-代数。在仿射开集上使用备注
\ref{derham-remark-de-rham-complex-graded} 中的讨论，得到
$$
(L^\star \to X)_*(\Omega^\bullet_{L^\star/S}) =
\bigoplus\nolimits_{n \in \mathbf{Z}} \Omega^\bullet_{L^\star/S, n}
$$
它与楔积以及上述 $(L^\star \to X)_*\mathcal{O}_{L^\star}$ 的分解相容，
并使 $\Omega_{X/S}$ 映入部分 $\Omega_{L^\star/S, 0}$；该部分的次数为 $0$。
复形 $\Omega^\bullet_{L^\star/S, 0}$ 对我们尤为重要。

\begin{lemma}
\label{derham-lemma-the-complex-for-L-star}
沿用上述记号，存在复形的短正合列
$$
0 \to \Omega^\bullet_{X/S} \to
\Omega^\bullet_{L^\star/S, 0} \to
\Omega^\bullet_{X/S}[-1] \to 0
$$
\end{lemma}

\begin{proof}
上面已经构造了映射
$\Omega^\bullet_{X/S} \to \Omega^\bullet_{L^\star/S, 0}$。

\medskip\noindent
构造
$\text{Res} : \Omega^\bullet_{L^\star/S, 0} \to \Omega^\bullet_{X/S}[-1]$。
设 $U \subset X$ 是开集，并设 $s \in \mathcal{L}(U)$ 和
$s' \in \mathcal{L}^{\otimes -1}(U)$ 是满足 $s' s = 1$ 的截面。
于是 $s$ 给出代数层 $(L^\star \to X)_*\mathcal{O}_{L^\star}$
在 $U$ 上的可逆截面，其逆为 $s' = s^{-1}$。因而可以考虑 $1$-形式
$\text{d}\log(s) = s' \text{d}(s)$；它属于
$\Omega^1_{L^\star/S, 0}(U)$，这是由我们在
$\Omega^1_{L^\star/S}$ 上构造的分次得到的。下面的仿射计算将表明，$1$ 和
$\text{d}\log(s)$ 自由生成 $\Omega^\bullet_{L^\star/S, 0}|_U$；
这里采用其作为 $\Omega^\bullet_{X/S}|_U$ 上右模的结构。
因此，可以按下列法则定义 $\text{Res}$；它定义在 $U$ 上：
$$
\text{Res}(\omega' + \text{d}\log(s) \wedge \omega) = \omega
$$
这里 $\omega', \omega \in \Omega^\bullet_{X/S}(U)$ 任意。
该映射与局部生成元 $s$ 的选择无关，因而可以黏合成整体映射。
事实上，$s$ 的另一个选择必为 $gs$ 的形式，其中
$g \in \mathcal{O}_X(U)$ 可逆；此时
$\text{d}\log(gs) = g^{-1}\text{d}(g) + \text{d}\log(s)$，
由此容易看出无关性。最后注意，$\text{Res}$ 的定义法则与微分相容，因为
$\text{d}(\omega' + \text{d}\log(s) \wedge \omega) =
\text{d}(\omega') - \text{d}\log(s) \wedge \text{d}(\omega)$
，并且按照《同调》，定义
\ref{homology-definition-shift-cochain} 中的符号约定，
$\Omega^\bullet_{X/S}[-1]$ 上的微分把 $\omega'$ 送到
$-\text{d}(\omega')$。

\medskip\noindent
局部计算。可以用仿射开集覆盖 $X$，这些开集记为 $U \subset X$，
并满足 $\mathcal{L}|_U \cong \mathcal{O}_U$，而且映入某个仿射开集
$V \subset S$。写成 $U = \Spec(A)$、$V = \Spec(R)$，并选取
生成元 $s$；它属于 $\mathcal{L}$。于是有
$$
L^\star \times_X U = \Spec(A[s, s^{-1}])
$$
计算微分可得
$$
\Omega^1_{A[s, s^{-1}]/R} =
A[s, s^{-1}] \otimes_A \Omega^1_{A/R} \oplus A[s, s^{-1}] \text{d}\log(s)
$$
因而取外幂得到
$$
\Omega^p_{A[s, s^{-1}]/R} =
A[s, s^{-1}] \otimes_A \Omega^p_{A/R}
\oplus
A[s, s^{-1}] \text{d}\log(s) \otimes_A \Omega^{p - 1}_{A/R}
$$
取次数为 $0$ 的部分，得到
$$
\Omega^p_{A[s, s^{-1}]/R, 0} =
\Omega^p_{A/R} \oplus \text{d}\log(s) \otimes_A \Omega^{p - 1}_{A/R}
$$
引理得证。
\end{proof}

\begin{lemma}
\label{derham-lemma-the-complex-for-L-star-gives-chern-class}
“边缘”映射
$\delta : \Omega^\bullet_{X/S} \to \Omega^\bullet_{X/S}[2]$
位于 $D(X, f^{-1}\mathcal{O}_S)$ 中，它来自引理
\ref{derham-lemma-the-complex-for-L-star} 的短正合列；它就是备注
\ref{derham-remark-cup-product-as-a-map} 中对应于
$\xi = c_1^{dR}(\mathcal{L})$ 的映射。
\end{lemma}

\begin{proof}
准确地说，考虑短正合列的平移
$$
0 \to \Omega^\bullet_{X/S}[1] \to
\Omega^\bullet_{L^\star/S, 0}[1] \to
\Omega^\bullet_{X/S} \to 0
$$
该短正合列来自引理 \ref{derham-lemma-the-complex-for-L-star}。
作为下列分次的次数零部分：
$(L^\star \to X)_*\Omega^\bullet_{L^\star/S}$
，$\Omega^\bullet_{L^\star/S, 0}$ 是微分分次
$\mathcal{O}_X$-代数，并且映射
$\Omega^\bullet_{X/S} \to \Omega^\bullet_{L^\star/S, 0}$
是微分分次 $\mathcal{O}_X$-代数的同态。
因此可以把 $\Omega^\bullet_{X/S}[1] \to
\Omega^\bullet_{L^\star/S, 0}[1]$ 视为右微分分次
$\Omega^\bullet_{X/S}$-模的映射；它定义在 $X$ 上。映射
$\text{Res} : \Omega^\bullet_{L^\star/S, 0}[1] \to \Omega^\bullet_{X/S}$
是右微分分次 $\Omega^\bullet_{X/S}$-模的映射，因为它局部地由法则
$\text{Res}(\omega' + \text{d}\log(s) \wedge \omega) = \omega$ 定义；
定义；参见引理 \ref{derham-lemma-the-complex-for-L-star} 的证明。
因此，由《微分分次层》，第 \ref{sdga-section-misc} 节中的讨论可知，
$\delta$ 来自映射
$\delta' : \Omega^\bullet_{X/S} \to \Omega^\bullet_{X/S}[2]$
；它位于德拉姆复形上的右微分分次模的导出范畴
$D(\Omega^\bullet_{X/S}, \text{d})$ 中。备注
\ref{derham-remark-cup-product-as-a-map} 所述的唯一性表明，只需证明
$\delta(1) = c_1^{dR}(\mathcal{L})$。

\medskip\noindent
我们断言存在交换图
$$
\xymatrix{
0 \ar[r] &
\mathcal{O}_X^* \ar[r] \ar[d]_{\text{d}\log} &
E \ar[r] \ar[d] &
\underline{\mathbf{Z}} \ar[d] \ar[r] &
0 \\
0 \ar[r] &
\Omega^\bullet_{X/S}[1] \ar[r] &
\Omega^\bullet_{L^\star/S, 0}[1] \ar[r] &
\Omega^\bullet_{X/S} \ar[r] &
0
}
$$
其中上行是阿贝尔层的短正合列，其边缘映射把 $1$ 送到
$\mathcal{L}$ 在 $H^1(X, \mathcal{O}_X^*)$ 中的类。
由边缘映射同短正合列之间映射的相容性，只需证明该断言。
把 $E$ 定义为下述法则的层化：
$$
U \longmapsto \{(s, n) \mid
n \in \mathbf{Z},\ s \in \mathcal{L}^{\otimes n}(U)\text{ 生成元}\}
$$
其群结构由 $(s, n) \cdot (t, m) = (s \otimes t, n + m)$ 给出。
中间的竖直映射把 $(s, n)$ 送到 $\text{d}\log(s)$。
这给出短正合列之间的映射，因为证明引理
\ref{derham-lemma-the-complex-for-L-star} 时构造的映射
$Res : \Omega^1_{L^\star/S, 0} \to \mathcal{O}_X$ 把
$\text{d}\log(s)$ 送到 $1$，只要 $s$ 是 $\mathcal{L}$ 的局部生成元。
为了计算上行中 $1$ 的边缘，按照第
\ref{derham-section-first-chern-class} 节，选取局部平凡化 $s_i$；
它们来自 $\mathcal{L}$，并定义在开集 $U_i$ 上。在交集
$U_{i_0i_1} = U_{i_0} \cap U_{i_1}$
上，存在可逆函数 $f_{i_0i_1}$，使得
$f_{i_0i_1} = s_{i_1}|_{U_{i_0i_1}} s_{i_0}|_{U_{i_0i_1}}^{-1}$
，而 $\mathcal{L}$ 的上同调类由 {\v C}ech 上闭链
$\{f_{i_0i_1}\}$ 给出。于是当然有
$$
(f_{i_0i_1}, 0) = (s_{i_1}, 1)|_{U_{i_0i_1}} \cdot
(s_{i_0}, 1)|_{U_{i_0i_1}}^{-1}
$$
它们是 $E$ 的截面。这就完成了证明。
\end{proof}

\begin{lemma}
\label{derham-lemma-push-omega-a}
沿用上述记号，有
\begin{enumerate}
\item $\Omega^p_{L^\star/S, n} =
\Omega^p_{L^\star/S, 0} \otimes_{\mathcal{O}_X} \mathcal{L}^{\otimes n}$
该等式对所有 $n \in \mathbf{Z}$ 成立；这里把两边视为拟凝聚
$\mathcal{O}_X$-模；
\item $\Omega^\bullet_{X/S} = \Omega^\bullet_{L/X, 0}$
这是复形的等式；并且
\item 对 $n > 0$ 和 $p \geq 0$，有
$\Omega^p_{L/X, n} = \Omega^p_{L^\star/S, n}$.
\end{enumerate}
\end{lemma}

\begin{proof}
在每一种情形下，都存在整体定义的典范映射；略去的局部计算表明它是同构。
\end{proof}

\begin{lemma}
\label{derham-lemma-line-bundle-characteristic-zero}
在上述情形下，假设存在态射 $S \to \Spec(\mathbf{Q})$。
则 $\Omega^\bullet_{X/S} \to \pi_*\Omega^\bullet_{L/S}$ 是拟同构，并且
$H_{dR}^*(X/S) = H_{dR}^*(L/S)$。
\end{lemma}

\begin{proof}
设 $R$ 是 $\mathbf{Q}$-代数，并设 $A$ 是 $R$-代数。
仿射局部的陈述是：映射
$$
\Omega^\bullet_{A/R} \longrightarrow \Omega^\bullet_{A[t]/R}
$$
是 $R$-模复形的拟同构。事实上，它是同伦等价；其同伦逆由把
$g \omega + g' \text{d}t \wedge \omega'$ 送到 $g(0)\omega$ 的映射给出；
这里 $g, g' \in A[t]$，且 $\omega, \omega' \in \Omega^\bullet_{A/R}$。
该同伦把 $g \omega + g' \text{d}t \wedge \omega'$ 送到
$(\int g') \omega'$；其中 $\int g' \in A[t]$ 是常数项为零、
关于 $t$ 的导数为 $g'$ 的多项式。当然，这里用到 $R$ 包含
$\mathbf{Q}$，因为 $\int t^n = (1/n)t^{n + 1}$。
\end{proof}

\begin{example}
\label{derham-example-affine-line}
引理 \ref{derham-lemma-line-bundle-characteristic-zero} 在正特征下不成立。
$\mathbf{A}^1_k = \Spec(k[x])$ 在域 $k$ 上的德拉姆复形具有如下直和形式：
$$
k \oplus
\bigoplus\nolimits_{n \geq 1}
(k \cdot t^n \xrightarrow{n}
k \cdot t^{n - 1} \text{d}t)
$$
因此，若 $k$ 的特征为 $p > 0$，则
$H^0_{dR}(\mathbf{A}^1_k/k)$ 和
$H^1_{dR}(\mathbf{A}^1_k/k)$
作为 $k$-向量空间都是无限维的。
\end{example}








\section{射影空间的德拉姆上同调}
\label{derham-section-projective-space}

\noindent
设 $A$ 是环，且 $n \geq 1$。结构态射
$\mathbf{P}^n_A \to \Spec(A)$ 是相对维数为 $n$ 的真光滑态射。
它光滑、相对维数为 $n$ 且为有限型，因为 $\mathbf{P}^n_A$ 有有限仿射开覆盖，
其中每个概形都同构于 $\mathbf{A}^n_A$；参见《构造》，引理
\ref{constructions-lemma-standard-covering-projective-space}。
由《构造》，引理 \ref{constructions-lemma-projective-space-separated}，
它还分离且普遍闭，因而为真。
用 $\mathcal{O}$ 和 $\mathcal{O}(d)$ 分别表示结构层
$\mathcal{O}_{\mathbf{P}^n_A}$ 与 Serre 扭曲
$\mathcal{O}_{\mathbf{P}^n_A}(d)$。
用 $\Omega = \Omega_{\mathbf{P}^n_A/A}$ 表示相对微分层，并用
$\Omega^p$ 表示其外幂。

\begin{lemma}
\label{derham-lemma-euler-sequence}
存在短正合列
$$
0 \to \Omega \to \mathcal{O}(-1)^{\oplus n + 1} \to \mathcal{O} \to 0
$$
\end{lemma}

\begin{proof}
为说明这一点，回忆
$\mathbf{P}^n_A = \text{Proj}(A[T_0, \ldots, T_n])$,
并采用形式记号
$$
\mathcal{O}(-1)^{\oplus n + 1} =
\bigoplus\nolimits_{j = 0, \ldots, n} \mathcal{O}(-1) \text{d}T_j
$$
上述短正合列中的第一个箭头
$$
\Omega \to
\bigoplus\nolimits_{j = 0, \ldots, n} \mathcal{O}(-1) \text{d}T_j
$$
在每个标准开集
$D_+(T_i) = \Spec(A[T_0/T_i, \ldots, T_n/T_i])$
上由下列法则给出：
$$
\sum\nolimits_{j \not = i} g_j \text{d}(T_j/T_i)
\longmapsto
\sum\nolimits_{j \not = i} g_j/T_i \text{d}T_j
- (\sum\nolimits_{j \not = i} g_jT_j/T_i^2) \text{d}T_i
$$
这是有意义的，因为 $1/T_i$ 是 $\mathcal{O}(-1)$ 在
$D_+(T_i)$ 上的截面。映射
$$
\bigoplus\nolimits_{j = 0, \ldots, n} \mathcal{O}(-1) \text{d}T_j
\to
\mathcal{O}
$$
把 $\text{d}T_j$ 送到 $T_j$。更准确地说，在 $D_+(T_i)$ 上，
把截面 $\sum g_j \text{d}T_j$ 送到 $\sum T_jg_j$。
略去验证这确实给出短正合列。
\end{proof}

\noindent
给定整数 $k \in \mathbf{Z}$ 和拟凝聚
$\mathcal{O}_{\mathbf{P}^n_A}$-模 $\mathcal{F}$，照例用
$\mathcal{F}(k)$ 表示第 $k$ 个 Serre 扭曲；被扭曲的模是 $\mathcal{F}$。
参见《构造》，定义 \ref{constructions-definition-twist}。

\begin{lemma}
\label{derham-lemma-twisted-hodge-cohomology-projective-space}
在上述情形下，上同调群满足：
\begin{enumerate}
\item $H^q(\mathbf{P}^n_A, \Omega^p) = 0$，除非
$0 \leq p = q \leq n$；
\item 对 $0 \leq p \leq n$，$A$-模
$H^p(\mathbf{P}^n_A, \Omega^p)$ 是秩为 $1$ 的自由模；
\item 对 $q > 0$、$k > 0$ 和任意 $p$，有
$H^q(\mathbf{P}^n_A, \Omega^p(k)) = 0$；并且
\item 在此添加更多内容。
\end{enumerate}
\end{lemma}

\begin{proof}
以下将不再逐次说明地使用《概形上同调》，引理
\ref{coherent-lemma-cohomology-projective-space-over-ring} 的结论。
特别地，相应陈述对 $H^q(\mathbf{P}^n_A, \mathcal{O}(k))$ 成立。

\medskip\noindent
先证明 $p = 1$ 的情形。考虑短正合列
$$
0 \to \Omega \to \mathcal{O}(-1)^{\oplus n + 1} \to \mathcal{O} \to 0
$$
它来自引理 \ref{derham-lemma-euler-sequence}。由于 $\mathcal{O}(-1)$
在所有次数上的上同调均为零，所以
$H^q(\mathbf{P}^n_A, \Omega)$ 仅在次数 $1$ 可能非零；
在该次数，它由 $1$ 的边缘自由生成；该元素位于
$H^0(\mathbf{P}^n_A, \mathcal{O})$ 中。

\medskip\noindent
假设 $p > 1$。把上述短正合列看作定义了一个 $2$ 步滤过；
该滤过位于 $\mathcal{O}(-1)^{\oplus n + 1}$ 上。
它在 $\wedge^p\mathcal{O}(-1)^{\oplus n + 1}$ 上诱导的滤过如下：
$$
0 \to \Omega^p \to \wedge^p\left(\mathcal{O}(-1)^{\oplus n + 1}\right)
\to \Omega^{p - 1} \to 0
$$
注意，$\wedge^p\mathcal{O}(-1)^{\oplus n + 1}$ 同构于
$n + 1$ 取 $p$ 个 $\mathcal{O}(-p)$ 副本的直和，
因而在所有次数上的上同调均为零。由归纳假设可知，
$H^q(\mathbf{P}^n_A, \Omega^p)$ 为零，除非 $q = p$；而
$H^p(\mathbf{P}^n_A, \Omega^p)$ 是秩为 $1$ 的自由模，其生成元是
$H^{p - 1}(\mathbf{P}^n_A, \Omega^{p - 1})$ 中生成元的边缘。

\medskip\noindent
设 $k > 0$。注意 $\Omega^n = \mathcal{O}(-n - 1)$；例如，
这可由上述短正合列取 $p = n + 1$ 得到。因此 $\Omega^n(k)$
在正次数上的上同调均为零。利用短正合列
$$
0 \to \Omega^p(k) \to \wedge^p\left(\mathcal{O}(-1)^{\oplus n + 1}\right)(k)
\to \Omega^{p - 1}(k) \to 0
$$
并对 $p$ 作{\it 降序}归纳，得到 $\Omega^p(k)$
在正次数上的上同调均为零；这对所有 $p$ 成立。
\end{proof}

\begin{lemma}
\label{derham-lemma-hodge-cohomology-projective-space}
有 $H^q(\mathbf{P}^n_A, \Omega^p) = 0$，除非
$0 \leq p = q \leq n$。对 $0 \leq p \leq n$，$A$-模
$H^p(\mathbf{P}^n_A, \Omega^p)$ 是秩为 $1$ 的自由模，其基元素为
$c_1^{Hodge}(\mathcal{O}(1))^p$。
\end{lemma}

\begin{proof}
消失性与自由性由引理
\ref{derham-lemma-twisted-hodge-cohomology-projective-space} 给出。
当 $p = 0$ 时，$1 \in H^0(\mathbf{P}^n_A, \mathcal{O})$
显然是生成元。

\medskip\noindent
先证明 $p = 1$ 的情形。考虑短正合列
$$
0 \to \Omega \to \mathcal{O}(-1)^{\oplus n + 1} \to \mathcal{O} \to 0
$$
它来自引理 \ref{derham-lemma-euler-sequence}。在引理
\ref{derham-lemma-twisted-hodge-cohomology-projective-space} 的证明中已经看到，
$H^1(\mathbf{P}^n_A, \Omega)$ 的生成元是边缘 $\xi$；
它来自 $1 \in H^0(\mathbf{P}^n_A, \mathcal{O})$。
如引理 \ref{derham-lemma-euler-sequence} 的证明那样，把
$\mathcal{O}(-1)^{\oplus n + 1}$ 与
$\bigoplus_{j = 0, \ldots, n} \mathcal{O}(-1)\text{d}T_j$ 等同。
考虑开覆盖
$$
\mathcal{U} : 
\mathbf{P}^n_A =
\bigcup\nolimits_{i = 0, \ldots, n} D_{+}(T_i)
$$
整体截面 $1$（它属于 $\mathcal{O}$）在
$U_i = D_+(T_i)$ 上的限制，可以提升为截面
$T_i^{-1} \text{d}T_i$；它是
$\bigoplus \mathcal{O}(-1)\text{d}T_j$ 在 $U_i$ 上的截面。
因此，表示 $\xi$ 的上闭链由下式给出：
$$
T_{i_1}^{-1} \text{d}T_{i_1} - T_{i_0}^{-1} \text{d}T_{i_0} =
\text{d}\log(T_{i_1}/T_{i_0}) \in \Omega(U_{i_0i_1})
$$
另一方面，对每个 $i$，截面 $T_i$ 是 $\mathcal{O}(1)$
在 $U_i$ 上的平凡化截面。因此可见
$f_{i_0i_1} = T_{i_1}/T_{i_0} \in \mathcal{O}^*(U_{i_0i_1})$
是表示 $\mathcal{O}(1)$ 的上闭链；它位于 $\Pic(\mathbf{P}^n_A)$ 中。
参见第 \ref{derham-section-first-chern-class} 节。因此
$c_1^{Hodge}(\mathcal{O}(1))$ 由上闭链
$\text{d}\log(T_{i_1}/T_{i_0})$ 给出，这与上面对 $\xi$ 的结果一致。

\medskip\noindent
对一般的 $p$ 作归纳证明。基本情形 $p = 0, 1$ 已在上面处理。
假设 $p > 1$。在引理
\ref{derham-lemma-twisted-hodge-cohomology-projective-space} 的证明中已经看到，
$H^p(\mathbf{P}^n_A, \Omega^p)$ 的生成元，是
$c_1^{Hodge}(\mathcal{O}(1))^{p - 1}$ 在与下列短正合列相伴的
上同调长正合列中的边缘：
$$
0 \to \Omega^p \to \wedge^p\left(\mathcal{O}(-1)^{\oplus n + 1}\right)
\to \Omega^{p - 1} \to 0
$$
由第 \ref{derham-section-first-chern-class} 节中的计算，上同调类
$c_1^{Hodge}(\mathcal{O}(1))^{p - 1}$ 在相差一个符号的意义下，
由各项如下的上闭链表示：
$$
\beta_{i_0 \ldots i_{p - 1}} =
\text{d}\log(T_{i_1}/T_{i_0}) \wedge
\text{d}\log(T_{i_2}/T_{i_1}) \wedge \ldots \wedge
\text{d}\log(T_{i_{p - 1}}/T_{i_{p - 2}})
$$
它位于 $\Omega^{p - 1}(U_{i_0 \ldots i_{p - 1}})$ 中。各个
$\beta_{i_0 \ldots i_{p - 1}}$ 可以提升为截面
$\tilde \beta_{i_0 \ldots i_{p -1}} =
T_{i_0}^{-1}\text{d}T_{i_0} \wedge \beta_{i_0 \ldots i_{p - 1}}$
；这是 $\wedge^p(\bigoplus \mathcal{O}(-1) \text{d}T_j)$
在 $U_{i_0 \ldots i_{p - 1}}$ 上的截面。由此可知，
$H^p(\mathbf{P}^n_A, \Omega^p)$ 的生成元由下列上闭链给出，
其分量为
\begin{align*}
\sum\nolimits_{a = 0}^p (-1)^a
\tilde \beta_{i_0 \ldots \hat{i_a} \ldots i_p}
& =
T_{i_1}^{-1}\text{d}T_{i_1} \wedge \beta_{i_1 \ldots i_p}
+ \sum\nolimits_{a = 1}^p (-1)^a
T_{i_0}^{-1}\text{d}T_{i_0} \wedge
\beta_{i_0 \ldots \hat{i_a} \ldots i_p} \\
& =
(T_{i_1}^{-1}\text{d}T_{i_1} - T_{i_0}^{-1}\text{d}T_{i_0}) \wedge
\beta_{i_1 \ldots i_p} +
T_{i_0}^{-1}\text{d}T_{i_0} \wedge \text{d}(\beta)_{i_0 \ldots i_p} \\
& =
\text{d}\log(T_{i_1}/T_{i_0}) \wedge \beta_{i_1 \ldots i_p}
\end{align*}
把它视为 $\Omega^p$ 在 $U_{i_0 \ldots i_p}$ 上的截面。
除去符号，它与表示 $c_1^{Hodge}(\mathcal{O}(1))^p$ 的上闭链相同。
证明完成。
\end{proof}

\begin{lemma}
\label{derham-lemma-de-rham-cohomology-projective-space}
对 $0 \leq i \leq n$，德拉姆上同调
$H^{2i}_{dR}(\mathbf{P}^n_A/A)$ 是自由 $A$-模，秩为 $1$，
其基元素为 $c_1^{dR}(\mathcal{O}(1))^i$。
$\mathbf{P}^n_A$ 在 $A$ 上的德拉姆上同调在所有其他次数均为零。
\end{lemma}

\begin{proof}
考虑第 \ref{derham-section-hodge-to-de-rham} 节中的霍奇--德拉姆谱序列。
由引理 \ref{derham-lemma-hodge-cohomology-projective-space} 中对
$\mathbf{P}^n_A$ 在 $A$ 上的霍奇上同调所作的计算可知，
该谱序列在 $E_1$ 叶退化。因此
$H^{2i}_{dR}(\mathbf{P}^n_A/A)$ 是自由 $A$-模，秩为 $1$；
这对 $0 \leq i \leq n$ 成立，而其他次数为零。
注意 $c_1^{dR}(\mathcal{O}(1))^i \in H^{2i}_{dR}(\mathbf{P}^n_A/A)$
（$i = 0, \ldots, n$），且当 $i = n$ 时，该元素是
$c_1^{Hodge}(\mathcal{L})^n$ 在下列复形映射下的像：
$$
\Omega^n_{\mathbf{P}^n_A/A}[-n]
\longrightarrow
\Omega^\bullet_{\mathbf{P}^n_A/A}
$$
例如，这可由备注 \ref{derham-remark-truncations} 中的讨论得到，也可由第
\ref{derham-section-first-chern-class} 节中表示这些类的上闭链之显式描述得到。
谱序列表明，诱导映射
$$
H^n(\mathbf{P}^n_A, \Omega^n_{\mathbf{P}^n_A/A}) \longrightarrow
H^{2n}_{dR}(\mathbf{P}^n_A/A)
$$
是同构。由于 $c_1^{Hodge}(\mathcal{L})^n$ 是源的生成元
（引理 \ref{derham-lemma-hodge-cohomology-projective-space}），
可知 $c_1^{dR}(\mathcal{L})^n$ 是靶的生成元。
由杯积的 $A$-双线性性，进一步得到
$c_1^{dR}(\mathcal{L})^i$ 是
$H^{2i}_{dR}(\mathbf{P}^n_A/A)$ 的生成元；这对
$0 \leq i \leq n$ 成立。
\end{proof}









\section{光滑态射的谱序列}
\label{derham-section-relative-spectral-sequence}

\noindent
考虑概形的交换图
$$
\xymatrix{
X \ar[rr]_f \ar[rd]_p & & Y \ar[ld]^q \\
& S
}
$$
其中 $f$ 是光滑态射。于是得到局部分裂的短正合列
$$
0 \to f^*\Omega_{Y/S} \to \Omega_{X/S} \to \Omega_{X/Y} \to 0
$$
，这是《态射》，引理 \ref{morphisms-lemma-triangle-differentials-smooth} 的结论。
把它看作递降滤过 $F$；该滤过定义在 $\Omega_{X/S}$ 上，其中
$F^0\Omega_{X/S} = \Omega_{X/S}$、$F^1\Omega_{X/S} = f^*\Omega_{Y/S}$，且
$F^2\Omega_{X/S} = 0$。应用函子 $\wedge^p$，对每个 $p$ 得到诱导滤过
$$
\Omega^p_{X/S} = F^0\Omega^p_{X/S} \supset
F^1\Omega^p_{X/S} \supset
F^2\Omega^p_{X/S} \supset \ldots \supset F^{p + 1}\Omega^p_{X/S} = 0
$$
其相继商为
$$
\text{gr}^k\Omega^p_{X/S} =
F^k\Omega^p_{X/S}/F^{k + 1}\Omega^p_{X/S} =
f^*\Omega^k_{Y/S} \otimes_{\mathcal{O}_X} \Omega^{p - k}_{X/Y} =
f^{-1}\Omega^k_{Y/S} \otimes_{f^{-1}\mathcal{O}_Y} \Omega^{p - k}_{X/Y}
$$
这里 $k = 0, \ldots, p$。事实上，读者可用莱布尼茨法则验证：
$F^k\Omega^\bullet_{X/S}$ 是 $\Omega^\bullet_{X/S}$ 的子复形。
由此，$\Omega^\bullet_{X/S}$ 具有滤过复形的结构。也可以通过观察下式看出：
$$
F^k\Omega^\bullet_{X/S} =
\Im\left(\wedge :
\text{Tot}(
f^{-1}\sigma_{\geq k}\Omega^\bullet_{Y/S} \otimes_{p^{-1}\mathcal{O}_S}
\Omega^\bullet_{X/S})
\longrightarrow
\Omega^\bullet_{X/S}\right)
$$
它是 $X$ 上复形映射的像。滤过复形
$$
\Omega^\bullet_{X/S} = F^0\Omega^\bullet_{X/S} \supset
F^1\Omega^\bullet_{X/S} \supset F^2\Omega^\bullet_{X/S} \supset \ldots
$$
具有如下相伴分次部分：
$$
\text{gr}^k\Omega^\bullet_{X/S} = 
f^{-1}\Omega^k_{Y/S}[-k] \otimes_{f^{-1}\mathcal{O}_Y} \Omega^\bullet_{X/Y}
$$
这由上述讨论得到。

\begin{lemma}
\label{derham-lemma-spectral-sequence-smooth}
设 $f : X \to Y$ 是基概形 $S$ 上拟紧、拟分离且光滑的概形态射。
存在有界谱序列，其第一叶为
$$
E_1^{p, q} =
H^q(\Omega^p_{Y/S} \otimes_{\mathcal{O}_Y}^\mathbf{L} Rf_*\Omega^\bullet_{X/Y})
$$
它收敛到 $R^{p + q}f_*\Omega^\bullet_{X/S}$。
\end{lemma}

\begin{proof}
把 $\Omega^\bullet_{X/S}$ 配以上述滤过，视为滤过复形。
该谱序列就是《上同调》，引理
\ref{cohomology-lemma-relative-spectral-sequence-filtered-object} 的谱序列。
由《概形的导出范畴》，引理
\ref{perfect-lemma-cohomology-de-rham-base-change}，有
$$
Rf_*\text{gr}^k\Omega^\bullet_{X/S} =
\Omega^k_{Y/S}[-k] \otimes_{\mathcal{O}_Y}^\mathbf{L} Rf_*\Omega^\bullet_{X/Y}
$$
从而得到结论。
\end{proof}

\begin{remark}
\label{derham-remark-gauss-manin}
在引理 \ref{derham-lemma-spectral-sequence-smooth} 中，考虑上同调层
$$
\mathcal{H}^q_{dR}(X/Y) = H^q(Rf_*\Omega^\bullet_{X/Y})
$$
若 $f$ 除光滑外还为真，且 $S$ 是 $\mathbf{Q}$ 上的概形，则
$\mathcal{H}^q_{dR}(X/Y)$ 有限局部自由（在此插入将来的参考文献）。
若只假设 $\mathcal{H}^q_{dR}(X/Y)$ 是平坦
$\mathcal{O}_Y$-模，则得到（略去一个很短的论证）
$$
E_1^{p, q} =
\Omega^p_{Y/S} \otimes_{\mathcal{O}_Y} \mathcal{H}^q_{dR}(X/Y)
$$
而谱序列中的微分为映射
$$
d_1^{p, q} :
\Omega^p_{Y/S} \otimes_{\mathcal{O}_Y} \mathcal{H}^q_{dR}(X/Y)
\longrightarrow
\Omega^{p + 1}_{Y/S} \otimes_{\mathcal{O}_Y} \mathcal{H}^q_{dR}(X/Y)
$$
特别地，当 $p = 0$ 时，得到映射
$d_1^{0, q} : \mathcal{H}^q_{dR}(X/Y) \to
\Omega^1_{Y/S} \otimes_{\mathcal{O}_Y} \mathcal{H}^q_{dR}(X/Y)$
；它实际上是可积联络 $\nabla$（在此插入将来的参考文献）。
复形
$$
\mathcal{H}^q_{dR}(X/Y) \to
\Omega^1_{Y/S} \otimes_{\mathcal{O}_Y} \mathcal{H}^q_{dR}(X/Y) \to
\Omega^2_{Y/S} \otimes_{\mathcal{O}_Y} \mathcal{H}^q_{dR}(X/Y) \to \ldots
$$
的微分由 $d_1^{\bullet, q}$ 给出，它是 $\nabla$ 的德拉姆复形。
联络 $\nabla$ 称为{\it 高斯--曼宁联络}。
\end{remark}






\section{Leray--Hirsch 型定理}
\label{derham-section-leray-hirsch}

\noindent
本节证明：对光滑真态射，有时可以用下方的德拉姆上同调表示上方的德拉姆上同调。

\begin{lemma}
\label{derham-lemma-relative-global-generation-on-fibres}
设 $f : X \to Y$ 是光滑真概形态射。设 $N$ 和
$n_1, \ldots, n_N \geq 0$ 为整数，并设
$\xi_i \in H^{n_i}_{dR}(X/Y)$（$1 \leq i \leq N$）。
假设对所有点 $y \in Y$，$\xi_1, \ldots, \xi_N$ 在
$H^*_{dR}(X_y/y)$ 中的像构成 $\kappa(y)$ 上的基。则与
$$
\bigoplus\nolimits_{i = 1}^N \mathcal{O}_Y[-n_i]
\longrightarrow
Rf_*\Omega^\bullet_{X/Y}
$$
$\xi_1, \ldots, \xi_N$ 相伴的上述映射是同构。
\end{lemma}

\begin{proof}
由引理 \ref{derham-lemma-proper-smooth-de-Rham}，
$Rf_*\Omega^\bullet_{X/Y}$ 是 $D(\mathcal{O}_Y)$ 中的完美对象，
且其构造与任意基变换相容。因此，本引理的映射是
完美对象之间的映射 $a : K \to L$；这些对象位于
$D(\mathcal{O}_Y)$ 中。根据对纤维的假设，它到任一点的导出限制都是同构。
于是锥 $C$（即 $a$ 的锥）是 $D(\mathcal{O}_Y)$ 中的完美对象
（《上同调》，引理 \ref{cohomology-lemma-two-out-of-three-perfect}），
且其到任一点的导出限制为零。因此，由《更多代数》，引理
\ref{more-algebra-lemma-lift-perfect-from-residue-field}，$C$ 为零，
而 $a$ 是同构。（这里还使用《概形的导出范畴》中的引理
\ref{perfect-lemma-affine-compare-bounded} 和
\ref{perfect-lemma-perfect-affine}，以转化到代数情形。）
\end{proof}

\noindent
我们先在如下特殊情形下证明本节的主要结果。

\begin{lemma}
\label{derham-lemma-global-generation-on-fibres}
设 $f : X \to Y$ 是基概形 $S$ 上的光滑真概形态射。假设
\begin{enumerate}
\item $Y$ 和 $S$ 都是仿射的；
\item 存在整数 $N$、$n_1, \ldots, n_N \geq 0$，以及
$\xi_i \in H^{n_i}_{dR}(X/S)$（$1 \leq i \leq N$），使得对所有点
$y \in Y$，$\xi_1, \ldots, \xi_N$ 在 $H^*_{dR}(X_y/y)$ 中的像构成
$\kappa(y)$ 上的基。
\end{enumerate}
则映射
$$
\bigoplus\nolimits_{i = 1}^N H^*_{dR}(Y/S) \longrightarrow
H^*_{dR}(X/S), \quad
(a_1, \ldots, a_N) \longmapsto  \sum \xi_i \cup f^*a_i
$$
是同构。
\end{lemma}

\begin{proof}
记 $Y = \Spec(A)$ 且 $S = \Spec(R)$。在此情形下，由引理
\ref{derham-lemma-de-rham-affine}，$\Omega^\bullet_{A/R}$ 计算
$R\Gamma(Y, \Omega^\bullet_{Y/S})$。取有限仿射开覆盖
$\mathcal{U} : X = \bigcup_{i \in I} U_i$。考虑复形
$$
K^\bullet =
\text{Tot}(\check{\mathcal{C}}^\bullet(\mathcal{U}, \Omega_{X/S}^\bullet))
$$
，如《上同调》第
\ref{cohomology-section-cech-cohomology-of-complexes} 节所述。
我们汇集关于这个复形的一些事实；其中大部分可在刚才给出的参考文献中找到：
\begin{enumerate}
\item $K^\bullet$ 是一个 $R$-模复形，其各项都是 $A$-模；
\item $K^\bullet$ 表示 $R\Gamma(X, \Omega^\bullet_{X/S})$ 这个 $D(R)$ 中的对象
（《概形上同调》，引理
\ref{coherent-lemma-quasi-coherent-affine-cohomology-zero}，以及
《上同调》，引理 \ref{cohomology-lemma-cech-complex-complex-computes}）；
\item 存在自然映射 $\Omega^\bullet_{A/R} \to K^\bullet$，
它是 $R$-模复形的映射，并且逐项是 $A$-线性的，
并在上同调上诱导拉回映射
$H^*_{dR}(Y/S) \to H^*_{dR}(X/S)$；
\item $K^\bullet$ 上有一个记作 $\wedge$ 的乘法，使其成为微分分次
$R$-代数；
\item $K^\bullet$ 上的乘法诱导 $H^*_{dR}(X/S)$ 上的杯积
（《上同调》，第 \ref{cohomology-section-cup-product} 节）；
\item 滤过 $F$ 在 $\Omega^*_{X/S}$ 上，并诱导如下滤过
$$
K^\bullet =
F^0K^\bullet \supset F^1K^\bullet \supset F^2K^\bullet \supset \ldots
$$
；它由 $K^\bullet$ 的子复形构成，并满足
\begin{enumerate}
\item $F^kK^n \subset K^n$ 是 $A$-子模；
\item $F^kK^\bullet \wedge F^lK^\bullet \subset F^{k + l}K^\bullet$；
\item $\text{gr}^kK^\bullet$ 是 $A$-模复形；
\item $\text{gr}^0K^\bullet =
\text{Tot}(\check{\mathcal{C}}^\bullet(\mathcal{U}, \Omega_{X/Y}^\bullet))$
，且表示 $R\Gamma(X, \Omega^\bullet_{X/Y})$ 这个 $D(A)$ 中的对象；
\item 乘法诱导同构
$\Omega^k_{A/R}[-k] \otimes_A \text{gr}^0K^\bullet \to \text{gr}^kK^\bullet$
\end{enumerate}
\end{enumerate}
我们略去这些陈述的详细证明；请参阅引理
\ref{derham-lemma-spectral-sequence-smooth} 中构造谱序列之前的讨论。

\medskip\noindent
对每个 $i = 1, \ldots, N$，选取上闭链 $x_i \in K^{n_i}$
来表示 $\xi_i$。接着考虑复形的映射
$$
\tilde x :
M^\bullet = \bigoplus\nolimits_{i = 1, \ldots, N}
\Omega^\bullet_{A/R}[-n_i]
\longrightarrow
K^\bullet
$$
，它把 $\omega$（位于第 $i$ 个直和项）映到 $x_i \wedge \omega$。
只需证明这个映射是拟同构。我们按如下规则赋予 $M^\bullet$ 滤过复形结构：
$$
F^kM^\bullet =
\bigoplus\nolimits_{i = 1, \ldots, N}
(\sigma_{\geq k}\Omega^\bullet_{A/R})[-n_i]
$$
。在此选择下，映射 $\tilde x$ 是滤过复形的态射。注意
$\text{gr}^0M^\bullet = \bigoplus A[-n_i]$，且乘法诱导同构
$\Omega^k_{A/R}[-k] \otimes_A \text{gr}^0M^\bullet \to \text{gr}^kM^\bullet$.
由构造及引理 \ref{derham-lemma-relative-global-generation-on-fibres} 可知
$$
\text{gr}^0\tilde x :
\text{gr}^0M^\bullet \longrightarrow
\text{gr}^0K^\bullet
$$
在 $D(A)$ 中是同构。因而对所有 $k \geq 0$，得到同构
$$
\text{gr}^k \tilde x :
\text{gr}^kM^\bullet = \Omega^k_{A/R}[-k] \otimes_A \text{gr}^0M^\bullet
\longrightarrow
\Omega^k_{A/R}[-k] \otimes_A \text{gr}^0K^\bullet =
\text{gr}^kK^\bullet
$$
于 $D(A)$ 中。事实上，复形
$\text{gr}^0K^\bullet =
\text{Tot}(\check{\mathcal{C}}^\bullet(\mathcal{U}, \Omega_{X/Y}^\bullet))$
由《概形的导出范畴》，引理 \ref{perfect-lemma-K-flat}，作为
$A$-模复形是 K-平坦的。因此右边的张量积就是导出张量积；
左边直接考察即可得到同一结论。最后，取导出张量积
$\Omega^k_{A/R}[-k] \otimes_A^\mathbf{L} -$ 是 $D(A)$ 上的函子，
故将同构映到同构。对 $k$ 作归纳，推出
$$
\tilde x : M^\bullet/F^kM^\bullet \to K^\bullet/F^kK^\bullet
$$
在 $D(R)$ 中是同构，因为有短正合列
$$
0 \to F^kM^\bullet/F^{k + 1}M^\bullet \to
M^\bullet/F^{k + 1}M^\bullet \to
\text{gr}^kM^\bullet \to 0
$$
，对 $K^\bullet$ 也有类似的短正合列。由于这些滤过在任一给定次数上
都是有限的，这就证明了 $\tilde x$ 是拟同构。
\end{proof}

\begin{proposition}
\label{derham-proposition-global-generation-on-fibres}
设 $f : X \to Y$ 是基概形 $S$ 上的光滑真概形态射。设 $N$ 和
$n_1, \ldots, n_N \geq 0$ 为整数，并设
$\xi_i \in H^{n_i}_{dR}(X/S)$（$1 \leq i \leq N$）。假设对所有点
$y \in Y$，$\xi_1, \ldots, \xi_N$ 在 $H^*_{dR}(X_y/y)$ 中的像构成
$\kappa(y)$ 上的基。则映射
$$
\tilde \xi = \bigoplus \tilde \xi_i[-n_i] :
\bigoplus \Omega^\bullet_{Y/S}[-n_i]
\longrightarrow
Rf_*\Omega^\bullet_{X/S}
$$
（见证明）在 $D(Y, (Y \to S)^{-1}\mathcal{O}_S)$ 中是同构；相应地，映射
$$
\bigoplus\nolimits_{i = 1}^N H^*_{dR}(Y/S) \longrightarrow
H^*_{dR}(X/S), \quad
(a_1, \ldots, a_N) \longmapsto  \sum \xi_i \cup f^*a_i
$$
是同构。
\end{proposition}

\begin{proof}
记结构态射为 $p : X \to S$ 和 $q : Y \to S$。令
$\xi'_i : \Omega^\bullet_{X/S} \to \Omega^\bullet_{X/S}[n_i]$
为评注 \ref{derham-remark-cup-product-as-a-map} 中与 $\xi_i$ 对应的映射。记
$$
\tilde \xi_i :
\Omega^\bullet_{Y/S} \to Rf_*\Omega^\bullet_{X/S}[n_i]
$$
为 $\xi'_i$ 与典范映射
$\Omega^\bullet_{Y/S} \to Rf_*\Omega^\bullet_{X/S}$ 的复合。利用
$$
R\Gamma(Y, Rf_*\Omega^\bullet_{X/S}) = R\Gamma(X, \Omega^\bullet_{X/S})
$$
，在上同调上，$\tilde \xi_i$ 是映射 $\eta \mapsto \xi_i \cup f^*\eta$，
它从 $H^m_{dR}(Y/S)$ 映到 $H^{m + n}_{dR}(X/S)$。
此外，由于 $\xi'_i$ 的构造与到开集的限制可交换，$\tilde \xi_i$
的构造也与到开集的限制可交换。

\medskip\noindent
因此可以考虑映射
$$
\tilde \xi = \bigoplus \tilde \xi_i[-n_i] :
\bigoplus \Omega^\bullet_{Y/S}[-n_i]
\longrightarrow
Rf_*\Omega^\bullet_{X/S}
$$
。要证明该引理，只需证明它在 $D(Y, q^{-1}\mathcal{O}_S)$ 中是同构。
若能证明 $\tilde \xi$ 来自滤过复形的映射（配备适当滤过），
便可诉诸引理 \ref{derham-lemma-spectral-sequence-smooth} 的谱序列完成证明。
这比所需工作更多，因此我们改为约化到仿射情形（其证明在某种意义上
确实使用了谱序列）。

\medskip\noindent
事实上，若 $Y' \subset Y$ 是任意开集，其逆像为 $X' \subset X$，
则 $\tilde \xi|_{X'}$ 诱导映射
$$
\bigoplus\nolimits_{i = 1}^N H^*_{dR}(Y'/S) \longrightarrow
H^*_{dR}(X'/S), \quad
(a_1, \ldots, a_N) \longmapsto  \sum \xi_i|_{X'} \cup f^*a_i
$$
于 $Y'$ 上的上同调中，见上文讨论。因此，只需为 $Y$ 上的拓扑找一个基，使该命题对基中的
各成员成立（特别地，此时可以不再考虑映射 $\tilde \xi$）。这就约化到
$Y$ 和 $S$ 都是仿射的情形；该情形由引理
\ref{derham-lemma-global-generation-on-fibres} 处理，证明完成。
\end{proof}





\section{射影空间丛公式}
\label{derham-section-projective-space-bundle-formula}

\noindent
标题已说明一切。

\begin{proposition}
\label{derham-proposition-projective-space-bundle-formula}
设 $X \to S$ 是概形态射。设 $\mathcal{E}$ 是局部自由
$\mathcal{O}_X$-模，且其秩恒为 $r$。考虑态射
$p : P = \mathbf{P}(\mathcal{E}) \to X$。
则映射
$$
\bigoplus\nolimits_{i = 0, \ldots, r - 1} H^*_{dR}(X/S)
\longrightarrow
H^*_{dR}(P/S)
$$
由规则
$$
(a_0, \ldots, a_{r - 1}) \longmapsto
\sum\nolimits_{i = 0, \ldots, r - 1} c_1^{dR}(\mathcal{O}_P(1))^i \cup p^*(a_i)
$$
给出，并且是同构。
\end{proposition}

\begin{proof}
取仿射开集 $\Spec(A) \subset X$，使 $\mathcal{E}$ 限制为平凡局部自由模
$\mathcal{O}_{\Spec(A)}^{\oplus r}$。于是
$P \times_X \Spec(A) = \mathbf{P}^{r - 1}_A$。由此可见 $p$ 真且光滑，
见第 \ref{derham-section-projective-space} 节。此外，诸类
$c_1^{dR}(\mathcal{O}_P(1))^i$（$i = 0, 1, \ldots, r - 1$）限制到纤维
$X_y = \mathbf{P}^{r - 1}_y$ 后，自由生成德拉姆上同调
$H^*_{dR}(X_y/y)$ 这一 $\kappa(y)$-空间，见引理
\ref{derham-lemma-de-rham-cohomology-projective-space}。因此，我们已经验证了命题
\ref{derham-proposition-global-generation-on-fibres} 的条件，结论得证。
\end{proof}

\begin{remark}
\label{derham-remark-projective-space-bundle-formula}
在命题 \ref{derham-proposition-projective-space-bundle-formula} 的情形下，进一步得到映射
$$
\tilde \xi :
\bigoplus\nolimits_{t = 0, \ldots, r - 1}
\Omega^\bullet_{X/S}[-2t]
\longrightarrow
Rp_*\Omega^\bullet_{P/S}
$$
在 $D(X, (X \to S)^{-1}\mathcal{O}_X)$ 中是同构；这直接由应用命题
\ref{derham-proposition-global-generation-on-fibres} 得到。注意，对 $t = 0$，
相应箭头恰为典范映射
$c_{P/X} : \Omega^\bullet_{X/S} \to Rp_*\Omega^\bullet_{P/S}$
，见第 \ref{derham-section-de-rham-complex} 节。事实上，在这个特殊情形下，
还可进一步明确该映射。考虑典范映射
$$
\xi' : \Omega^\bullet_{P/S} \to \Omega^\bullet_{P/S}[2]
$$
，它是评注 \ref{derham-remark-cup-product-as-a-map} 中对应于
$c_1^{dR}(\mathcal{O}_P(1))$ 的映射。于是
$$
\xi'[2(t - 1)] \circ \ldots \circ \xi'[2] \circ \xi' : 
\Omega^\bullet_{P/S} \to \Omega^\bullet_{P/S}[2t]
$$
是评注 \ref{derham-remark-cup-product-as-a-map} 中对应于
$c_1^{dR}(\mathcal{O}_P(1))^t$ 的映射。追踪命题
\ref{derham-proposition-global-generation-on-fibres} 的证明中所作的选择，得到
$$
\tilde \xi|_{\Omega^\bullet_{X/S}[-2t]} =
Rp_*\xi'[-2] \circ \ldots \circ Rp_*\xi'[-2(t - 1)] \circ
Rp_*\xi'[-2t] \circ c_{P/X}[-2t]
$$
，这就是该同构到直和项 $\Omega^\bullet_{X/S}[-2t]$ 的限制。
由此得到下面将用到的简单推论：令
$$
M = \bigoplus\nolimits_{t = 1, \ldots, r - 1} \Omega^\bullet_{X/S}[-2t]
\quad\text{且}\quad
K = \bigoplus\nolimits_{t = 0, \ldots, r - 2} \Omega^\bullet_{X/S}[-2t]
$$
；将它们视为箭头 $\tilde \xi$ 的源中的子复形。由上文讨论形式地推出
$$
c_{P/X} \oplus
\tilde \xi|_M :
\Omega^\bullet_{X/S} \oplus M \longrightarrow
Rp_*\Omega^\bullet_{P/S}
$$
是同构，且图表
$$
\xymatrix{
K \ar[d]_{\tilde \xi|_K} \ar[r]_{\text{id}} &
M[2] \ar[d]^{(\tilde \xi|_M)[2]} \\
Rp_*\Omega^\bullet_{P/S} \ar[r]^{Rp_*\xi'} &
Rp_*\Omega^\bullet_{P/S}[2]
}
$$
交换；其中 $\text{id} : K \to M[2]$ 把对应于 $t$ 的直和项
（在 $K$ 的分解中）与对应于 $t + 1$ 的直和项（在 $M$ 的分解中）等同。
\end{remark}









\section{沿除子的对数极点}
\label{derham-section-divisor}

\noindent
设 $X \to S$ 是概形态射。设 $Y \subset X$ 是有效 Cartier 除子。
若 $X$ 沿 $Y$ 在平展局部看起来像 $Y \times \mathbf{A}^1$，
则存在复形的典范短正合列
$$
0 \to \Omega^\bullet_{X/S} \to
\Omega^\bullet_{X/S}(\log Y) \to
\Omega^\bullet_{Y/S}[-1] \to 0
$$
，它具有许多良好性质，本节将加以讨论。这个构造还有一个从正规交叉除子
出发的变体（《平展态射》，定义
\ref{etale-definition-strict-normal-crossings}），我们将在别处讨论
（此处插入未来的参考文献）。

\begin{definition}
\label{derham-definition-local-product}
设 $X \to S$ 是概形态射。设 $Y \subset X$ 是有效 Cartier 除子。
称{\it 对 $Y \subset X$（在 $S$ 上）定义了对数极点德拉姆复形}，
如果对所有 $y \in Y$ 以及局部方程 $f \in \mathcal{O}_{X, y}$
（它定义 $Y$），都有
\begin{enumerate}
\item $\mathcal{O}_{X, y} \to \Omega_{X/S, y}$，$g \mapsto g \text{d}f$
是分裂单射；
\item $\Omega^p_{X/S, y}$ 无 $f$-挠，对所有 $p$ 均如此。
\end{enumerate}
\end{definition}

\noindent
一个简单的局部计算表明，对每个 $y \in Y$，只需找到一个满足条件
(1) 和 (2) 的局部方程 $f$。

\begin{lemma}
\label{derham-lemma-log-complex}
设 $X \to S$ 是概形态射。设 $Y \subset X$ 是有效 Cartier 除子。
假设对 $Y \subset X$（在 $S$ 上）定义了对数极点德拉姆复形。
则存在复形的典范短正合列
$$
0 \to \Omega^\bullet_{X/S} \to
\Omega^\bullet_{X/S}(\log Y) \to
\Omega^\bullet_{Y/S}[-1] \to 0
$$
\end{lemma}

\begin{proof}
我们的假设是：对每个 $y \in Y$ 以及局部方程
$f \in \mathcal{O}_{X, y}$（它定义 $Y$），都有
$$
\Omega_{X/S, y} = \mathcal{O}_{X, y}\text{d}f \oplus M
\quad\text{且}\quad
\Omega^p_{X/S, y} = \wedge^{p - 1}(M)\text{d}f \oplus \wedge^p(M)
$$
，其中某个模为 $M$，并且对 $f$ 而言，$\wedge^p(M)$ 无挠。因而
$$
\Omega^p_{Y/S, y} = \wedge^p(M/fM) = \wedge^p(M)/f\wedge^p(M)
$$
。下文将默认使用这些事实。特别地，层 $\Omega^p_{X/S}$ 没有支撑在
$Y$ 上的非零局部截面，并且有典范包含
$$
\Omega^p_{X/S} \subset \Omega^p_{X/S}(Y)
$$
，见《平坦性的进一步结果》，第 \ref{flat-section-eta} 节。令
$U = \Spec(A)$ 为仿射开子概形，使得 $Y \cap U = V(f)$，其中
$f \in A$ 是非零因子。考虑 $\mathcal{O}_U$-子模，它位于
$\Omega^p_{X/S}(Y)|_U$ 中，
它由 $\Omega^p_{X/S}|_U$ 和
$\text{d}\log(f) \wedge \Omega^{p - 1}_{X/S}$ 生成，其中
$\text{d}\log(f) = f^{-1}\text{d}(f)$。这与 $f$ 的选择无关：
$Y$ 在 $U$ 上的理想的另一个生成元等于 $uf$，其中 $u \in A$ 是单位，且
$$
\text{d}\log(uf) - \text{d}\log(f) = \text{d}\log(u) = u^{-1}\text{d}u
$$
是 $\Omega_{X/S}$ 在 $U$ 上的截面。这些局部层可粘合成拟凝聚子模
$$
\Omega^p_{X/S} \subset \Omega^p_{X/S}(\log Y) \subset \Omega^p_{X/S}(Y)
$$
。约定把 $\Omega^p_{Y/S}$ 看作拟凝聚 $\mathcal{O}_X$-模。
存在唯一的满射 $\mathcal{O}_X$-线性映射
$$
\text{Res} : \Omega^p_{X/S}(\log Y) \to \Omega^{p - 1}_{Y/S}
$$
，其定义规则为
$$
\text{Res}(\eta' + \text{d}\log(f) \wedge \eta) = \eta|_{Y \cap U}
$$
；这里 $U$ 遍历如上的开集，且
$\eta' \in \Omega^p_{X/S}(U)$ 和 $\eta \in \Omega^{p - 1}_{X/S}(U)$ 任取。
若微分形式 $\eta$ 在 $U$ 上，并且限制到 $Y \cap U$ 后为零，则有
$\eta = \text{d}f \wedge \eta' + f\eta''$，其中 $\eta'$ 和 $\eta''$
是 $U$ 上的某些微分形式。因此得到短正合列
$$
0 \to \Omega^p_{X/S} \to \Omega^p_{X/S}(\log Y) \to \Omega^{p - 1}_{Y/S} \to 0
$$
，对所有 $p$ 均成立。还需定义微分
$\Omega^p_{X/S}(\log Y) \to \Omega^{p + 1}_{X/S}(\log Y)$。
在子层 $\Omega^p_{X/S}$ 上，使用 $X$ 在 $S$ 上的德拉姆复形的微分。
最后定义
$\text{d}(\text{d}\log(f) \wedge \eta) = -\text{d}\log(f) \wedge \text{d}\eta$.
这个符号由 Leibniz 法则（在 $\Omega^\bullet_{X/S}$ 上）所迫使，
并与 $\Omega^\bullet_{Y/S}[-1]$ 上的微分相容；按照《同调》，定义
\ref{homology-definition-shift-cochain} 中的符号约定，后者正是
$-\text{d}_{Y/S}$。这样便得到引理所述的复形短正合列。
\end{proof}

\begin{definition}
\label{derham-definition-log-complex}
设 $X \to S$ 是概形态射。设 $Y \subset X$ 是有效 Cartier 除子。
假设对 $Y \subset X$（在 $S$ 上）定义了对数极点德拉姆复形。则复形
$$
\Omega^\bullet_{X/S}(\log Y)
$$
（构造见引理 \ref{derham-lemma-log-complex}）称为
{\it $Y \subset X$（在 $S$ 上）的对数极点德拉姆复形}。
\end{definition}

\noindent
这个复形具有许多良好性质。

\begin{lemma}
\label{derham-lemma-multiplication-log}
设 $p : X \to S$ 是概形态射。设 $Y \subset X$ 是有效 Cartier 除子。
假设对 $Y \subset X$（在 $S$ 上）定义了对数极点德拉姆复形。
\begin{enumerate}
\item 映射
$\wedge : \Omega^p_{X/S} \times \Omega^q_{X/S} \to \Omega^{p + q}_{X/S}$
唯一延拓为 $\mathcal{O}_X$-双线性映射
$$
\wedge : \Omega^p_{X/S}(\log Y) \times \Omega^q_{X/S}(\log Y)
\to \Omega^{p + q}_{X/S}(\log Y)
$$
，并满足 Leibniz 法则
$
\text{d}(\omega \wedge \eta) = \text{d}(\omega) \wedge \eta +
(-1)^{\deg(\omega)} \omega \wedge \text{d}(\eta)$,
\item 对于 (1) 中的乘法，映射
$\Omega^\bullet_{X/S} \to \Omega^\bullet_{X/S}(\log(Y)$
是微分分次 $\mathcal{O}_S$-代数的同态；
\item 经由 (1) 中的映射，有 $\Omega^p_{X/S}(\log Y) =
\wedge^p(\Omega^1_{X/S}(\log Y))$；
\item 映射
$\text{Res} : \Omega^\bullet_{X/S}(\log Y) \to \Omega^\bullet_{Y/S}[-1]$
满足
$$
\text{Res}(\omega \wedge \eta) = \text{Res}(\omega) \wedge \eta|_Y
$$
，其中 $\omega$ 是 $\Omega^p_{X/S}(\log Y)$ 的局部截面，而 $\eta$
是 $\Omega^q_{X/S}$ 的局部截面。
\end{enumerate}
\end{lemma}

\begin{proof}
这由引理 \ref{derham-lemma-log-complex} 的证明中该复形的局部构造直接计算得到。
细节从略。
\end{proof}

\noindent
考虑交换图表
$$
\xymatrix{
X' \ar[r]_f \ar[d] & X \ar[d] \\
S' \ar[r] & S
}
$$
，其中各对象都是概形。设 $Y \subset X$ 是有效 Cartier 除子，且其拉回
$Y' = f^*Y$ 有定义（《除子》，定义
\ref{divisors-definition-pullback-effective-Cartier-divisor}）。
假设对 $Y \subset X$（在 $S$ 上）定义了对数极点德拉姆复形，并且
对 $Y' \subset X'$（在 $S'$ 上）定义了对数极点德拉姆复形。
在这种情况下，得到复形短正合列之间的映射
$$
\xymatrix{
0 \ar[r] &
f^{-1}\Omega^\bullet_{X/S} \ar[r] \ar[d] &
f^{-1}\Omega^\bullet_{X/S}(\log Y) \ar[r] \ar[d] &
f^{-1}\Omega^\bullet_{Y/S}[-1] \ar[r] \ar[d] &
0 \\
0 \ar[r] &
\Omega^\bullet_{X'/S'} \ar[r] &
\Omega^\bullet_{X'/S'}(\log Y') \ar[r] &
\Omega^\bullet_{Y'/S'}[-1] \ar[r] &
0
}
$$
。线性化后，对每个 $p$ 得到线性映射
$f^*\Omega^p_{X/S}(\log Y) \to \Omega^p_{X'/S'}(\log Y')$.

\begin{lemma}
\label{derham-lemma-gysin-via-log-complex}
设 $f : X \to S$ 是概形态射。设 $Y \subset X$ 是有效 Cartier 除子。
假设对 $Y \subset X$（在 $S$ 上）定义了对数极点德拉姆复形。记
$$
\delta : \Omega^\bullet_{Y/S} \to \Omega^\bullet_{X/S}[2]
$$
为 $D(X, f^{-1}\mathcal{O}_S)$ 中来自引理
\ref{derham-lemma-log-complex} 的短正合列的“边缘”映射。记
$$
\xi' : \Omega^\bullet_{X/S} \to \Omega^\bullet_{X/S}[2]
$$
为 $D(X, f^{-1}\mathcal{O}_S)$ 中评注
\ref{derham-remark-cup-product-as-a-map} 的映射，它对应于
$\xi = c_1^{dR}(\mathcal{O}_X(-Y))$。记
$$
\zeta' : \Omega^\bullet_{Y/S} \to \Omega^\bullet_{Y/S}[2]
$$
为 $D(Y, f|_Y^{-1}\mathcal{O}_S)$ 中评注
\ref{derham-remark-cup-product-as-a-map} 的映射，它对应于
$\zeta = c_1^{dR}(\mathcal{O}_X(-Y)|_Y)$。则图表
$$
\xymatrix{
\Omega^\bullet_{X/S} \ar[d]_{\xi'} \ar[r] &
\Omega^\bullet_{Y/S} \ar[d]^{\zeta'} \ar[ld]_\delta \\
\Omega^\bullet_{X/S}[2] \ar[r] &
\Omega^\bullet_{Y/S}[2]
}
$$
在 $D(X, f^{-1}\mathcal{O}_S)$ 中交换。
\end{lemma}

\begin{proof}
更准确地说，我们把 $\delta$ 定义为与如下移位短正合列相伴的边缘映射：
$$
0 \to \Omega^\bullet_{X/S}[1] \to
\Omega^\bullet_{X/S}(\log Y)[1] \to
\Omega^\bullet_{Y/S} \to 0
$$
。只需证明每个三角形都交换。置
$\mathcal{L} = \mathcal{O}_X(-Y)$。记 $\pi : L \to X$ 为满足
$\pi_*\mathcal{O}_L = \bigoplus_{n \geq 0} \mathcal{L}^{\otimes n}$ 的线丛。

\medskip\noindent
左上三角形的交换性。由引理
\ref{derham-lemma-the-complex-for-L-star-gives-chern-class}，映射 $\xi'$
是引理 \ref{derham-lemma-the-complex-for-L-star} 所给三角形的边缘映射。
由函子性，只需证明存在短正合列的态射
$$
\xymatrix{
0 \ar[r] &
\Omega^\bullet_{X/S}[1] \ar[r] \ar[d] &
\Omega^\bullet_{L^\star/S, 0}[1] \ar[r] \ar[d] &
\Omega^\bullet_{X/S} \ar[r] \ar[d] &
0 \\
0 \ar[r] &
\Omega^\bullet_{X/S}[1] \ar[r] &
\Omega^\bullet_{X/S}(\log Y)[1] \ar[r] &
\Omega^\bullet_{Y/S} \ar[r] &
0
}
$$
，其中左、右竖直箭头是显然的。可用如下规则定义中间的竖直箭头：
$$
\omega' + \text{d}\log(s) \wedge \omega \longmapsto
\omega' + \text{d}\log(f) \wedge \omega
$$
。其中 $\omega', \omega$ 是 $\Omega^\bullet_{X/S}$ 的局部截面，
$s$ 是 $\mathcal{L}$ 的局部生成元，而 $f \in \mathcal{O}_X(-Y)$
是相应的截面；它属于 $Y$ 在 $X$ 中的理想层。由于引理
\ref{derham-lemma-the-complex-for-L-star} 与 \ref{derham-lemma-log-complex}
中映射的构造完全吻合，此定义可行。

\medskip\noindent
右下三角形的交换性。记 $\overline{L}$ 为 $L$ 到 $Y$ 的限制。
由引理 \ref{derham-lemma-the-complex-for-L-star-gives-chern-class}，映射
$\zeta'$ 是引理 \ref{derham-lemma-the-complex-for-L-star} 中使用线丛
$\overline{L}$（定义在 $Y$ 上）所给三角形的边缘映射。由函子性，
只需证明存在短正合列的态射
$$
\xymatrix{
0 \ar[r] &
\Omega^\bullet_{X/S}[1] \ar[r] \ar[d] &
\Omega^\bullet_{X/S}(\log Y)[1] \ar[r] \ar[d] &
\Omega^\bullet_{Y/S} \ar[r] \ar[d] &
0 \\
0 \ar[r] &
\Omega^\bullet_{Y/S}[1] \ar[r] &
\Omega^\bullet_{\overline{L}^\star/S, 0}[1] \ar[r] &
\Omega^\bullet_{Y/S} \ar[r] &
0 \\
}
$$
，其中左、右竖直箭头是显然的。可用如下规则定义中间的竖直箭头：
$$
\omega' + \text{d}\log(f) \wedge \omega \longmapsto
\omega'|_Y + \text{d}\log(\overline{s}) \wedge \omega|_Y
$$
。其中 $\omega', \omega$ 是 $\Omega^\bullet_{X/S}$ 的局部截面，
$f$ 是 $\mathcal{O}_X(-Y)$ 的局部生成元，并被视为 $X$ 上的函数；
$\overline{s}$ 是 $f|_Y$，并被视为
$\mathcal{L}|_Y = \mathcal{O}_X(-Y)|_Y$ 的截面。由于引理
\ref{derham-lemma-the-complex-for-L-star} 与 \ref{derham-lemma-log-complex}
中映射的构造完全吻合，此定义可行。
\end{proof}

\begin{lemma}
\label{derham-lemma-log-complex-consequence}
设 $X \to S$ 是概形态射。设 $Y \subset X$ 是有效 Cartier 除子。
假设对 $Y \subset X$（在 $S$ 上）定义了对数极点德拉姆复形。
设 $b \in H^m_{dR}(X/S)$ 是德拉姆上同调类，且其到 $Y$ 的限制为零。
则 $c_1^{dR}(\mathcal{O}_X(Y)) \cup b = 0$ 于 $H^{m + 2}_{dR}(X/S)$ 中。
\end{lemma}

\begin{proof}
这直接由引理 \ref{derham-lemma-gysin-via-log-complex} 得到。事实上，
$$
c_1^{dR}(\mathcal{O}_X(Y)) \cup b =
-c_1^{dR}(\mathcal{O}_X(-Y)) \cup b = -\xi'(b) = -\delta(b|_Y) = 0
$$
，正如所需。第二个等号见评注 \ref{derham-remark-cup-product-as-a-map}。
\end{proof}

\begin{lemma}
\label{derham-lemma-check-log-smooth}
设 $X \to T \to S$ 是概形态射。设 $Y \subset X$ 是有效 Cartier 除子。
若 $X \to T$ 与 $Y \to T$ 都光滑，则对 $Y \subset X$（在 $S$ 上）
定义了对数极点德拉姆复形。
\end{lemma}

\begin{proof}
设 $y \in Y$ 是一点。由《态射的进一步结果》，引理
\ref{more-morphisms-lemma-etale-local-structure}，存在整数 $0 \geq m$
及交换图表
$$
\xymatrix{
Y \ar[d] &
V \ar[l] \ar[d] \ar[r] &
\mathbf{A}^m_T
\ar[d]^{(a_1, \ldots, a_m) \mapsto (a_1, \ldots, a_m, 0)} \\
X &
U \ar[l] \ar[r]^-\pi &
\mathbf{A}^{m + 1}_T
}
$$
，其中 $U \subset X$ 是开集，$V = Y \cap U$，$\pi$ 是平展的，
$V = \pi^{-1}(\mathbf{A}^m_T)$，且 $y \in V$。记
$z \in \mathbf{A}^m_T$ 为 $y$ 的像。于是
$$
\Omega^p_{X/S, y} = \Omega^p_{\mathbf{A}^{m + 1}_T/S, z}
\otimes_{\mathcal{O}_{\mathbf{A}^{m + 1}_T, z}} \mathcal{O}_{X, x}
$$
，由引理 \ref{derham-lemma-etale}。记 $x_1, \ldots, x_{m + 1}$ 为
$\mathbf{A}^{m + 1}_T$ 上的坐标函数。由于定义
\ref{derham-definition-local-product} 中的条件 (1) 和 (2) 不依赖于局部坐标的选择，
只需在 $f$ 为 $x_{m + 1}$ 经平坦局部环同态
$\mathcal{O}_{\mathbf{A}^{m + 1}_T, z} \to \mathcal{O}_{X, x}$
所得的像时检验条件 (1) 和 (2)。这样可见，只需对
$\mathbf{A}^m_T \subset \mathbf{A}^{m + 1}_T$ 及点 $z$ 检验条件 (1) 和 (2)。
为证明这一情形，可以假设 $S = \Spec(A)$ 和 $T = \Spec(B)$ 都是仿射的。
设 $A \to B$ 是态射 $T \to S$ 所对应的环同态，并置
$P = B[x_1, \ldots, x_{m + 1}]$，从而
$\mathbf{A}^{m + 1}_T = \Spec(P)$。有
$$
\Omega_{P/A} = \Omega_{B/A} \otimes_B P \oplus
\bigoplus\nolimits_{j = 1, \ldots, m} P \text{d}x_j \oplus
P \text{d}x_{m + 1}
$$
。因此映射 $P \to \Omega_{P/A}$，$g \mapsto g \text{d}x_{m + 1}$
是分裂单射，且 $x_{m + 1}$ 在 $\Omega^p_{P/A}$ 上是非零因子，
对所有 $p \geq 0$ 均如此。关于与 $z$ 对应的素理想局部化，证明即完成。
\end{proof}

\begin{remark}
\label{derham-remark-check-log-completion-1}
设 $S$ 是局部 Noether 概形。设 $X$ 在 $S$ 上局部有限型。
设 $Y \subset X$ 是有效 Cartier 除子。若映射
$$
\mathcal{O}_{X, y}^\wedge \longrightarrow \mathcal{O}_{Y, y}^\wedge
$$
对所有 $y \in Y$ 都有截面，则对 $Y \subset X$（在 $S$ 上）
定义了对数极点德拉姆复形。若将来需要此结果，我们会给出精确陈述，
并在此补上证明。
\end{remark}

\begin{remark}
\label{derham-remark-check-log-completion-2}
设 $S$ 是局部 Noether 概形。设 $X$ 在 $S$ 上局部有限型。
设 $Y \subset X$ 是有效 Cartier 除子。若对每个 $y \in Y$，
都能找到 $S$ 上概形的图表
$$
X \xleftarrow{\varphi} U \xrightarrow{\psi} V
$$
，其中 $\varphi$ 平展，且
$\psi|_{\varphi^{-1}(Y)} : \varphi^{-1}(Y) \to V$ 平展，
则对 $Y \subset X$（在 $S$ 上）定义了对数极点德拉姆复形。
一个特殊情形是序对 $(X, Y)$ 在平展局部看起来像
$(V \times \mathbf{A}^1, V \times \{0\})$。若将来需要此结果，
我们会给出精确陈述，并在此补上证明。
\end{remark}

















\section{计算}
\label{derham-section-calculations}

\noindent
本节计算一个标准吹起的若干霍奇上同调群和德拉姆上同调群。

\medskip\noindent
固定环 $R$，并置 $S = \Spec(R)$。固定整数 $0 \leq m$ 和 $1 \leq n$。
考虑闭浸入
$$
Z = \mathbf{A}^m_S \longrightarrow \mathbf{A}^{m + n}_S = X,\quad
(a_1, \ldots, a_m) \mapsto (a_1, \ldots, a_m, 0, \ldots 0).
$$
我们要考虑吹起 $L$：它是 $X$ 沿闭子概形 $Z$ 的吹起。写作
$$
X =
\mathbf{A}^{m + n}_S =
\Spec(A)
\quad\text{其中}\quad
A = R[x_1, \ldots, x_m, y_1, \ldots, y_n]
$$
。令 $A = R[x_1, \ldots, x_m, y_1, \ldots, y_n]$ 成为分次 $R$-代数，
其中规定 $\deg(x_i) = 0$ 且 $\deg(y_j) = 1$。对于这个分次，有
$$
P =
\text{Proj}(A) =
\mathbf{A}^m_S \times_S \mathbf{P}^{n - 1}_S =
Z \times_S \mathbf{P}^{n - 1}_S =
\mathbf{P}^{n - 1}_Z
$$
。注意，截出 $Z$（在 $X$ 中）的理想是 $A_+$。因此 $L$ 是 Rees 代数
$$
A \oplus A_+ \oplus (A_+)^2 \oplus \ldots =
\bigoplus\nolimits_{d \geq 0} A_{\geq d}
$$
的 Proj。因此 $L$ 是《态射的进一步结果》第
\ref{more-morphisms-section-proj-spec} 节中更一般地研究的现象的一个例子；
我们将不再说明而直接使用那里得到的观察。特别地，有交换图表
$$
\xymatrix{
P \ar[r]_0 \ar[d]_p &
L \ar[r]_-\pi \ar[d]^b &
P \ar[d]^p \\
Z \ar[r]^i &
X \ar[r] &
Z
}
$$
，使得 $\pi : L \to P$ 是 $P = Z \times_S \mathbf{P}^{n - 1}_S$
上的线丛，其零截面为 $0$；它的像 $E = 0(P) \subset L$
是吹起 $b$ 的例外除子。

\begin{lemma}
\label{derham-lemma-comparison}
对 $a \geq 0$，有
\begin{enumerate}
\item 映射 $\Omega^a_{X/S} \to b_*\Omega^a_{L/S}$ 是同构；
\item 映射 $\Omega^a_{Z/S} \to p_*\Omega^a_{P/S}$ 是同构；
\item 映射 $Rb_*\Omega^a_{L/S} \to i_*Rp_*\Omega^a_{P/S}$
在次数 $\geq 1$ 的上同调层上是同构。
\end{enumerate}
\end{lemma}

\begin{proof}
先证明 (2)。由于
$P = Z \times_S \mathbf{P}^{n - 1}_S$
，有
$$
\Omega^a_{P/S} = \bigoplus\nolimits_{a = r + s}
\text{pr}_1^*\Omega^r_{Z/S} \otimes
\text{pr}_2^*\Omega^s_{\mathbf{P}^{n - 1}_S/S}
$$
。回忆 $p = \text{pr}_1$，由投影公式（《上同调》，引理
\ref{cohomology-lemma-projection-formula}），得到
$$
p_*\Omega^a_{P/S} = \bigoplus\nolimits_{a = r + s}
\Omega^r_{Z/S} \otimes
\text{pr}_{1, *}\text{pr}_2^*\Omega^s_{\mathbf{P}^{n - 1}_S/S}
$$
。由第 \ref{derham-section-projective-space} 节中的计算，尤其是引理
\ref{derham-lemma-hodge-cohomology-projective-space} 的证明，有
$\text{pr}_{1, *}\text{pr}_2^*\Omega^s_{\mathbf{P}^{n - 1}_S/S} = 0$；
唯一的例外是 $s = 0$，此时得到
$\text{pr}_{1, *}\mathcal{O}_P = \mathcal{O}_Z$。这证明了 (2)。

\medskip\noindent
由第 \ref{derham-section-line-bundle} 节中的材料，尤其是引理
\ref{derham-lemma-push-omega-a}，有
$\pi_*\Omega^a_{L/S} = \Omega^a_{P/S} \oplus
\bigoplus_{k \geq 1} \Omega^a_{L/S, k}$.
由于上图中的复合 $\pi \circ 0$ 是 $P$ 上的恒等态射，为证明 (3)，
只需证明 $\Omega^a_{L/S, k}$ 对 $k > 0$ 有消失的高阶上同调。
由引理 \ref{derham-lemma-the-complex-for-L-star} 和 \ref{derham-lemma-push-omega-a}，
存在短正合列
$$
0 \to \Omega^a_{P/S} \otimes \mathcal{O}_P(k)
\to \Omega^a_{L/S, k} \to
\Omega^{a - 1}_{P/S} \otimes \mathcal{O}_P(k) \to 0
$$
，其中 $\Omega^{a - 1}_{P/S} = 0$，若 $a = 0$。由于
$P = Z \times_S \mathbf{P}^{n - 1}_S$，有
$$
\Omega^a_{P/S} = \bigoplus\nolimits_{i + j = a}
\Omega^i_{Z/S} \boxtimes \Omega^j_{\mathbf{P}^{n - 1}_S/S}
$$
，由引理 \ref{derham-lemma-de-rham-complex-product}。由于 $\Omega^i_{Z/S}$
是有限秩自由模，只需证明
$\mathcal{O}_Z \boxtimes \Omega^j_{\mathbf{P}^{n - 1}_S/S}(k)$
对 $k > 0$ 的高阶上同调为零。这由引理
\ref{derham-lemma-twisted-hodge-cohomology-projective-space} 应用于
$P = Z \times_S \mathbf{P}^{n - 1}_S = \mathbf{P}^{n - 1}_Z$ 得到，
从而 (3) 的证明完成。

\medskip\noindent
还需证明 (1)。若 $n = 1$，则我们吹起的是一个有效 Cartier 除子，
$b$ 是同构，故 (1) 成立。若 $n > 1$，则复合
$$
\Gamma(X, \Omega^a_{X/S})
\to
\Gamma(L, \Omega^a_{L/S})
\to
\Gamma(L \setminus E, \Omega^a_{L/S})
=
\Gamma(X \setminus Z, \Omega^a_{X/S})
$$
是同构，因为 $\Omega^a_{X/S}$ 是有限自由的（略去一个小细节）。
因此 (1) 失败的唯一可能是：$\Gamma(L, \Omega^a_{L/S})$ 中存在非零元素，
它在 $E = 0(P)$ 外消失。由于 $L$ 是 $P$ 上的线丛，零截面为
$0 : P \to L$，只需证明：在线丛上，微分形式层不存在这样的非零截面，
即它在零截面外恒为零。读者可以通过局部计算（这是较好的方法），
或利用 $\Omega_{L/S, k} \subset \Omega_{L^\star/S, k}$（见上述参考）
看出这一点。
\end{proof}

\noindent
此时我们建议读者跳到下一节。

\begin{lemma}
\label{derham-lemma-comparison-bis}
对 $a \geq 0$，存在典范映射
$$
b^*\Omega^a_{X/S} \longrightarrow
\Omega^a_{L/S} \longrightarrow
b^*\Omega^a_{X/S} \otimes_{\mathcal{O}_L} \mathcal{O}_L((n - 1)E)
$$
，其复合由包含
$\mathcal{O}_L \subset \mathcal{O}_L((n - 1)E)$.
诱导。
\end{lemma}

\begin{proof}
显示公式中的第一个箭头见第 \ref{derham-section-de-rham-complex} 节。
为得到第二个箭头，必须证明：若把 $\Omega^a_{L/S}$ 的局部截面视为
$b^*\Omega^a_{X/S}$ 的“亚纯截面”，则它的极点阶至多为
$n - 1$（沿 $E$）。为此，在 $L$ 上的仿射局部坐标图中工作。回忆 $L$ 被仿射吹起代数
$A[\frac{I}{y_i}]$ 的谱覆盖，其中 $I = A_{+}$ 是由
$y_1, \ldots, y_n$ 生成的理想。见《代数》第
\ref{algebra-section-blow-up} 节及《除子》，引理
\ref{divisors-lemma-blowing-up-affine}。由对称性，只需在对应于 $i = 1$
的坐标图上工作。此时
$$
A[\frac{I}{y_1}] = R[x_1, \ldots, x_m, y_1, t_2, \ldots, t_n]
$$
，其中 $t_i = y_i/y_1$，见《代数的进一步结果》，引理
\ref{more-algebra-lemma-blowup-regular-sequence}。因此模 $\Omega^1_{L/S}$
在相应仿射开集上由
$\text{d}x_1, \ldots, \text{d}x_m$、$\text{d}y_1$ 和
$\text{d}t_1, \ldots, \text{d}t_n$ 自由生成。当然，这些生成元中的前
$m + 1$ 个来自 $b^*\Omega^1_{X/S}$；对余下的 $n - 1$ 个，有
$$
\text{d}t_j =
\text{d}\frac{y_j}{y_1} =
\frac{1}{y_1}\text{d}y_j - \frac{y_j}{y_1^2}\text{d}y_1 =
\frac{\text{d}y_j - t_j \text{d}y_1}{y_1}
$$
，它有 $1$ 阶极点（沿 $E$），因为在此坐标图上 $E$ 由 $y_1$ 截出。
由于这些元素中任取 $a$ 个的外积给出 $\Omega^a_{L/S}$ 在此坐标图上的一组基，
且其中至多涉及 $n - 1$ 个 $\text{d}t_j$，证明完成。
\end{proof}

\begin{lemma}
\label{derham-lemma-blowup-twist-same-cohomology}
设 $E = 0(P)$ 是吹起 $b$ 的例外除子。对任意局部自由
$\mathcal{O}_X$-模 $\mathcal{E}$ 以及 $0 \leq i \leq n - 1$，映射
$$
\mathcal{E}
\longrightarrow
Rb_*(b^*\mathcal{E} \otimes_{\mathcal{O}_L} \mathcal{O}_L(iE))
$$
在 $D(\mathcal{O}_X)$ 中是同构。
\end{lemma}

\begin{proof}
由投影公式，只需对 $\mathcal{E} = \mathcal{O}_X$ 证明此结论，
见《上同调》，引理 \ref{cohomology-lemma-projection-formula}。
由于 $X$ 是仿射的，只需证明映射
$$
H^0(X, \mathcal{O}_X) \to
H^0(L, \mathcal{O}_L) \to
H^0(L, \mathcal{O}_L(iE))
$$
都是同构，并且 $H^j(X, \mathcal{O}_L(iE)) = 0$，其中 $j > 0$ 且
$0 \leq i \leq n - 1$；见《概形上同调》，引理
\ref{coherent-lemma-quasi-coherence-higher-direct-images-application}。
由于 $\pi$ 是仿射的，取 $\pi_*$ 后即可计算整体截面和上同调，
见《概形上同调》，引理 \ref{coherent-lemma-relative-affine-cohomology}。
若 $n = 1$，则 $L \to X$ 是同构且 $i = 0$，所以第一个陈述成立。
若 $n > 1$，则考虑复合
$$
H^0(X, \mathcal{O}_X) \to H^0(L, \mathcal{O}_L) \to
H^0(L, \mathcal{O}_L(iE)) \to H^0(L \setminus E, \mathcal{O}_L) =
H^0(X \setminus Z, \mathcal{O}_X)
$$
。由于
$H^0(X \setminus Z, \mathcal{O}_X) = H^0(X, \mathcal{O}_X)$
在此情形成立，因为 $Z$ 的余维为 $n \geq 2$（在 $X$ 中；细节从略），
从而第一个陈述成立。对于第二个陈述，回忆
$\mathcal{O}_L(E) = \mathcal{O}_L(-1)$，见《除子》，引理
\ref{divisors-lemma-blowing-up-gives-effective-Cartier-divisor}。因此
$$
\pi_*\mathcal{O}_L(iE) =
\pi_*\mathcal{O}_L(-i) =
\bigoplus\nolimits_{k \geq -i} \mathcal{O}_P(k)
$$
，如《态射的进一步结果》第 \ref{more-morphisms-section-proj-spec} 节所述。
最后，由 $P = \mathbf{P}^{n - 1}_Z$ 上结构层扭子的上同调消失即可得出结论；
该消失见《概形上同调》，引理
\ref{coherent-lemma-cohomology-projective-space-over-ring}。
\end{proof}























\section{吹起与德拉姆上同调}
\label{derham-section-blowing-up}

\noindent
固定基概形 $S$、光滑态射 $X \to S$，以及闭子概形
$Z \subset X$；它在 $S$ 上也光滑。记 $b : X' \to X$ 为 $X$ 沿 $Z$ 的吹起。
记 $E \subset X'$ 为例外除子。图示为
\begin{equation}
\label{derham-equation-blowup}
\vcenter{
\xymatrix{
E \ar[r]_j \ar[d]_p & X' \ar[d]^b \\
Z \ar[r]^i & X
}
}
\end{equation}
本节的目标是证明映射
$b^* : H_{dR}^*(X/S) \longrightarrow H_{dR}^*(X'/S)$
是单射（其实还可得到更多结论）。

\begin{lemma}
\label{derham-lemma-comparison-general}
采用《态射的进一步结果》，引理 \ref{more-morphisms-lemma-blowup} 中的记号。
对 $a \geq 0$，有
\begin{enumerate}
\item 映射 $\Omega^a_{X/S} \to b_*\Omega^a_{X'/S}$ 是同构；
\item 映射 $\Omega^a_{Z/S} \to p_*\Omega^a_{E/S}$ 是同构；
\item 映射 $Rb_*\Omega^a_{X'/S} \to i_*Rp_*\Omega^a_{E/S}$
在次数 $\geq 1$ 的上同调层上是同构。
\end{enumerate}
\end{lemma}

\begin{proof}
设 $\epsilon : X_1 \to X$ 是满平展态射。记
$i_1 : Z_1 \to X_1$、$b_1 : X'_1 \to X_1$、$E_1 \subset X'_1$ 以及
$p_1 : E_1 \to Z_1$ 为《态射的进一步结果》，引理
\ref{more-morphisms-lemma-blowup} 中所考虑对象的基变换。注意，$i_1$
是 $S$ 上光滑概形之间的闭浸入，并且由《除子》，引理
\ref{divisors-lemma-flat-base-change-blowing-up}，$b_1$ 是以 $Z_1$
为中心的吹起。假设能对态射 $b_1$、$p_1$ 和 $i_1$ 证明 (1)、(2)、(3)。
那么由引理 \ref{derham-lemma-etale}，(1)、(2)、(3) 中的映射经 $\epsilon$
拉回后都是同构。由于 $\epsilon$ 是满平坦态射，结论随即成立。
因此，由《态射的进一步结果》，引理
\ref{more-morphisms-lemma-etale-local-structure}，可以在平展局部工作，
并假设处于第 \ref{derham-section-calculations} 节讨论的情形。此时本引理就是
引理 \ref{derham-lemma-comparison}。
\end{proof}

\begin{lemma}
\label{derham-lemma-distinguished-triangle-blowup}
采用《态射的进一步结果》，引理 \ref{more-morphisms-lemma-blowup} 中的记号，
并记结构态射为 $f : X \to S$。存在典范可区别三角
$$
\Omega^\bullet_{X/S} \to
Rb_*(\Omega^\bullet_{X'/S}) \oplus i_*\Omega^\bullet_{Z/S} \to
i_*Rp_*(\Omega^\bullet_{E/S}) \to
\Omega^\bullet_{X/S}[1]
$$
于 $D(X, f^{-1}\mathcal{O}_S)$ 中，其中四个映射
$$
\begin{matrix}
\Omega^\bullet_{X/S} & \to & Rb_*(\Omega^\bullet_{X'/S}), \\
\Omega^\bullet_{X/S} & \to & i_*\Omega^\bullet_{Z/S}, \\
Rb_*(\Omega^\bullet_{X'/S}) & \to & i_*Rp_*(\Omega^\bullet_{E/S}), \\
i_*\Omega^\bullet_{Z/S} & \to & i_*Rp_*(\Omega^\bullet_{E/S})
\end{matrix}
$$
都是典范映射（第 \ref{derham-section-de-rham-complex} 节），
但其中一个的符号反转。
\end{lemma}

\begin{proof}
选取可区别三角
$$
C \to Rb_*\Omega^\bullet_{X'/S} \oplus i_*\Omega^\bullet_{Z/S}
\to i_*Rp_*\Omega^\bullet_{E/S} \to C[1]
$$
于 $D(X, f^{-1}\mathcal{O}_S)$ 中。只需证明
$\Omega^\bullet_{X/S}$ 与 $C$ 同构，并且该同构与典范映射相容。
由三角范畴公理，存在可区别三角之间的映射
$$
\xymatrix{
C' \ar[r] \ar[d] &
b_*\Omega^\bullet_{X'/S} \oplus i_*\Omega^\bullet_{Z/S} \ar[r] \ar[d] &
i_*p_*\Omega^\bullet_{E/S} \ar[r] \ar[d] &
C'[1] \ar[d] \\
C \ar[r] &
Rb_*\Omega^\bullet_{X'/S} \oplus i_*\Omega^\bullet_{Z/S} \ar[r] &
i_*Rp_*\Omega^\bullet_{E/S} \ar[r] &
C[1]
}
$$
由引理 \ref{derham-lemma-comparison-general} 的 (3) 以及《导出范畴》，命题
\ref{derived-proposition-9}，推出 $C' \to C$ 是同构。由引理
\ref{derham-lemma-comparison-general} 的 (2)，映射
$i_*\Omega^\bullet_{Z/S} \to i_*p_*\Omega^\bullet_{E/S}$ 是同构。
因此在导出范畴中 $C' = b_*\Omega^\bullet_{X'/S}$。最后，引理
\ref{derham-lemma-comparison-general} 的 (1) 告诉我们，它等于
$\Omega^\bullet_{X/S}$。我们略去其与典范映射相容的验证。
\end{proof}

\begin{proposition}
\label{derham-proposition-blowup-split}
采用《态射的进一步结果》，引理 \ref{more-morphisms-lemma-blowup} 中的记号。
映射 $\Omega^\bullet_{X/S} \to Rb_*\Omega^\bullet_{X'/S}$
在 $D(X, (X \to S)^{-1}\mathcal{O}_S)$ 中有分裂。
\end{proposition}

\begin{proof}
考虑引理 \ref{derham-lemma-distinguished-triangle-blowup} 中构造的三角形。
我们声称映射
$$
Rb_*(\Omega^\bullet_{X'/S}) \oplus i_*\Omega^\bullet_{Z/S} \to
i_*Rp_*(\Omega^\bullet_{E/S})
$$
有一个分裂，其像包含直和项 $i_*\Omega^\bullet_{Z/S}$。由《导出范畴》，
引理 \ref{derived-lemma-split}，这表明三角形的第一个箭头有一个分裂，
且该分裂在直和项 $i_*\Omega^\bullet_{Z/S}$ 上为零，从而证明该引理。
我们将把 $Rp_*\Omega^\bullet_{E/S}$ 分解为直和来证明此声称：
第一部分对应于 $\Omega^\bullet_{Z/S}$，第二部分可以通过
$Rb_*\Omega^\bullet_{X'/S}$ 提升。

\medskip\noindent
该声称的证明。可以把 $X$ 分解为相对 $S$ 具有固定维数的开闭子概形，
见《态射》，引理 \ref{morphisms-lemma-smooth-omega-finite-locally-free}。
由于导出范畴 $D(X, f^{-1}\mathcal{O})_S)$ 相应地分解为范畴的积，
可以假设 $X$ 具有固定相对维数 $N$（在 $S$ 上）。可以把
$Z = \coprod Z_m$ 分解为相对维数为 $m \geq 0$（在 $S$ 上）的开闭子概形。
态射 $i_m : Z_m \to X$ 是 $i$ 到 $Z_m$ 的限制，并且是余维为
$N - m$ 的正则浸入，
见《除子》，引理
\ref{divisors-lemma-immersion-smooth-into-smooth-regular-immersion}。
令 $E = \coprod E_m$ 为相应分解，即置 $E_m = p^{-1}(Z_m)$。
若 $p_m : E_m \to Z_m$ 表示 $p$ 到 $E_m$ 的限制，则有典范同构
$$
\tilde \xi_m :
\bigoplus\nolimits_{t = 0, \ldots, N - m - 1}
\Omega^\bullet_{Z_m/S}[-2t]
\longrightarrow
Rp_{m, *}\Omega^\bullet_{E_m/S}
$$
于 $D(Z_m, (Z_m \to S)^{-1}\mathcal{O}_S)$ 中；在次数 $0$ 上，它是典范映射
$\Omega^\bullet_{Z_m/S} \to Rp_{m, *}\Omega^\bullet_{E_m/S}$.
见评注 \ref{derham-remark-projective-space-bundle-formula}。因此有同构
$$
\tilde \xi :
\bigoplus\nolimits_m
\bigoplus\nolimits_{t = 0, \ldots, N - m - 1}
\Omega^\bullet_{Z_m/S}[-2t]
\longrightarrow
Rp_*(\Omega^\bullet_{E/S})
$$
于 $D(Z, (Z \to S)^{-1}\mathcal{O}_S)$ 中；它到源的直和项
$\Omega^\bullet_{Z/S} = \bigoplus \Omega^\bullet_{Z_m/S}$ 的限制
是典范映射 $\Omega^\bullet_{Z/S} \to Rp_*(\Omega^\bullet_{E/S})$。
考虑复形的子复形 $M_m$ 和 $K_m$，这里该复形为
$\bigoplus\nolimits_{t = 0, \ldots, N - m - 1} \Omega^\bullet_{Z_m/S}[-2t]$
，它们在评注 \ref{derham-remark-projective-space-bundle-formula} 中引入。置
$$
M = \bigoplus M_m
\quad\text{且}\quad
K = \bigoplus K_m
$$
。有 $M = K[-2]$，并且由构造，映射
$$
c_{E/Z} \oplus \tilde \xi|_M :
\Omega^\bullet_{Z/S} \oplus M
\longrightarrow
Rp_*(\Omega^\bullet_{E/S})
$$
是同构（见上述评注）。

\medskip\noindent
考虑映射
$$
\delta : \Omega^\bullet_{E/S}[-2] \longrightarrow \Omega^\bullet_{X'/S}
$$
，它是 $D(X', (X' \to S)^{-1}\mathcal{O}_S)$ 中引理
\ref{derham-lemma-gysin-via-log-complex} 的映射，并具有如下性质：复合
$$
\Omega^\bullet_{E/S}[-2] \longrightarrow \Omega^\bullet_{X'/S}
\longrightarrow
\Omega^\bullet_{E/S}
$$
是映射 $\theta'$；它见评注 \ref{derham-remark-cup-product-as-a-map}，并对应于
$c_1^{dR}(\mathcal{O}_{X'}(-E))|_E) = c_1^{dR}(\mathcal{O}_E(1))$。
评注 \ref{derham-remark-projective-space-bundle-formula}
的最后一个断言表明图表
$$
\xymatrix{
K[-2] \ar[d]_{(\tilde \xi|_K)[-2]} \ar[r]_{\text{id}} &
M \ar[d]^{\tilde x|_M} \\
Rp_*\Omega^\bullet_{E/S}[-2] \ar[r]^-{Rp_*\theta'} &
Rp_*\Omega^\bullet_{E/S}
}
$$
交换。因此，可用如下映射得到该声称所需的分裂：
\begin{align*}
Rp_*(\Omega^\bullet_{E/S})
& \xrightarrow{(c_{E/Z} \oplus \tilde \xi|_M)^{-1}}
\Omega^\bullet_{Z/S} \oplus M \\
& \xrightarrow{\text{id} \oplus \text{id}^{-1}}
\Omega^\bullet_{Z/S} \oplus K[-2] \\
& \xrightarrow{\text{id} \oplus (\tilde \xi|_K)[-2]}
\Omega^\bullet_{Z/S} \oplus Rp_*\Omega^\bullet_{E/S}[-2] \\
& \xrightarrow{\text{id} \oplus Rb_*\delta}
\Omega^\bullet_{Z/S} \oplus Rb_*\Omega^\bullet_{X'/S}
\end{align*}
上面所述 $\theta'$ 与 $\delta$ 的关系，再加上上面包含
$\theta'$、$\tilde \xi|_K$ 和 $\tilde \xi|_M$ 的交换图表，恰好说明
这是典范映射
$\Omega^\bullet_{Z/S} \oplus Rb_*(\Omega^\bullet_{X'/S}) \to
Rp_*(\Omega^\bullet_{E/S})$ 的一个截面。该声称的证明完成。
\end{proof}

\noindent
引理 \ref{derham-lemma-splitting-on-omega-a} 表明，在霍奇上同调上构造分裂，
比命题 \ref{derham-proposition-blowup-split} 的结果容易得多。
我们力劝读者跳到下一节。

\begin{lemma}
\label{derham-lemma-ext-zero}
设 $i : Z \to X$ 是余维为 $c$ 的正则概形闭浸入。则
$\Ext^q_{\mathcal{O}_X}(i_*\mathcal{F}, \mathcal{E}) = 0$，只要 $q < c$；
这里 $\mathcal{E}$ 在 $X$ 上局部自由，而 $\mathcal{F}$ 是任意
$\mathcal{O}_Z$-模。
\end{lemma}

\begin{proof}
由 $\Ext$ 的局部到整体谱序列，只需在 $X$ 上仿射局部地证明，
见《上同调》第 \ref{cohomology-section-ext} 节。因此可以假设
$X = \Spec(A)$，且存在正则序列 $f_1, \ldots, f_c$，它位于 $A$ 中，使得
$Z = \Spec(A/(f_1, \ldots, f_c))$。可以假设 $c \geq 1$。于是
$f_1 : \mathcal{E} \to \mathcal{E}$ 是单射。由于 $i_*\mathcal{F}$
被 $f_1$ 零化，这表明引理对 $i = 0$ 成立，并且有满射
$$
\Ext^{q - 1}_{\mathcal{O}_X}(i_*\mathcal{F}, \mathcal{E}/f_1\mathcal{E})
\longrightarrow
\Ext^q_{\mathcal{O}_X}(i_*\mathcal{F}, \mathcal{E})
$$
。所以只需证明该箭头的源为零。接着重复这个论证：若 $c \geq 2$，则映射
$f_2 : \mathcal{E}/f_1\mathcal{E} \to \mathcal{E}/f_1\mathcal{E}$
是单射。由于 $i_*\mathcal{F}$ 被 $f_2$ 零化，这表明引理对 $q = 1$
成立，并且有满射
$$
\Ext^{q - 2}_{\mathcal{O}_X}(i_*\mathcal{F},
\mathcal{E}/f_1\mathcal{E} + f_2\mathcal{E})
\longrightarrow
\Ext^{q - 1}_{\mathcal{O}_X}(i_*\mathcal{F}, \mathcal{E}/f_1\mathcal{E})
$$
如此继续即可证明引理。
\end{proof}

\begin{lemma}
\label{derham-lemma-splitting-on-omega-a}
采用《态射的进一步结果》，引理 \ref{more-morphisms-lemma-blowup} 中的记号。
对 $a \geq 0$，存在唯一箭头
$Rb_*\Omega^a_{X'/S} \to \Omega^a_{X/S}$ 于 $D(\mathcal{O}_X)$ 中；它与
$\Omega^a_{X/S} \to Rb_*\Omega^a_{X'/S}$ 的复合是
$\Omega^a_{X/S}$ 上的恒等态射。
\end{lemma}

\begin{proof}
可以把 $X$ 分解为相对 $S$ 具有固定维数的开闭子概形，见《态射》，引理
\ref{morphisms-lemma-smooth-omega-finite-locally-free}。由于导出范畴
$D(X, f^{-1}\mathcal{O})_S)$ 相应地分解为范畴的积，可以假设
$X$ 具有固定相对维数 $N$（在 $S$ 上）。可以把
$Z = \coprod Z_m$ 分解为相对维数为 $m \geq 0$（在 $S$ 上）的开闭子概形。
态射 $i_m : Z_m \to X$ 是 $i$ 到 $Z_m$ 的限制，并且是余维为
$N - m$ 的正则浸入，见《除子》，引理
\ref{divisors-lemma-immersion-smooth-into-smooth-regular-immersion}。
令 $E = \coprod E_m$ 为相应分解，即置 $E_m = p^{-1}(Z_m)$。
我们声称存在自然映射
$$
b^*\Omega^a_{X/S} \to \Omega^a_{X'/S} \to
b^*\Omega^a_{X/S} \otimes_{\mathcal{O}_{X'}}
\mathcal{O}_{X'}(\sum (N - m - 1)E_m)
$$
，其复合由包含
$\mathcal{O}_{X'} \to \mathcal{O}_{X'}(\sum (N - m - 1)E_m)$.
诱导。为证明这一点，只需说明第一个箭头的余核在 $X'$ 上局部地被有效
Cartier 除子 $\sum (N - m - 1)E_m$ 的一个局部方程零化。为此，
可以像引理 \ref{derham-lemma-comparison-general} 的证明那样在 $X$ 上平展局部地
工作，并应用引理 \ref{derham-lemma-comparison-bis}。利用引理
\ref{derham-lemma-blowup-twist-same-cohomology} 在平展局部计算，可见所诱导的复合
$$
\Omega^a_{X/S} \to Rb_*\Omega^a_{X'/S} \to
Rb_*\left(b^*\Omega^a_{X/S} \otimes_{\mathcal{O}_{X'}}
\mathcal{O}_{X'}(\sum (N - m - 1)E_m)\right)
$$
在 $D(\mathcal{O}_X)$ 中是同构；由此得到引理中映射的存在性。

\medskip\noindent
对于唯一性，只需证明不存在从 $\tau_{\geq 1}Rb_*\Omega_{X'/S}$
到 $\Omega^a_{X/S}$ 的非零映射（在 $D(\mathcal{O}_X)$ 中）。
为此，只需再证明不存在从 $R^qb_*\Omega_{X'/s}[-q]$
到 $\Omega^a_{X/S}$ 的非零映射（在 $D(\mathcal{O}_X)$ 中），其中
$q \geq 1$（细节从略）。由引理 \ref{derham-lemma-comparison-general} 可见，
$R^qb_*\Omega_{X'/s} \cong i_*R^qp_*\Omega^a_{E/S}$ 是定义在
$Z = \coprod Z_m$ 上的一个模的推前。此外，注意
$R^qp_*\Omega^a_{E/S}$ 到 $Z_m$ 的限制只有在 $q < N - m$ 时才可能非零。
事实上，$E_m \to Z_m$ 的纤维维数为 $N - m - 1$，可以应用《极限》，引理
\ref{limits-lemma-higher-direct-images-zero-above-dimension-fibre}。
因此所需消失由引理 \ref{derham-lemma-ext-zero} 得到。
\end{proof}














\section{比较微分形式层}
\label{derham-section-quasi-finite-syntomic}

\noindent
本节的目标是：对每个局部拟有限的 syntomic 概形态射 $f : Y \to X$，
构造映射
$$
c^p_{Y/X} :
\Omega^p_{Y/\mathbf{Z}}
\longrightarrow
f^*\Omega^p_{X/\mathbf{Z}} \otimes_{\mathcal{O}_Y} \det(\NL_{Y/X})
$$
，对所有 $p \geq 0$ 均满足下列性质：
\begin{enumerate}
\item
\label{derham-item-degree-zero}
$c_{Y/X}^0(1) = \delta(\NL_{Y/X})$，见《判别式》第
\ref{discriminant-section-tate-map} 节；
\item
\label{derham-item-multiplicative}
$c_{Y/X}^{q + p}(\omega \wedge \eta) = \omega \wedge c_{Y/X}^p(\eta)$
，其中 $\omega$ 是 $f^*\Omega^q_{X/\mathbf{Z}}$ 的局部截面，
而 $\eta$ 是 $\Omega^p_{Y/\mathbf{Z}}$ 的局部截面；
\item
\label{derham-item-base-change}
$c^p_{Y/X}$ 的构造与到开集的限制以及基变换相容
（表述见引理 \ref{derham-lemma-base-change-Garel-upstairs}）。
\end{enumerate}
注意，这些条件蕴含映射 $c^p_{Y/X}$ 是 $\mathcal{O}_Y$-线性的，
并且它与
$f^*\Omega^p_{X/\mathbf{Z}} \to \Omega^p_{Y/\mathbf{Z}}$
的复合是乘以 $\delta(\NL_{Y/X})$。事实上，我们将证明存在唯一的一族算子
$c^p_{Y/X}$ 具有性质 (1)、(2)、(3)。若 $f$ 是有限的，第
\ref{derham-section-trace} 节会把这个结果用于构造德拉姆复形上的迹映射。

\begin{lemma}
\label{derham-lemma-funny-map}
设 $R$ 是环，并考虑交换图表
$$
\xymatrix{
0 \ar[r] &
K^0 \ar[r] &
L^0 \ar[r] &
M^0 \ar[r] & 0 \\
& & L^{-1} \ar[u]_\partial \ar@{=}[r] &
M^{-1} \ar[u]
}
$$
，其中各项为 $R$-模，顶行正合，而 $M^0$ 和 $M^{-1}$
是秩相同的有限自由模。对 $p \geq 0$，存在典范映射
$$
c^p : 
\wedge^p(H^0(L^\bullet))
\longrightarrow
\wedge^p(K^0) \otimes_R \det(M^\bullet)
$$
，具有下列性质：
\begin{enumerate}
\item $c^0(1) = \delta(M^\bullet)$；
\item $c^{q + p}(\omega \wedge \eta) = \omega \wedge c^p(\eta)$
，其中 $\omega \in \wedge^q(K^0)$ 且
$\eta \in \wedge^p(H^0(L^\bullet))$。
\end{enumerate}
特别地，$c^p$ 与映射
$\wedge^p(K^0) \to \wedge^p(H^0(L^\bullet))$ 的复合是乘以
$\delta(M^\bullet)$。
\end{lemma}

\begin{proof}
设 $M^0$ 和 $M^{-1}$ 都是秩为 $n$ 的自由模。对每个 $p \geq 0$，
存在唯一的满射
$$
\pi_p :
\wedge^{p + n}(L^0)
\longrightarrow
\wedge^p(K^0) \otimes \wedge^n(M^0)
$$
，具有下列两个性质；(i) 有
$
\pi_p(k_1 \wedge \ldots \wedge k_p \wedge l_1 \wedge \ldots \wedge l_n) =
k_1 \wedge \ldots \wedge k_p \otimes
m_1 \wedge \ldots \wedge m_n
$
，如果 $k_1, \ldots, k_p \in K^0$，并且
$l_1, \ldots, l_n \in L^0$ 映到 $m_1, \ldots, m_n \in M^0$；
(ii) $\pi_p$ 的核是由外积 $l_1 \wedge \ldots \wedge l_{p + n}$
生成的子模，其中 $> p$ 个 $l_j$ 属于 $K^0$。取一组基
$e_1, \ldots, e_n$，它是 $L^{-1} = M^{-1}$ 的基。然后按如下规则定义映射：
$$
\eta \mapsto
\pi_p(\tilde \eta \wedge \partial(e_1) \wedge \ldots \wedge \partial(e_n))
\otimes (e_1 \wedge \ldots \wedge e_n)^{\otimes -1}
$$
，其中 $\eta \in \wedge^p(H^0(L^\bullet))$，而
$\tilde \eta \in \wedge^p(L^0)$ 是 $\eta$ 的提升。右边的表达式有意义，因为
$\det(M^\bullet) = \wedge^n(M^0) \otimes \wedge^n(M^{-1})^{\otimes -1}$.
右边的表达式与 $\tilde \eta$ 的选择无关，因为
$\partial(\xi) \wedge \partial(e_1) \wedge \ldots \wedge \partial(e_n) = 0$
，对任意 $\xi \in L^{-1}$。直接计算表明，该映射与
$L^{-1} = M^{-1}$ 的基的选择无关。由《代数的进一步结果》第
\ref{more-algebra-section-determinants-complexes} 节中
$\delta(M^\bullet)$ 的定义，立即得到
$c^0(1) = \delta(M^\bullet)$。最后，陈述 (2) 中的规则可由定义直接验证。
\end{proof}

\begin{lemma}
\label{derham-lemma-funny-map-functorial}
设 $R_1 \to R_2$ 是环同态。对 $i = 1, 2$，考虑交换图表
$$
\xymatrix{
0 \ar[r] &
K_i^0 \ar[r] &
L_i^0 \ar[r] &
M_i^0 \ar[r] & 0 \\
& & L_i^{-1} \ar[u]_{\partial_i} \ar@{=}[r] &
M_i^{-1} \ar[u]
}
$$
，其中各项为 $R_i$-模，如引理 \ref{derham-lemma-funny-map} 所述。假设有映射
$K_1^0 \to K_2^0$、$L_1^j \to L_2^j$、$M_1^j \to M_2^j$，
它们与给定环同态 $R_1 \to R_2$ 以及显示图表中的映射相容。若映射
$M_1^j \otimes_{R_1} R_2 \to M_2^j$ 都是同构，则图表
$$
\xymatrix{ 
\wedge^p(H^0(L_1^\bullet)) \ar[d]
\ar[r]_-{c^p} &
\wedge^p(K_1^0) \otimes_{R_1} \det(M_1^\bullet) \ar[d] \\
\wedge^p(H^0(L_2^\bullet))
\ar[r]^-{c^p} &
\wedge^p(K_2^0) \otimes_{R_2} \det(M_2^\bullet)
}
$$
交换。
\end{lemma}

\begin{proof}
这由引理 \ref{derham-lemma-funny-map} 的证明中给出的映射的显式描述得到。
注意，需要假设 $M_1^j \otimes_{R_1} R_2 \to M_2^j$ 都是同构，
才能看出我们用于 $M_1^{-1}$ 的基映到 $M_2^{-1}$ 的一组基
（因而在检验两个构造的相容性时可以使用“同一”组基）。
\end{proof}

\begin{remark}
\label{derham-remark-local-description}
设 $A$ 是一个环，并令 $P = A[x_1, \ldots, x_n]$。取
$f_1, \ldots, f_n \in P$，并置 $B = P/(f_1, \ldots, f_n)$。
假设 $A \to B$ 是拟有限的。于是
$B$ 是 $A$ 上的相对全局完全交（代数，定义
\ref{algebra-definition-relative-global-complete-intersection}），且由代数，引理
\ref{algebra-lemma-relative-global-complete-intersection-conormal}，
$(f_1, \ldots, f_n)/(f_1, \ldots, f_n)^2$ 是以各类
$\overline{f}_i$ 为生成元的自由模。考虑下图
$$
\xymatrix{
\Omega_{A/\mathbf{Z}} \otimes_A B \ar[r] &
\Omega_{P/\mathbf{Z}} \otimes_P B \ar[r] &
\Omega_{P/A} \otimes_P B \\
&
(f_1, \ldots, f_n)/(f_1, \ldots, f_n)^2 \ar[u] \ar@{=}[r] &
(f_1, \ldots, f_n)/(f_1, \ldots, f_n)^2 \ar[u]
}
$$
右列表示 $\NL_{B/A}$ 在 $D(B)$ 中的一个代表，因而有上同调
$\Omega_{B/A}$ 位于次数 $0$。顶行是分裂短正合列
$0 \to \Omega_{A/\mathbf{Z}} \otimes_A B \to
\Omega_{P/\mathbf{Z}} \otimes_P B \to \Omega_{P/A} \otimes_P B \to 0$.
由代数，引理 \ref{algebra-lemma-differential-seq}，中列有上同调
$\Omega_{B/\mathbf{Z}}$ 位于次数 $0$。
因此由引理 \ref{derham-lemma-funny-map}，我们得到典范的 $B$-模映射
$$
\Omega^p_{B/\mathbf{Z}} \longrightarrow
\Omega^p_{A/\mathbf{Z}} \otimes_A \det(\NL_{B/A})
$$
其中 $p \geq 0$，这些映射满足性质 (1) 和 (2)；特别地，它们与映射
$\Omega^p_{A/\mathbf{Z}} \to \Omega^p_{B/\mathbf{Z}}$
的复合是乘以 $\delta(\NL_{B/A})$。
\end{remark}

\begin{lemma}
\label{derham-lemma-base-change-Garel-upstairs}
考虑概形的交换图
$$
\xymatrix{
Y' \ar[d]_{f'} \ar[r]_b & Y \ar[d]^f \\
X' \ar[r]^a & X
}
$$
它将 $Y'$ 同构到 $X' \times_X Y$ 的一个开子概形。假设 $f$ 是
局部拟有限的 syntomic 态射，并且给定映射
$c^p_{Y/X} : \Omega^p_{Y/\mathbf{Z}} \to
f^*\Omega^p_{X/\mathbf{Z}} \otimes_{\mathcal{O}_Y} \det(\NL_{Y/X})$
（$p \geq 0$），满足
(\ref{derham-item-degree-zero}) 和 (\ref{derham-item-multiplicative})。
那么至多存在一族映射
$c^p_{Y'/X'} : \Omega^p_{Y'/\mathbf{Z}} \to
(f')^*\Omega^p_{X'/\mathbf{Z}} \otimes_{\mathcal{O}_{Y'}} \det(\NL_{Y'/X'})$
（$p \geq 0$），满足 (\ref{derham-item-degree-zero}) 和
(\ref{derham-item-multiplicative})，并使诸图
$$
\xymatrix{
b^*\Omega^p_{Y/\mathbf{Z}} \ar[rr]_-{b^*c^p_{Y/X}} \ar[d] & &
b^*(f^*\Omega^p_{X/\mathbf{Z}} \otimes_{\mathcal{O}_Y} \det(\NL_{Y/X}))
\ar@{=}[r] &
(f')^*a^*\Omega^p_{X/\mathbf{Z}} \otimes_{\mathcal{O}_Y} b^*\det(\NL_{Y/X}))
\ar[d] \\
\Omega^p_{Y'/\mathbf{Z}} \ar[rrr]^-{c^p_{Y'/X'}} & & &
(f')^*\Omega^p_{X'/\mathbf{Z}} \otimes_{\mathcal{O}_{Y'}} \det(\NL_{Y'/X'})
}
$$
对所有 $p \geq 0$ 交换。这里的竖直箭头使用第
\ref{derham-section-de-rham-complex} 节中的映射
$b^*\Omega^p_{Y/\mathbf{Z}} \to \Omega^p_{Y'/\mathbf{Z}}$ 和
$a^*\Omega^p_{X/\mathbf{Z}} \to \Omega^p_{X'/\mathbf{Z}}$
以及判别式，第 \ref{discriminant-section-tate-map} 节中的等同
$b^*\det(\NL_{Y/X}) = \det(\NL_{Y'/X'})$。
\end{lemma}

\begin{proof}
映射
$$
(f')^*\Omega_{X'/\mathbf{Z}} \oplus b^*\Omega_{Y/\mathbf{Z}}
\longrightarrow
\Omega_{Y'/\mathbf{Z}}
$$
是满射，因为 $Y'$ 是该纤维积的一个开子概形；相应的代数陈述是：
$\Omega_{B \otimes_A C/\mathbf{Z}}$ 作为模由
$\Omega_{B/\mathbf{Z}}$ 和 $\Omega_{C/\mathbf{Z}}$ 的像生成。
因此对所有 $p$，映射
$$
\bigoplus\nolimits_{i = 0, \ldots, p}
(f')^*\Omega^i_{X'/\mathbf{Z}} \otimes
b^*\Omega^{p - i}_{Y/\mathbf{Z}}
\longrightarrow
\Omega^p_{Y'/\mathbf{Z}}
$$
是满射。条件 (\ref{derham-item-degree-zero})、(\ref{derham-item-multiplicative})
连同引理中图的交换性表明，映射 $c^p_{Y'/X'}$
已被唯一确定（但其存在性并非立即可得）。
\end{proof}

\begin{lemma}
\label{derham-lemma-existence-base-change-Garel-upstairs}
考虑交换图
$$
\xymatrix{
Y' \ar[d]_{f'} \ar[r]_b & Y \ar[d]^f \\
X' \ar[r]^a & X
}
$$
它将 $Y'$ 同构到 $X' \times_X Y$ 的一个开子概形。假设 $f$ 是局部拟有限的
syntomic 态射。那么对每个 $y' \in Y'$，都可以找到开集
$V' \subset Y'$、$V \subset Y$、$U' \subset X$、$U \subset X$，
使得 $y' \in V'$，并且
$f'(V') \subset U'$, $b(V') \subset V$, $a(U') \subset U$,
$f(V) \subset U$；而且对 $p \geq 0$ 存在映射
$c^p_{V/U}$ 和 $c^p_{V'/U'}$，它们满足 (\ref{derham-item-degree-zero})
和 (\ref{derham-item-multiplicative})，并与图
$$
\xymatrix{
V' \ar[d] \ar[r]_b & V \ar[d] \\
U' \ar[r] & U
}
$$
相容，其含义如引理 \ref{derham-lemma-base-change-Garel-upstairs} 中所述。
\end{lemma}

\begin{proof}
显然，为了证明这一点，可以把 $Y'$ 替换为 $X' \times_X Y$。
按照判别式，引理 \ref{discriminant-lemma-syntomic-quasi-finite} 的 (5)，
取仿射开集 $V \subset Y$ 和 $U \subset X$，其中 $V$ 包含 $y'$ 的像。
取仿射开集 $U' \subset X'$，它包含 $y'$ 的像并映入 $U$。
把 $X, Y, X', Y'$ 分别替换为 $U, V, U', U' \times_U V$ 后，
我们可以假设处于下一段所述的情形。

\medskip\noindent
假设 $X' = \Spec(A')$、$X = \Spec(A)$、$Y = \Spec(B)$，且
$Y' = \Spec(A' \otimes_A B)$，其中
$B = A[x_1, \ldots, x_n]/(f_1, \ldots, f_n)$
是 $A$ 上的相对全局完全交。那么
$B' = A'[x_1, \ldots, x_n]/(f'_1, \ldots, f'_n)$
也是 $A'$ 上的相对全局完全交，其中 $f'_i$ 是 $f_i$ 在
$A'[x_1, \ldots, x_n]$ 中的像；见代数，引理
\ref{algebra-lemma-base-change-relative-global-complete-intersection}。
备注 \ref{derham-remark-local-description} 中的构造给出了映射
$c^p_{V/U}$ 和 $c^p_{V'/U'}$。
由引理 \ref{derham-lemma-funny-map-functorial} 以及 $B$、$B'$ 分别在
$A$、$A'$ 上的表示之间的相容性，这些映射彼此相容。略去若干细节。
\end{proof}

\begin{lemma}
\label{derham-lemma-Garel-upstairs}
存在唯一的规则，它为每个概形的局部拟有限 syntomic 态射
$f : Y \to X$ 指定 $\mathcal{O}_Y$-模映射
$$
c^p_{Y/X} :
\Omega^p_{Y/\mathbf{Z}}
\longrightarrow
f^*\Omega^p_{X/\mathbf{Z}} \otimes_{\mathcal{O}_Y} \det(\NL_{Y/X})
$$
其中 $p \geq 0$，并满足 (\ref{derham-item-degree-zero})、
(\ref{derham-item-multiplicative}) 和 (\ref{derham-item-base-change})。特别地，
$c^p_{Y/X}$ 与 $f^*\Omega^p_{X/\mathbf{Z}} \to \Omega^p_{Y/\mathbf{Z}}$
的复合是乘以 $\delta(\NL_{Y/X})$。
\end{lemma}

\begin{proof}
此证明与判别式，命题 \ref{discriminant-proposition-tate-map}
的证明非常相似，建议读者先阅读该证明。

\medskip\noindent
我们先重新表述该陈述。考虑范畴 $\mathcal{C}$：其对象记为
$Y/X$，是概形的局部拟有限 syntomic 态射 $f : Y \to X$；
其态射 $b/a : Y'/X' \to Y/X$ 是交换图
$$
\xymatrix{
Y' \ar[d]_{f'} \ar[r]_b & Y \ar[d]^f \\
X' \ar[r]^a & X
}
$$
它将 $Y'$ 同构到 $X' \times_X Y$ 的一个开子概形。引理的含义是：
对每个对象 $Y/X$（属于 $\mathcal{C}$），都有映射
$c^p_{Y/X}$（$p \geq 0$）满足性质 (\ref{derham-item-degree-zero})
和 (\ref{derham-item-multiplicative})；并且对每个态射
$b/a : Y'/X' \to Y/X$（属于 $\mathcal{C}$），映射
$c^p_{Y/X}$ 和 $c^p_{Y'/X'}$
按引理 \ref{derham-lemma-base-change-Garel-upstairs} 的意义相容。

\medskip\noindent
给定 $Y/X$（属于 $\mathcal{C}$）以及 $y \in Y$，可以找到
仿射开集 $V \subset Y$ 和 $U \subset X$，满足 $f(V) \subset U$，
使得存在一些映射
$$
\Omega^p_{Y/\mathbf{Z}}|_V
\longrightarrow
\left(
f^*\Omega^p_{X/\mathbf{Z}} \otimes_{\mathcal{O}_Y} \det(\NL_{Y/X})
\right)|_V
$$
满足性质 (\ref{derham-item-degree-zero}) 和 (\ref{derham-item-multiplicative})。
这是因为可以按照判别式，引理
\ref{discriminant-lemma-syntomic-quasi-finite} 的 (5) 选取仿射开集，
再使用备注 \ref{derham-remark-local-description} 中的构造。

\medskip\noindent
注意，$f$ 的平展轨迹恰好是 $\delta(\NL_{Y/X})$ 不消失的点集；
见判别式，第 \ref{discriminant-section-tate-map} 节中的讨论。特别地，
$f^*\Omega^p_{X/\mathbf{Z}} \to \Omega^p_{Y/\mathbf{Z}}$
在逆转 $\delta(\NL_{Y/X})$ 后成为同构。由此可知，如果
$\Omega^p_{X/\mathbf{Z}}$ 是有限局部自由的，且
$\delta(\NL_{Y/X})$ 在 $\mathcal{O}_Y$ 中的零化子为零，
那么上一段构造的局部映射是唯一的，并会自动粘合！

\medskip\noindent
令 $\mathcal{C}_{nice} \subset \mathcal{C}$ 表示由满足下列条件的
$Y/X$ 构成的全子范畴：
\begin{enumerate}
\item $\Omega_{X/\mathbf{Z}}$ 是局部自由的；
\item $\delta(\NL_{Y/X})$ 在 $\mathcal{O}_Y$ 中的零化子为零。
\end{enumerate}
由上一段的说明，对每个对象 $Y/X$（属于 $\mathcal{C}_{nice}$），
都有唯一的映射 $c^p_{Y/X}$（$p \geq 0$）满足条件
(\ref{derham-item-degree-zero}) 和 (\ref{derham-item-multiplicative})。
如果 $b/a : Y'/X' \to Y/X$ 是 $\mathcal{C}_{nice}$ 的一个态射，
那么映射 $c^p_{Y/X}$ 和 $c^p_{Y'/X'}$ 按照引理
\ref{derham-lemma-base-change-Garel-upstairs} 的意义相容：事实上，由引理
\ref{derham-lemma-existence-base-change-Garel-upstairs}，局部上确实存在相容映射；
而刚才的唯一性说明这些映射分别等于
$c^p_{Y/X}$ 和 $c^p_{Y'/X'}$。
换言之，我们已经在全子范畴 $\mathcal{C}_{nice}$ 上解决了该问题。
对 $Y/X$（属于 $\mathcal{C}_{nice}$），仍用 $c^p_{Y/X}$
表示刚刚得到的解。

\medskip\noindent
事实上，更一般地，假设有态射 $b/a : Y'/X' \to Y/X$
（属于 $\mathcal{C}$），并且 $Y/X$ 是 $\mathcal{C}_{nice}$ 的对象。
同一论证（这次使用引理 \ref{derham-lemma-base-change-Garel-upstairs}
中的唯一性）告诉我们，存在映射 $c^p_{Y'/X'}$（$p \geq 0$），
它们满足条件 (\ref{derham-item-degree-zero}) 和 (\ref{derham-item-multiplicative})，
并与已经构造的映射 $c^p_{Y/X}$ 相容。我们把它称为基变换
$(b/a)^*c^p_{Y/X}$。

\medskip\noindent
考虑态射
$$
Y_1/X_1 \xleftarrow{b_1/a_1} Y/X \xrightarrow{b_2/a_2} Y_2/X_2
$$
（属于 $\mathcal{C}$），并假设 $Y_1/X_1$ 和 $Y_2/X_2$
是 $\mathcal{C}_{nice}$ 的对象。
{\bf 断言。} 两个基变换 $(b_i/a_i)^*c^p_{Y_i/X_i}$（$i = 1, 2$）
相等。我们先说明该断言蕴含引理，然后再证明该断言。

\medskip\noindent
令 $d, n \geq 1$，并考虑判别式，例
\ref{discriminant-example-universal-quasi-finite-syntomic}
中构造的局部拟有限 syntomic 态射 $Y_{n, d} \to X_{n, d}$。
那么 $Y_{n, d}$ 和 $Y_{n, d}$ 是 $\mathbf{Z}$ 上光滑的
有限型不可约概形。事实上，$X_{n, d}$ 是 $\mathbf{Z}$ 上某个多项式环的谱，
而 $Y_{n, d}$ 是这种谱的一个开子概形。态射
$Y_{n, d} \to X_{n, d}$ 是局部拟有限的 syntomic 态射，并在某个稠密开集上平展；
见判别式，引理
\ref{discriminant-lemma-universal-quasi-finite-syntomic-etale}。
因此 $\delta(\NL_{Y_{n, d}/X_{n, d}})$ 非零：例如，判别式，备注
\ref{discriminant-remark-local-description-delta} 给出了
$\delta(\NL_{Y/X})$ 的局部描述，而态射，引理
\ref{morphisms-lemma-etale-at-point} 的 (8) 给出了平展态射的局部描述。
不可约正则概形上的可逆模的非零截面具有零零化子。因此
$Y_{n, d}/X_{n, d}$ 是 $\mathcal{C}_{nice}$ 的对象。

\medskip\noindent
令 $Y/X$ 是 $\mathcal{C}$ 的任意对象，并令 $y \in Y$。
由判别式，引理 \ref{discriminant-lemma-locally-comes-from-universal}，
可以找到 $n, d \geq 1$ 以及态射
$$
Y/X \leftarrow V/U \xrightarrow{b/a} Y_{n, d}/X_{n, d}
$$
（属于 $\mathcal{C}$），使得 $V \subset Y$ 和 $U \subset X$ 都是开集。
于是有基变换 $c^p_{V/U} = (b/a)^*c^p_{Y_{n, d}/X_{n, d}}$。
该断言保证这些局部构造的映射 $c^p_{V/U}$ 能够粘合！
因此得到满足性质 (\ref{derham-item-degree-zero}) 和
(\ref{derham-item-multiplicative}) 的良定义全局映射 $c^p_{Y/X}$。
最后，令 $b/a : Y'/X' \to Y/X$ 是 $\mathcal{C}$ 的任意态射。
本段已经构造了映射 $c^p_{Y'/X'}$（$p \geq 0$）和
$c^p_{Y/X}$（$p \geq 0$）。为了检验它们按引理
\ref{derham-lemma-base-change-Garel-upstairs} 的意义相容，可以在
$Y'/X'$ 和 $Y/X$ 上局部地工作。因此可以假设存在态射
$Y/X \to Y_{n, d}/X_{n, d}$。依构造，族
$c^p_{Y'/X'}$ 和族 $c^p_{X/Y}$ 都与指向
$Y_{n, d}/X_{n, d}$ 的映射相容。由此（以及引理
\ref{derham-lemma-base-change-Garel-upstairs} 中的唯一性）容易推出
$c^p_{Y'/X'}$ 与 $c^p_{Y/X}$ 相容。因此只需证明该断言。

\medskip\noindent
在证明的余下部分中，我们证明该断言。可以取一点 $y \in Y$，
并证明这些映射在 $y$ 的某个开邻域中相等。因此，可以用开邻域替换
$Y_1$、$Y_2$；这些开邻域分别包含 $y$ 在 $Y_1$ 和 $Y_2$ 中的像。
从而可以假设 $Y, X, Y_1, X_1, Y_2, X_2$ 都是仿射的，例如分别是环
$B, A, B_1, A_1, B_2, A_2$ 的谱。图示为
$$
\xymatrix{
B_1 \ar[r] & B & B_2 \ar[l] \\
A_1 \ar[u] \ar[r] & A \ar[u] & A_2 \ar[l] \ar[u]
}
$$
依假设，$B$ 的谱同时是 $A \otimes_{A_1} B_1$ 的谱和
$A \otimes_{A_2} B_2$ 的谱的仿射开集。进一步缩小后，可以假设存在元素
$g_i \in A \otimes_{A_i} B_i$，使我们的映射给出同构
$(A \otimes_{A_i} B_i)_{g_i} = B$；见
性质，引理 \ref{properties-lemma-standard-open-two-affines}。
令 $x_\alpha$、$y_\beta$ 为足够大的一组变量，使得可以选取满射
$$
A'_1 = A_1[x_\alpha] \to A
\quad\text{且}\quad
A'_2 = A_2[y_\beta] \to A
$$
，它们分别是 $A_1$-代数和 $A_2$-代数的满射。随后可以选取提升
$h_i \in A'_i \otimes_{A_i} B_i$，它提升
$g_i \in A \otimes_{A_i} B_i$；并考虑图
$$
\xymatrix{
(A'_1 \otimes_{A_1} B_1)_{h_1} \ar[r] & B &
(A'_2 \otimes_{A_1} B_2)_{h_2} \ar[l] \\
A'_1 \ar[u] \ar[r] & A \ar[u] & A'_2 \ar[l] \ar[u]
}
$$
依构造，两个方块都是余笛卡尔的。接着考虑环映射
$$
A' = A'_1 \times_A A'_2 \longrightarrow
B' =
(A'_1 \otimes_{A_1} B_1)_{h_1} \times_B
(A'_2 \otimes_{A_1} B_2)_{h_2}
$$
由更多代数，引理
\ref{more-algebra-lemma-module-over-fibre-product}，有
$A'_1 \otimes_{A'} B' = (A'_1 \otimes_{A_1} B_1)_{h_1}$ 和
$A'_2 \otimes_{A'} B' = (A'_2 \otimes_{A_2} B_2)_{h_2}$。特别地，态射
$Y' = \Spec(B') \to \Spec(A') = X'$ 的纤维是态射
$Y_i \to X_i$ 的纤维经基变换所得概形的开子概形，因而是有限完全局部交；
见判别式，引理 \ref{discriminant-lemma-syntomic-quasi-finite}。
由更多代数，引理
\ref{more-algebra-lemma-properties-algebras-over-fibre-product}，
环映射 $A' \to B'$ 是平坦且有限表示的。因此由判别式，引理
\ref{discriminant-lemma-syntomic-quasi-finite} 的 (3)，
可知 $Y'/X'$ 是 $\mathcal{C}$ 的对象。现在考虑交换图
$$
\xymatrix{
& & Y/X \ar[ld] \ar@/_2pc/[lldd]_{b_1/a_1} \ar[rd] \ar@/^2pc/[rrdd]^{b_2/a_2} \\
& Y'_1/X'_1 \ar[rd] \ar[ld] & & Y'_2/X'_2 \ar[ld] \ar[rd] \\
Y_1/X_1 & & Y'/X' & & Y_2/X_2
}
$$
其中 $Y'_i/X'_i$ 对应于
$A'_i \to (A'_i \otimes_{A_i} B_i)_{h_i}$。注意，
$Y'_i/X'_i$ 是 $\mathcal{C}_{nice}$ 的对象，因为它由一个无穷维仿射空间与
$Y_i/X_i$ 的积取开子概形而得；略去一个小细节。特别地，
$c^p_{Y_i/X_i}$ 经 $(b_i/a_i)$ 拉回后，与
$c^p_{Y'_i/X'_i}$ 经 $Y/X \to Y'_i/X'_i$ 拉回后所得的映射相同。
如果 $Y'/X'$ 是 $\mathcal{C}_{nice}$ 的对象，我们现在就已经完成了，
但实际几乎从不如此。因为在这种情况下，沿方块的两边拉回
$c^p_{Y'/X'}$ 就会给出所需结论。为了绕开
$Y'/X'$ 不属于 $\mathcal{C}_{nice}$ 的问题，注意上述论证表明：
在可能缩小所有概形 $X, Y, X'_1, Y'_1, X'_2, Y'_2, X', Y'$ 后，
可以找到某些 $n, d \geq 1$，并把图扩张如下：
$$
\xymatrix{
& Y/X \ar[ld] \ar[rd] \\
Y'_1/X'_1 \ar[rd] & & Y'_2/X'_2 \ar[ld] \\
& Y'/X' \ar[d] \\
& Y_{n, d}/X_{n, d}
}
$$
随后可以通过从 $c^p_{Y_{n, d}/X_{n, d}}$ 拉回来使用已经给出的论证。
证明完毕。
\end{proof}








\section{德拉姆复形上的迹映射}
\label{derham-section-trace}

\noindent
本节部分内容可参见 \cite{Garel}。设 $S$ 是一个概形，并设
$f : Y \to X$ 是 $S$ 上概形之间的有限局部自由态射。于是存在迹映射
$\text{Trace}_f : f_*\mathcal{O}_Y \to \mathcal{O}_X$；见判别式，第
\ref{discriminant-section-discriminant} 节。在这种情形下，
德拉姆复形上的迹映射是复形映射
$$
\Theta_{Y/X} : f_*\Omega^\bullet_{Y/S} \longrightarrow \Omega^\bullet_{X/S}
$$
使得 $\Theta_{Y/X}$ 等于 $\text{Trace}_f$（在次数 $0$ 上），并且满足
$$
\Theta_{Y/X}(\omega \wedge \eta) = \omega \wedge \Theta_{Y/X}(\eta)
$$
其中 $\omega$ 是 $\Omega^\bullet_{X/S}$ 的局部截面，$\eta$ 是
$f_*\Omega^\bullet_{Y/S}$ 的局部截面。我们尚不清楚这样的迹映射
$\Theta_{Y/X}$ 是否对每个有限局部自由态射 $Y \to X$ 都存在；如果你有反例或证明，
请发邮件至
\href{mailto:stacks.project@gmail.com}{stacks.project@gmail.com}
。

\begin{example}
\label{derham-example-no-trace}
下面给出一个德拉姆复形上不存在迹映射的例子。考虑
$\mathbf{C}$-代数 $B = \mathbf{C}[x, y]$，群 $G = \{\pm 1\}$
通过 $x \mapsto -x$ 和 $y \mapsto -y$ 作用于它。不变量
$A = B^G$ 构成 $\mathbf{C}$ 上的有限型正规整环，并由
$x^2, xy, y^2$ 生成。我们断言，对于包含 $A \subset B$，不存在合理的
迹映射
$\Omega_{B/\mathbf{C}} \to \Omega_{A/\mathbf{C}}$
作用在 $1$-形式上。事实上，考虑元素
$\omega = x \text{d} y \in \Omega_{B/\mathbf{C}}$.
由于 $\omega$ 在 $G$ 的作用下不变，如果存在“合理的”迹映射，
那么 $2\omega$ 应当属于
$\Omega_{A/\mathbf{C}} \to \Omega_{B/\mathbf{C}}$ 的像。但事实并非如此：
无法把 $2\omega$ 写成
$\text{d}(x^2)$、$\text{d}(xy)$ 和 $\text{d}(y^2)$ 的线性组合，
即使允许系数取自 $B$。
这个例子与 \cite{Zannier} 中的主定理相矛盾。
\end{example}

\begin{lemma}
\label{derham-lemma-Garel}
存在唯一的规则，它为概形的每个有限 syntomic 态射
$f : Y \to X$ 指定 $\mathcal{O}_X$-模映射
$$
\Theta^p_{Y/X} :
f_*\Omega^p_{Y/\mathbf{Z}}
\longrightarrow
\Omega^p_{X/\mathbf{Z}}
$$
并满足下列性质：
\begin{enumerate}
\item 与映射
$\Omega^p_{X/\mathbf{Z}} \otimes_{\mathcal{O}_X} f_*\mathcal{O}_Y
\to f_*\Omega^p_{Y/\mathbf{Z}}$ 的复合等于
$\text{id} \otimes \text{Trace}_f$
，其中 $\text{Trace}_f : f_*\mathcal{O}_Y \to \mathcal{O}_X$
是判别式，第 \ref{discriminant-section-discriminant} 节中的映射；
\item 该规则与基变换相容。
\end{enumerate}
\end{lemma}

\begin{proof}
首先，假设 $X$ 是局部诺特的。由引理
\ref{derham-lemma-Garel-upstairs}，有典范映射
$$
c^p_{Y/X} : \Omega_{Y/\mathbf{Z}}^p
\longrightarrow
f^*\Omega_{X/\mathbf{Z}}^p \otimes_{\mathcal{O}_Y} \det(\NL_{Y/X})
$$
由判别式，命题 \ref{discriminant-proposition-tate-map}，有典范同构
$$
c_{Y/X} : \det(\NL_{Y/X}) \to \omega_{Y/X}
$$
它把 $\delta(\NL_{Y/X})$ 映到 $\tau_{Y/X}$。合并这些映射得到
$$
c^p_{Y/X} \otimes c_{Y/X} :
\Omega_{Y/\mathbf{Z}}^p
\longrightarrow
f^*\Omega_{X/\mathbf{Z}}^p \otimes_{\mathcal{O}_Y} \omega_{Y/X}
$$
由判别式，第 \ref{discriminant-section-finite-morphisms} 节，
这等价于给出映射
$$
\Theta_{Y/X}^p :
f_*\Omega_{Y/\mathbf{Z}}^p
\longrightarrow
\Omega_{X/\mathbf{Z}}^p
$$
回忆，$c^p_{Y/X} \otimes c_{Y/X}$ 与 $\Theta_{Y/X}^p$
之间的关系使用求值映射
$f_*\omega_{Y/X} \to \mathcal{O}_X$
，它把 $\tau_{Y/X}$ 映到 $\text{Trace}_f(1)$；见判别式，第
\ref{discriminant-section-finite-morphisms} 节。因此性质 (1) 成立。
性质 (2) 对沿 $X' \to X$ 的基变换成立（其中 $X'$ 局部诺特），
因为 $c^p_{Y/X}$ 和 $c_{Y/X}$ 都与这种基变换相容。对有限 syntomic 态射
$f : Y \to X$，若 $X$ 局部诺特，我们仍用 $\Theta^p_{Y/X}$
表示刚刚得到的解。

\medskip\noindent
唯一性。假设有有限 syntomic 态射 $f: Y \to X$，使得
$X$ 在 $\Spec(\mathbf{Z})$ 上光滑，并且 $f$ 在 $X$ 的某个稠密开集上平展。
我们断言，在这种情形下，性质 (1) 唯一确定 $\Theta^p_{Y/X}$。
事实上，考虑映射
$$
\Omega^p_{X/\mathbf{Z}} \otimes_{\mathcal{O}_X} f_*\mathcal{O}_Y \to
f_*\Omega^p_{Y/\mathbf{Z}} \to
\Omega^p_{X/\mathbf{Z}}
$$
层 $\Omega^p_{X/\mathbf{Z}}$ 是无挠的（由所假设的光滑性），
所以只需检验 $\Theta^p_{Y/X}$ 在 $f$ 为平展的稠密开集上的限制被唯一确定；
也就是说，可以假设 $f$ 是平展的。然而，如果 $f$ 是平展的，那么
$f^*\Omega_{X/\mathbf{Z}} = \Omega_{Y/\mathbf{Z}}$
，因而所示等式中的第一个箭头是同构。由于复合已经确定，这就保证了唯一性。

\medskip\noindent
令 $f : Y \to X$ 是局部诺特概形之间的有限 syntomic 态射，并令
$x \in X$。由判别式，引理
\ref{discriminant-lemma-locally-comes-from-universal-finite}
，可以找到 $d \geq 1$ 以及交换图
$$
\xymatrix{
Y \ar[d] &
V \ar[d] \ar[l] \ar[r] &
V_d \ar[d] \\
X &
U \ar[l] \ar[r] &
U_d
}
$$
使得 $x \in U \subset X$ 是开集，$V = f^{-1}(U)$，并且
$V = U \times_{U_d} V_d$。因此 $\Theta^p_{Y/X}|_V$
是映射 $\Theta^p_{V_d/U_d}$ 的拉回。不过，由上面对唯一性的讨论以及
判别式，引理
\ref{discriminant-lemma-universal-finite-syntomic-smooth} 和
\ref{discriminant-lemma-universal-finite-syntomic-etale}
，映射 $\Theta^p_{V_d/U_d}$ 被要求 (1) 唯一确定。因此唯一性成立。

\medskip\noindent
至此，我们已经知道：对于所有有限 syntomic 态射 $Y \to X$，
只要 $X$ 局部诺特，解就存在且唯一。现在可以给出与引理
\ref{derham-lemma-Garel-upstairs} 的证明类似的论证，将结果推广到一般的 $X$。
不过，也可以直接使用绝对诺特逼近来完成证明。事实上，为构造
$\Theta^p_{Y/X}$，只需在 $X$ 上扎里斯基局部地构造它
（同时还需证明唯一性）。因此可以假设 $X$ 是仿射的（略去一个小细节）。
然后可将 $X = \lim_{i \in I} X_i$ 写成由有向集 $I$ 标号的诺特仿射概形的极限。
由代数，引理 \ref{algebra-lemma-colimit-category-fp-algebras}，
可以找到 $0 \in I$ 以及仿射概形之间的有限表示态射
$f_0 : Y_0 \to X_0$，它在 $X$ 上的基变换是 $Y \to X$。
增大 $0$ 后，可以假设 $Y_0 \to X_0$ 是有限 syntomic 态射；见代数，引理
\ref{algebra-lemma-colimit-lci} 和
\ref{algebra-lemma-colimit-finite}。对于 $i \geq 0$，基变换
$f_i : Y_i = Y_0 \times_{X_0} X_i \to X_i$ 也为有限 syntomic 态射。于是
$$
\Gamma(X, f_*\Omega^p_{Y/\mathbf{Z}}) =
\Gamma(Y, \Omega^p_{Y/\mathbf{Z}}) =
\colim_{i \geq 0} \Gamma(Y_i, \Omega^p_{Y_i/\mathbf{Z}}) =
\colim_{i \geq 0} \Gamma(X_i, f_{i, *}\Omega^p_{Y_i/\mathbf{Z}})
$$
因此我们可以（而且必须）把 $\Theta^p_{Y/X}$ 定义为映射
$\Theta^p_{Y_i/X_i}$ 的余极限。此映射与任意笛卡尔图
$$
\xymatrix{
Y' \ar[r] \ar[d] & Y \ar[d] \\
X' \ar[r] & X
}
$$
相容，其中 $X'$ 为仿射概形；因为由上述论证，我们已知这对诺特仿射概形成立
（略去一个小细节；提示：如果还写成
$X' = \lim_{j \in J} X'_j$，那么对每个 $i \in I$，存在 $j \in J$
以及态射 $X'_j \to X_i$，它与态射 $X' \to X$ 相容）。证明完毕。
\end{proof}

\begin{proposition}
\label{derham-proposition-Garel}
\begin{reference}
\cite{Garel}
\end{reference}
设 $f : Y \to X$ 是概形的有限 syntomic 态射。引理
\ref{derham-lemma-Garel} 中的映射 $\Theta^p_{Y/X}$ 定义复形映射
$$
\Theta_{Y/X} :
f_*\Omega^\bullet_{Y/\mathbf{Z}}
\longrightarrow
\Omega^\bullet_{X/\mathbf{Z}}
$$
并具有下列性质：
\begin{enumerate}
\item 在次数 $0$ 上得到
$\text{Trace}_f : f_*\mathcal{O}_Y \to \mathcal{O}_X$；见判别式，第
\ref{discriminant-section-discriminant} 节；
\item 有
$\Theta_{Y/X}(\omega \wedge \eta) = \omega \wedge \Theta_{Y/X}(\eta)$
，其中 $\omega$ 属于 $\Omega^\bullet_{X/\mathbf{Z}}$，$\eta$
属于 $f_*\Omega^\bullet_{Y/\mathbf{Z}}$；
\item 如果 $f$ 是基概形 $S$ 上的态射，那么
$\Theta_{Y/X}$ 诱导复形映射
$f_*\Omega^\bullet_{Y/S} \to \Omega^\bullet_{X/S}$.
\end{enumerate}
\end{proposition}

\begin{proof}
由判别式，引理
\ref{discriminant-lemma-locally-comes-from-universal-finite}
，对每个 $x \in X$，可以找到 $d \geq 1$ 以及交换图
$$
\xymatrix{
Y \ar[d] &
V \ar[d] \ar[l] \ar[r] &
V_d \ar[d] \ar[r] &
Y_d = \Spec(B_d) \ar[d] \\
X &
U \ar[l] \ar[r] &
U_d \ar[r] &
X_d = \Spec(A_d)
}
$$
使得 $x \in U \subset X$ 是仿射开集，$V = f^{-1}(U)$，并且
$V = U \times_{U_d} V_d$。写成 $U = \Spec(A)$ 和 $V = \Spec(B)$；
注意 $B = A \otimes_{A_d} B_d$，并回忆
$B_d = A_d e_1 \oplus \ldots \oplus A_d e_d$。假设给定
$a_1, \ldots, a_r \in A$ 和 $b_1, \ldots, b_s \in B$。
可以写成 $b_j = \sum a_{j, l} e_l$，其中 $a_{j, l} \in A$。
置 $N = r + sd$，并考虑分解
$$
\xymatrix{
V \ar[r] \ar[d] &
V' = \mathbf{A}^N \times V_d \ar[r] \ar[d] &
V_d \ar[d] \\
U \ar[r]&
U' = \mathbf{A}^N \times U_d \ar[r] &
U_d
}
$$
这里右下方的水平箭头由态射 $U \to U_d$（来自前面的图）和态射
$U \to \mathbf{A}^N$ 给出，后者由
$a_1, \ldots, a_r, a_{1, 1}, \ldots, a_{s, d}$ 给定。
于是函数 $a_1, \ldots, a_r$ 属于映射
$\Gamma(U', \mathcal{O}_{U'}) \to \Gamma(U, \mathcal{O}_U)$
的像，而函数 $b_1, \ldots, b_s$ 属于映射
$\Gamma(V', \mathcal{O}_{V'}) \to \Gamma(V, \mathcal{O}_V)$.
的像。这样可见，对这些群中任意有限多个元素\footnote{毕竟，
所有这些元素都是如下形式的元素的有限和：
$a_0 \text{d}a_1 \wedge \ldots \wedge \text{d}a_i$，其中
$a_0, \ldots, a_i \in A$，或者如下形式的元素的有限和：
$b_0 \text{d}b_1 \wedge \ldots \wedge \text{d}b_j$，其中
$b_0, \ldots, b_j \in B$。}：
$$
\Gamma(V, \Omega^i_{Y/\mathbf{Z}}),\quad i = 0, 1, 2, \ldots
\quad\text{且}\quad
\Gamma(U, \Omega^j_{X/\mathbf{Z}}),\quad j = 0, 1, 2, \ldots
$$
可以找到如上的分解 $V \to V' \to V_d$ 和
$U \to U' \to U_d$，其中 $V' = \mathbf{A}^N \times V_d$ 且
$U' = \mathbf{A}^N \times U_d$，使得这些截面都是下列群中截面的拉回：
$$
\Gamma(V', \Omega^i_{V'/\mathbf{Z}}),\quad i = 0, 1, 2, \ldots
\quad\text{且}\quad
\Gamma(U', \Omega^j_{U'/\mathbf{Z}}),\quad j = 0, 1, 2, \ldots
$$
其结果是，为检验
$\text{d} \circ \Theta_{Y/X} = \Theta_{Y/X} \circ \text{d}$
，只需检验这对 $\Theta_{V'/U'}$ 成立。引理的性质 (2) 也同样如此。

\medskip\noindent
由判别式，引理
\ref{discriminant-lemma-universal-finite-syntomic-smooth} 和
\ref{discriminant-lemma-universal-finite-syntomic-etale}
，概形 $U_d$ 是光滑的，且态射 $V_d \to U_d$ 在 $U_d$
的某个稠密开集上平展。因此态射 $V' \to U'$ 也具有同样的性质。
由于 $\Omega_{U'/\mathbf{Z}}$ 是局部自由的，故
$\Omega^p_{U'/\mathbf{Z}}$ 是无挠的；所以只需把所需关系限制到
$V'$ 为有限平展的开集后加以检验。随后可以在满平展基变换之后检验这些关系。
因而可以分裂该有限平展覆盖，并假设所考察的态射形如
$$
\coprod\nolimits_{i = 1, \ldots, d} W \longrightarrow W
$$
，其中 $W$ 在 $\mathbf{Z}$ 上光滑。在这种情形下，我们构造的任何局部性质
（只要确实成立）都可直接检验。这就完成了 (1) 和 (2) 的证明。

\medskip\noindent
最后注意到，(3) 由 (2) 推出，因为 $\Omega_{Y/S}$ 是
$\Omega_{Y/\mathbf{Z}}$ 关于如下子模的商：该子模由
$\Omega_{S/\mathbf{Z}}$ 的局部截面的拉回生成。
\end{proof}

\begin{example}
\label{derham-example-Garel}
设 $A$ 是一个环，并令
$f = x^d + \sum_{0 \leq i < d} a_{d - i} x^i \in A[x]$。
令 $B = A[x]/(f)$。由命题 \ref{derham-proposition-Garel}，有复形态射
$$
\Theta_{B/A} : \Omega^\bullet_B \longrightarrow \Omega^\bullet_A
$$
特别地，如果 $t \in B$ 表示 $x \in A[x]$ 的像，就可以考虑元素
$$
\Theta_{B/A}(t^i\text{d}t) \in \Omega^1_A,\quad i = 0, \ldots, d - 1
$$
这些元素是什么？按照命题 \ref{derham-proposition-Garel} 的证明中使用的同一原理，
只需在泛情形下计算，即 $A = \mathbf{Z}[a_1, \ldots, a_d]$ 时；
甚至当 $A$ 被分式域 $\mathbf{Q}(a_1, \ldots, a_d)$ 代替时。
形式地写成
$$
f = \prod\nolimits_{i = 1, \ldots, d} (x - \alpha_i)
$$
可见在 $\mathbf{Q}(\alpha_1, \ldots, \alpha_d)$ 上，代数 $B$ 分裂为：
$$
\mathbf{Q}(a_0, \ldots, a_{d - 1})[x]/(f) 
\longrightarrow
\prod\nolimits_{i = 1, \ldots, d} \mathbf{Q}(\alpha_1, \ldots, \alpha_d),
\quad
t \longmapsto (\alpha_1, \ldots, \alpha_d)
$$
因此，例如
$$
\Theta(\text{d}t) = \sum \text{d} \alpha_i = - \text{d}a_1
$$
接下来有
$$
\Theta(t\text{d}t) = \sum \alpha_i \text{d}\alpha_i =
a_1 \text{d} a_1 - \text{d}a_2
$$
接下来有
$$
\Theta(t^2\text{d}t) = \sum \alpha_i^2 \text{d}\alpha_i =
- a_1^2 \text{d} a_1 + a_1 \text{d}a_2 + a_2 \text{d}a_1 - \text{d}a_3
$$
（不计可能的计算错误），依此类推。这提示：如果 $f(x) = x^d - a$，那么
$$
\Theta_{B/A}(t^i\text{d}t) =
\left\{
\begin{matrix}
0 & \text{若} & i = 0, \ldots, d - 2 \\
\text{d}a & \text{若} & i = d - 1
\end{matrix}
\right.
$$
在 $\Omega_A$ 中成立。这确实如此，因为在这个特例中，可以对扩张
$\mathbf{Q}(a)[x]/(x^d - a)$ 进行计算来直接验证。
\end{example}

\begin{lemma}
\label{derham-lemma-Garel-map-frobenius-smooth-char-p}
设 $p$ 是素数。设 $X \to S$ 是相对维数为 $d$ 的光滑态射，
其中这些概形的特征为 $p$。相对 Frobenius 态射
$F_{X/S} : X \to X^{(p)}$（属于 $X/S$；簇，定义
\ref{varieties-definition-relative-frobenius}）是有限 syntomic 态射，
且相应的映射
$$
\Theta_{X/X^{(p)}} :
F_{X/S, *}\Omega^\bullet_{X/S} \to \Omega^\bullet_{X^{(p)}/S}
$$
除次数 $d$ 外在所有次数上都为零，而在该次数上定义一个满射。
\end{lemma}

\begin{proof}
由簇，引理 \ref{varieties-lemma-relative-frobenius-finite}，
$F_{X/S}$ 是有限态射。为了证明 $F_{X/S}$ 平坦，只需证明所有纤维间的态射
$F_{X/S, s} : X_s \to X^{(p)}_s$ 对每个 $s \in S$ 都平坦；
见更多态射，定理 \ref{more-morphisms-theorem-criterion-flatness-fibre}。
$X_s \to X^{(p)}_s$ 的平坦性由代数，引理
\ref{algebra-lemma-CM-over-regular-flat} 得出（并使用已经证明的有限性）。
由更多态射，引理 \ref{more-morphisms-lemma-lci-permanence}，
态射 $F_{X/S}$ 是局部完全交态射。因此 $F_{X/S}$ 是有限 syntomic 态射
（见更多态射，引理 \ref{more-morphisms-lemma-flat-lci}）。

\medskip\noindent
对每一点 $x \in X$，可以选取交换图
$$
\xymatrix{
X \ar[d] & U \ar[l] \ar[d]_\pi \\
S & \mathbf{A}^d_S \ar[l]
}
$$
其中 $\pi$ 是平展的，且 $x \in U$ 是 $X$ 中的开集；见态射，引理
\ref{morphisms-lemma-smooth-etale-over-affine-space}。注意，
$\mathbf{A}^d_S \to \mathbf{A}^d_S$, $(x_1, \ldots, x_d) \mapsto
(x_1^p, \ldots, x_d^p)$ 是 $\mathcal{A}^d_S$ 在 $S$ 上的相对 Frobenius。
交换图
$$
\xymatrix{
U \ar[d]_\pi \ar[r]_{F_{X/S}} & U^{(p)} \ar[d]^{\pi^{(p)}} \\
\mathbf{A}^d_S \ar[r]^{x_i \mapsto x_i^p} & \mathbf{A}^d_S
}
$$
来自簇，引理 \ref{varieties-lemma-relative-frobenius-endomorphism-identity}；
对 $\pi : U \to \mathbf{A}^d_S$，由平展态射，引理
\ref{etale-lemma-relative-frobenius-etale}，该图是笛卡尔图。
由于 $\Theta$ 的构造与基变换相容，且
$\Omega_{U/S} = \pi^*\Omega_{\mathbf{A}^d_S/S}$
（引理 \ref{derham-lemma-etale}），所以只需对 $\mathbf{A}^d_S$ 证明本引理。

\medskip\noindent
设 $A$ 是特征 $p$ 的环。考虑唯一的 $A$-代数同态
$A[y_1, \ldots, y_d] \to A[x_1, \ldots, x_d]$，它把 $y_i$ 映到
$x_i^p$。上述论证把问题化为计算映射
$$
\Theta^i : \Omega^i_{A[x_1, \ldots, x_d]/A} \to
\Omega^i_{A[y_1, \ldots, y_d]/A}
$$
我们建议读者亲自完成这一情形下的计算。如同例 \ref{derham-example-Garel}，
可以把它化为在泛情形下计算 $\Theta^i$ 的公式
$$
\mathbf{Z}[y_1, \ldots, y_d] \to \mathbf{Z}[x_1, \ldots, x_d],\quad
y_i \mapsto x_i^p
$$
进而，可以在复数情形下计算 $\Theta^i$ 的公式，即对 $\mathbf{C}$-代数映射
$$
\mathbf{C}[y_1, \ldots, y_d] \to \mathbf{C}[x_1, \ldots, x_d],\quad
y_i \mapsto x_i^p
$$
甚至可以逆转 $x_1, \ldots, x_d$ 和 $y_1, \ldots, y_d$。
在这种情形下，有 $\text{d}x_i = p^{-1} x_i^{- p + 1}\text{d}y_i$。
因此可见
\begin{align*}
\Theta^i(
x_1^{e_1} \ldots x_d^{e_d} \text{d}x_1 \wedge \ldots \wedge \text{d}x_i)
& =
p^{-i} \Theta^i(
x_1^{e_1 - p + 1} \ldots x_i^{e_i - p + 1} x_{i + 1}^{e_{i + 1}} \ldots
x_d^{e_d} \text{d}y_1 \wedge \ldots \wedge \text{d}y_i ) \\
& =
p^{-i} \text{Trace}(x_1^{e_1 - p + 1} \ldots x_i^{e_i - p + 1}
x_{i + 1}^{e_{i + 1}} \ldots x_d^{e_d})
\text{d}y_1 \wedge \ldots \wedge \text{d}y_i
\end{align*}
，这里使用了 $\Theta^i$ 的性质。初等计算表明，上式中的迹为零，除非
$e_1, \ldots, e_i$ 都同余于 $-1$（模 $p$），且
$e_{i + 1}, \ldots, e_d$ 都能被 $p$ 整除。此外，在这种情形下得到
$$
p^{d - i} y_1^{(e_1 - p + 1)/p} \ldots y_i^{(e_i - p + 1)/p}
y_{i + 1}^{e_{i + 1}/p} \ldots y_d^{e_d/p}
\text{d}y_1 \wedge \ldots \wedge \text{d}y_i
$$
由此可知，在特征 $p$ 中，除非 $d = i$，所得结果为零；
而在这种情况下，可以得到每个可能的 $d$-形式。
\end{proof}








\section{庞加莱对偶}
\label{derham-section-poincare-duality}

\noindent
本节证明域上的光滑固有概形的德拉姆上同调所满足的庞加莱对偶。
我们先说明霍奇上同调的情形。

\begin{lemma}
\label{derham-lemma-duality-hodge}
设 $k$ 是一个域。设 $X$ 是 $k$ 上非空、光滑、固有且等维的概形，
其维数为 $d$。存在 $k$-线性映射
$$
t : H^d(X, \Omega^d_{X/k}) \longrightarrow k
$$
；它在先与 $H^0(X, \mathcal{O}_X)$ 中某个单位的乘法复合的意义下唯一，
并具有如下性质：对所有 $p, q$，配对
$$
H^q(X, \Omega^p_{X/k}) \times H^{d - q}(X, \Omega^{d - p}_{X/k})
\longrightarrow
k, \quad
(\xi, \xi') \longmapsto t(\xi \cup \xi')
$$
是完美配对。
\end{lemma}

\begin{proof}
由概形对偶，引理 \ref{duality-lemma-duality-proper-over-field}，有
$\omega_X^\bullet = \Omega^d_{X/k}[d]$。由于 $\Omega_{X/k}$
是秩为 $d$ 的局部自由模（态射，引理
\ref{morphisms-lemma-smooth-omega-finite-locally-free}），有
$$
\Omega^d_{X/k} \otimes_{\mathcal{O}_X} (\Omega^p_{X/k})^\vee
\cong
\Omega^{d - p}_{X/k}
$$
因此，由概形对偶，引理
\ref{duality-lemma-duality-proper-over-field-perfect}，得到使该陈述成立的
$k$-线性映射 $t : H^d(X, \Omega^d_{X/k}) \to k$。特别地，配对
$H^0(X, \mathcal{O}_X) \times H^d(X, \Omega^d_{X/k}) \to k$
是完美的。这意味着任意 $k$-线性映射
$t' : H^d(X, \Omega^d_{X/k}) \to k$ 都形如
$\xi \mapsto t(g\xi)$，其中 $g \in H^0(X, \mathcal{O}_X)$。
当然，为使 $t'$ 仍在 $H^0(X, \mathcal{O}_X)$ 与
$H^d(X, \Omega^d_{X/k})$ 之间给出对偶，$g$ 必须是单位。
把配对 $\langle -, - \rangle_{p, q}$ 定义为用 $t$ 构造的配对，
把配对 $\langle -, - \rangle'_{p, q}$ 定义为用 $t'$ 构造的配对。显然有
$$
\langle \xi, \xi' \rangle'_{p, q} =
\langle g\xi, \xi' \rangle_{p, q}
$$
，其中 $\xi \in H^q(X, \Omega^p_{X/k})$ 且
$\xi' \in H^{d - q}(X, \Omega^{d - p}_{X/k})$。由于 $g$ 是单位，
即为可逆元，所以用 $t'$ 代替 $t$，对所有 $p, q$
仍得到完美配对。
\end{proof}

\begin{lemma}
\label{derham-lemma-bottom-part-degenerates}
设 $k$ 是一个域，设 $X$ 是 $k$ 上光滑固有概形。映射
$$
\text{d} : H^0(X, \mathcal{O}_X) \to H^0(X, \Omega^1_{X/k})
$$
为零。
\end{lemma}

\begin{proof}
由于 $X$ 在 $k$ 上光滑，它在 $k$ 上几何既约；见簇，引理
\ref{varieties-lemma-smooth-geometrically-normal}。因此
$H^0(X, \mathcal{O}_X) = \prod k_i$ 是有限多个有限可分域扩张
$k_i/k$ 的乘积；见簇，引理
\ref{varieties-lemma-proper-geometrically-reduced-global-sections}。由此
$\Omega_{H^0(X, \mathcal{O}_X)/k} = \prod \Omega_{k_i/k} = 0$
（例如见代数，引理
\ref{algebra-lemma-characterize-separable-algebraic-field-extensions}）。
由于引理中的映射分解为
$$
H^0(X, \mathcal{O}_X) \to
\Omega_{H^0(X, \mathcal{O}_X)/k} \to
H^0(X, \Omega_{X/k})
$$
（由德拉姆复形的函子性；见第 \ref{derham-section-de-rham-complex} 节），结论成立。
\end{proof}

\begin{lemma}
\label{derham-lemma-top-part-degenerates}
设 $k$ 是一个域。设 $X$ 是 $k$ 上维数为 $d$ 的等维光滑固有概形。映射
$$
\text{d} : H^d(X, \Omega^{d - 1}_{X/k}) \to H^d(X, \Omega^d_{X/k})
$$
为零。
\end{lemma}

\begin{proof}
很容易猜想这可由引理 \ref{derham-lemma-bottom-part-degenerates} 与
\ref{derham-lemma-duality-hodge} 结合推出。然而这并不可行，因为映射
$\mathcal{O}_X \to \Omega^1_{X/k}$ 和
$\Omega^{d - 1}_{X/k} \to \Omega^d_{X/k}$ 不是 $\mathcal{O}_X$-线性的，
所以不能利用概形对偶，备注
\ref{duality-remark-coherent-duality-proper-over-field}
中讨论的函子性来断定引理 \ref{derham-lemma-bottom-part-degenerates}
中的映射与本引理中的映射互为对偶。

\medskip\noindent
可以用 $X$ 的一个连通分支代替 $X$，因此可以假设 $X$ 不可约。由簇，引理
\ref{varieties-lemma-smooth-geometrically-normal} 和
\ref{varieties-lemma-proper-geometrically-reduced-global-sections}
可知 $k' = H^0(X, \mathcal{O}_X)$ 是有限可分扩张 $k'/k$。
由于 $\Omega_{k'/k} = 0$（例如见代数，引理
\ref{algebra-lemma-characterize-separable-algebraic-field-extensions}），
有 $\Omega_{X/k} = \Omega_{X/k'}$（见态射，引理
\ref{morphisms-lemma-triangle-differentials}）。
因此可以把 $k$ 替换为 $k'$，并假设 $H^0(X, \mathcal{O}_X) = k$。

\medskip\noindent
假设 $H^0(X, \mathcal{O}_X) = k$。由引理 \ref{derham-lemma-duality-hodge}，
有 $\dim H^d(X, \Omega^d_{X/k}) = 1$。
先假设 $k$ 的特征是素数 $p$。把相对 Frobenius 态射
$F_{X/k} : X \to X^{(p)}$（属于 $X$ 在 $k$ 上）记为上述符号；请记住引理
\ref{derham-lemma-Garel-map-frobenius-smooth-char-p} 中关于该态射证明的事实。
考虑交换图
$$
\xymatrix{
H^d(X, \Omega^{d - 1}_{X/k}) \ar[d] \ar[r] &
H^d(X^{(p)}, F_{X/k, *}\Omega^{d - 1}_{X/k}) \ar[d] \ar[r]_{\Theta^{d - 1}} &
H^d(X^{(p)}, \Omega^{d - 1}_{X^{(p)}/k}) \ar[d] \\
H^d(X, \Omega^d_{X/k}) \ar[r] &
H^d(X^{(p)}, F_{X/k, *}\Omega^d_{X/k}) \ar[r]^{\Theta^d} &
H^d(X^{(p)}, \Omega^d_{X^{(p)}/k})
}
$$
由于 $F_{X/k}$ 是有限的，左边两个水平箭头是同构；见概形上同调，引理
\ref{coherent-lemma-relative-affine-cohomology}。右方块交换，因为
$\Theta_{X^{(p)}/X}$ 是复形态射，且 $\Theta^{d - 1}$ 为零。
因此只需证明 $\Theta^d$ 非零（由上面的讨论，映射 $\Theta^d$
的源维数为 $1$）。然而我们知道
$$
\Theta^d : F_{X/k, *}\Omega^d_{X/k} \to \Omega^d_{X^{(p)}/k}
$$
是满射，因而应用右正合函子 $H^d(X^{(p)}, -)$ 后仍为满射
（右正合性来自次数超过 $d$ 的上同调消失；见上同调，命题
\ref{cohomology-proposition-vanishing-Noetherian}）。
最后，$H^d(X^{(d)}, \Omega^d_{X^{(d)}/k})$ 非零，例如因为它对偶于
$H^0(X^{(d)}, \mathcal{O}_{X^{(p)}})$；这是由引理
\ref{derham-lemma-duality-hodge} 应用于 $X^{(p)}$（在 $k$ 上）得出的。
这就完成了该情形的证明。

\medskip\noindent
最后，假设 $k$ 的特征为 $0$。可以把 $k$ 写成其有限型
$\mathbf{Z}$-子代数 $R$ 的滤过余极限。对其中某一个，可以找到概形的笛卡尔图
$$
\xymatrix{
X \ar[d] \ar[r] & Y \ar[d] \\
\Spec(k) \ar[r] & \Spec(R)
}
$$
使得 $Y \to \Spec(R)$ 是相对维数为 $d$ 的光滑固有态射。
见极限，引理 \ref{limits-lemma-descend-finite-presentation}、
\ref{limits-lemma-descend-smooth}、\ref{limits-lemma-descend-dimension-d} 和
\ref{limits-lemma-eventually-proper}。模
$M^{i, j} = H^j(Y, \Omega^i_{Y/R})$ 是有限 $R$-模；见概形上同调，引理
\ref{coherent-lemma-proper-over-affine-cohomology-finite}。
因此，用 $R$ 的一个局部化代替它之后，可以假设所有这些模都是有限自由的。由
$M^{i, j} \otimes_R k = H^j(X, \Omega^i_{X/k})$
；这是平坦基变换（概形上同调，引理
\ref{coherent-lemma-flat-base-change-cohomology}）的结果。
因此只需证明 $M^{d - 1, d} \to M^{d, d}$ 为零。这是整环上的有限自由模之间的映射，
所以只需找到稠密的素理想集合 $\mathfrak p \subset R$，使得与
$\kappa(\mathfrak p)$ 张量后所得映射为零。由于 $R$ 在 $\mathbf{Z}$
上有限型，可以取剩余域具有正特征的素理想 $\mathfrak p$
所组成的集合（略去细节）。注意，
$$
M^{d - 1, d} \otimes_R \kappa(\mathfrak p) =
H^d(Y_{\kappa(\mathfrak p)},
\Omega^{d - 1}_{Y_{\kappa(\mathfrak p)}/\kappa(\mathfrak p)})
$$
；例如这由极限，引理
\ref{limits-lemma-higher-direct-images-zero-above-dimension-fibre} 得出。
对 $M^{d, d}$ 也同样如此。因此
$M^{d - 1, d} \otimes_R \kappa(\mathfrak p) \to
M^{d, d} \otimes_R \kappa(\mathfrak p)$
由上面已经处理的正特征情形，此映射为零。
\end{proof}

\begin{proposition}
\label{derham-proposition-poincare-duality}
设 $k$ 是一个域。设 $X$ 是 $k$ 上非空、光滑、固有且等维的概形，
其维数为 $d$。存在 $k$-线性映射
$$
t : H^{2d}_{dR}(X/k) \longrightarrow k
$$
；它在先与 $H^0(X, \mathcal{O}_X)$ 中某个单位的乘法复合的意义下唯一，
并具有如下性质：对所有 $i$，配对
$$
H^i_{dR}(X/k) \times H_{dR}^{2d - i}(X/k)
\longrightarrow
k, \quad
(\xi, \xi') \longmapsto t(\xi \cup \xi')
$$
是完美配对。
\end{proposition}

\begin{proof}
由霍奇到德拉姆谱序列（第 \ref{derham-section-hodge-to-de-rham} 节）、
$\Omega^i_{X/k}$ 对 $i > d$ 的消失、上同调，命题
\ref{cohomology-proposition-vanishing-Noetherian} 中的消失，以及引理
\ref{derham-lemma-bottom-part-degenerates} 和
\ref{derham-lemma-top-part-degenerates} 的结果，可知
$H^0_{dR}(X/k) = H^0(X, \mathcal{O}_X)$，且
$H^d(X, \Omega^d_{X/k}) = H_{dR}^{2d}(X/k)$。
更确切地说，这些等同来自复形映射
$$
\Omega^\bullet_{X/k} \to \mathcal{O}_X[0]
\quad\text{且}\quad
\Omega^d_{X/k}[-d] \to \Omega^\bullet_{X/k}
$$
选择 $t : H_{dR}^{2d}(X/k) \to k$，使它经此等同对应于引理
\ref{derham-lemma-duality-hodge} 中那样的一个 $t$。于是无论如何，
该引理中所示的配对在 $i = 0$ 时是完美的。

\medskip\noindent
把 $\underline{k}$ 记为取值 $k$ 的常值层（在 $X$ 上）。
简记 $\Omega^\bullet = \Omega^\bullet_{X/k}$。
考虑映射 (\ref{derham-equation-wedge})，在当前情形下它写成
$$
\wedge :
\text{Tot}(\Omega^\bullet \otimes_{\underline{k}} \Omega^\bullet)
\longrightarrow
\Omega^\bullet
$$
对每个整数 $p = 0, 1, \ldots, d$，由于次数的原因，此映射在子复形
$\text{Tot}(\sigma_{> p} \Omega^\bullet \otimes_{\underline{k}}
\sigma_{\geq d - p} \Omega^\bullet)$ 上为零。因此，$\wedge$ 在子复形
$\text{Tot}(\Omega^\bullet \otimes_{\underline{k}}
\sigma_{\geq d - p}\Omega^\bullet)$ 上的限制通过复形映射
$$
\gamma_p :
\text{Tot}(\sigma_{\leq p} \Omega^\bullet \otimes_{\underline{k}}
\sigma_{\geq d - p} \Omega^\bullet)
\longrightarrow
\Omega^\bullet
$$
使用第 \ref{derham-section-cup-product} 节中的相同步骤，得到杯积
$$
H^i(X, \sigma_{\leq p} \Omega^\bullet) \times
H^{2d - i}(X, \sigma_{\geq d - p}\Omega^\bullet)
\longrightarrow
H_{dR}^{2d}(X, \Omega^\bullet)
$$
我们将对 $p$ 归纳证明：这些杯积经 $t$ 在
$H^i(X, \sigma_{\leq p} \Omega^\bullet)$ 与
$H^{2d - i}(X, \sigma_{\geq d - p}\Omega^\bullet)$ 之间诱导完美配对。
当 $p = d$ 时，这正是本命题的断言。

\medskip\noindent
基础情形是 $p = 0$。此时得到的恰是引理
\ref{derham-lemma-duality-hodge} 中 $H^i(X, \mathcal{O}_X)$ 与
$H^{d - i}(X, \Omega^d)$ 之间的配对，因而结论成立。

\medskip\noindent
归纳步骤。假设已知结论对 $p$ 成立。考虑可区别三角
$$
\Omega^{p + 1}[-p - 1] \to
\sigma_{\leq p + 1}\Omega^\bullet \to
\sigma_{\leq p}\Omega^\bullet \to
\Omega^{p + 1}[-p]
$$
以及可区别三角
$$
\sigma_{\geq d - p}\Omega^\bullet \to
\sigma_{\geq d - p - 1}\Omega^\bullet \to
\Omega^{d - p - 1}[-d + p + 1] \to
(\sigma_{\geq d - p}\Omega^\bullet)[1]
$$
注意，两者都是 $\underline{k}$-模层复形的同伦范畴中的可区别三角；
特别地，映射 $\sigma_{\leq p}\Omega^\bullet \to \Omega^{p + 1}[-p]$ 和
$\Omega^{d - p - 1}[-d + p + 1] \to (\sigma_{\geq d - p}\Omega^\bullet)[1]$
由实际的复形映射给出，即分别使用微分
$\Omega^p \to \Omega^{p + 1}$ 和
$\Omega^{d - p - 1} \to \Omega^{d - p}$。
考虑与这些可区别三角相伴的上同调长正合列
$$
\xymatrix{
H^{i - 1}(X, \sigma_{\leq p}\Omega^\bullet) \ar[d]_a \\
H^i(X, \Omega^{p + 1}[-p - 1]) \ar[d]_b \\
H^i(X, \sigma_{\leq p + 1}\Omega^\bullet) \ar[d]_c \\
H^i(X, \sigma_{\leq p}\Omega^\bullet) \ar[d]_d \\
H^{i + 1}(X, \Omega^{p + 1}[-p - 1])
}
\quad\quad
\xymatrix{
H^{2d - i  + 1}(X, \sigma_{\geq d - p}\Omega^\bullet) \\
H^{2d - i}(X, \Omega^{d - p - 1}[-d + p + 1]) \ar[u]_{a'} \\
H^{2d - i}(X, \sigma_{\geq d - p - 1}\Omega^\bullet) \ar[u]_{b'} \\
H^{2d - i}(X, \sigma_{\geq d - p}\Omega^\bullet) \ar[u]_{c'} \\
H^{2d - i - 1}(X, \Omega^{d - p - 1}[-d + p + 1]) \ar[u]_{d'}
}
$$
由归纳假设和引理 \ref{derham-lemma-duality-hodge}，我们知道上面在第一、第二、
第四和第五行的 $k$-向量空间之间构造的配对是完美的。由 $5$-引理，
为了证明中间行上同调群之间的配对是完美的，只需证明各对
$(a, a')$、$(b, b')$、$(c, c')$ 和 $(d, d')$
与给定配对相容（见下文）。

\medskip\noindent
先对 $(c, c')$ 证明这一点。这里只需注意到存在交换图
$$
\xymatrix{
\text{Tot}(\sigma_{\leq p} \Omega^\bullet \otimes_{\underline{k}}
\sigma_{\geq d - p} \Omega^\bullet) \ar[d]_{\gamma_p} &
\text{Tot}(\sigma_{\leq p + 1} \Omega^\bullet \otimes_{\underline{k}}
\sigma_{\geq d - p} \Omega^\bullet) \ar[l] \ar[d] \\
\Omega^\bullet &
\text{Tot}(\sigma_{\leq p + 1} \Omega^\bullet \otimes_{\underline{k}}
\sigma_{\geq d - p - 1} \Omega^\bullet) \ar[l]_-{\gamma_{p + 1}}
}
$$
因此，如果有 $\alpha \in H^i(X, \sigma_{\leq p + 1}\Omega^\bullet)$ 和
$\beta \in H^{2d - i}(X, \sigma_{\geq d - p}\Omega^\bullet)$，那么由杯积的函子性得到
$\gamma_p(\alpha \cup c'(\beta)) = \gamma_{p + 1}(c(\alpha) \cup \beta)$
。

\medskip\noindent
类似地，对 $(b, b')$ 使用交换图
$$
\xymatrix{
\text{Tot}(\sigma_{\leq p + 1} \Omega^\bullet \otimes_{\underline{k}}
\sigma_{\geq d - p - 1} \Omega^\bullet) \ar[d]_{\gamma_{p + 1}} &
\text{Tot}(\Omega^{p + 1}[-p - 1] \otimes_{\underline{k}}
\sigma_{\geq d - p - 1} \Omega^\bullet) \ar[l] \ar[d] \\
\Omega^\bullet &
\Omega^{p + 1}[-p - 1]
\otimes_{\underline{k}}
\Omega^{d - p - 1}[-d + p + 1] \ar[l]_-\wedge
}
$$
并作同样的论证。

\medskip\noindent
对 $(d, d')$ 使用交换图
$$
\xymatrix{
\Omega^{p + 1}[-p] \otimes_{\underline{k}}
\Omega^{d - p - 1}[-d + p] \ar[d] &
\text{Tot}(\sigma_{\leq p}\Omega^\bullet \otimes_{\underline{k}}
\Omega^{d - p - 1}[-d + p]) \ar[l] \ar[d] \\
\Omega^\bullet &
\text{Tot}(\sigma_{\leq p}\Omega^\bullet \otimes_{\underline{k}}
\sigma_{\geq d - p}\Omega^\bullet) \ar[l]
}
$$
并考察下列群中的上同调类：
$H^i(X, \sigma_{\leq p}\Omega^\bullet)$ 和
$H^{2d - i}(X, \Omega^{d - p - 1}[-d + p])$.
把 $i$ 改为 $i - 1$，就得到对 $(a, a')$ 的结论，
从而完成了这些配对是完美配对的证明。

\medskip\noindent
我们略去如下论证：$t$ 在先与 $H^0(X, \mathcal{O}_X)$
中某个单位的乘法复合的意义下唯一。
\end{proof}







\section{陈类}
\label{derham-section-chern-classes}

\noindent
目前已经证明的结果足以利用 Weil 上同调理论，第
\ref{weil-section-chern} 节中的讨论，在德拉姆上同调中构造陈类。

\begin{lemma}
\label{derham-lemma-chern-classes}
存在唯一的规则，它为每个拟紧拟分离概形 $X$ 指定全陈类
$$
c^{dR} :
K_0(\textit{Vect}(X))
\longrightarrow
\prod\nolimits_{i \geq 0} H^{2i}_{dR}(X/\mathbf{Z})
$$
并具有下列性质：
\begin{enumerate}
\item 有 $c^{dR}(\alpha + \beta) = c^{dR}(\alpha) c^{dR}(\beta)$，
其中 $\alpha, \beta \in K_0(\textit{Vect}(X))$；
\item 如果 $f : X \to X'$ 是拟紧拟分离概形之间的态射，那么
$c^{dR}(f^*\alpha) = f^*c^{dR}(\alpha)$；
\item 给定 $\mathcal{L} \in \Pic(X)$，有
$c^{dR}([\mathcal{L}]) = 1 + c_1^{dR}(\mathcal{L})$
\end{enumerate}
\end{lemma}

\noindent
该构造很容易推广到所有概形，但为此需要稍微加强 Weil 上同调理论，第
\ref{weil-section-chern} 节中的讨论。

\begin{proof}
应用 Weil 上同调理论，命题 \ref{weil-proposition-chern-class} 即可得到。

\medskip\noindent
令 $\mathcal{C}$ 为所有拟紧拟分离概形组成的范畴。它显然满足
Weil 上同调理论，第 \ref{weil-section-chern} 节的条件
(1)、(2) 以及 (3) 的 (a)、(b)、(c)。

\medskip\noindent
取逆变函子 $A$（从 $\mathcal{C}$ 到分次代数范畴），它把 $X$ 映到
$A(X) = \bigoplus_{i \geq 0} H_{dR}^{2i}(X/\mathbf{Z})$
，并赋予后者杯积。函子性在第 \ref{derham-section-de-rham-cohomology} 节讨论，
杯积在第 \ref{derham-section-cup-product} 节讨论。对于加性映射 $c_1^A$，
取第 \ref{derham-section-first-chern-class} 节构造的 $c_1^{dR}$。

\medskip\noindent
事实上，由引理 \ref{derham-lemma-cup-product-graded-commutative}
得到交换代数，这说明 $A$ 满足公理 (1)。

\medskip\noindent
为了检验 $A$ 的公理 (2)，只需检验
$H^*_{dR}(X \coprod Y/\mathbf{Z}) =  H^*_{dR}(X/\mathbf{Z}) \times
H^*_{dR}(Y/\mathbf{Z})$.
这是因为德拉姆上同调是取微分分次代数层（在扎里斯基拓扑中）的上同调构造的。

\medskip\noindent
对 $A$ 而言，公理 (3) 正是说取可逆模的第一陈类与拉回相容。
这由更一般的引理 \ref{derham-lemma-pullback-c1} 得出。

\medskip\noindent
对 $A$ 而言，公理 (4) 是我们在命题
\ref{derham-proposition-projective-space-bundle-formula} 中证明的射影空间丛公式。

\medskip\noindent
公理 (5)。设 $X$ 是拟紧拟分离概形，并设
$\mathcal{E} \to \mathcal{F}$ 是有限局部自由 $\mathcal{O}_X$-模之间的满射，
两者的秩分别为 $r + 1$ 和 $r$。把相应的包含态射记为
$i : P' = \mathbf{P}(\mathcal{F}) \to \mathbf{P}(\mathcal{E}) = P$。
这是 $X$ 上光滑射影概形之间的态射，它把 $P'$ 表示为 $P$
上的有效 Cartier 除子。因此由引理 \ref{derham-lemma-check-log-smooth}，
对 $P' \subset P$ 而言，其在 $\mathbf{Z}$ 上的对数极点复形有定义。
所以对于 $a \in A(P)$，若 $i^*a = 0$，则由引理
\ref{derham-lemma-log-complex-consequence}，有
$a \cup c_1^A(\mathcal{O}_P(P')) = 0$。证明完毕。
\end{proof}

\begin{remark}
\label{derham-remark-splitting-principle}
Weil 上同调理论，引理
\ref{weil-lemma-splitting-principle}（分裂原理）和
\ref{weil-lemma-chern-classes-E-tensor-L}（张量积的陈类）的类似结论，
对拟紧拟分离概形上的德拉姆陈类成立。这是显然的，因为我们已经在引理
\ref{derham-lemma-chern-classes} 的证明中说明，Weil 上同调理论，第
\ref{weil-section-chern} 节的所有公理均成立。
\end{remark}

\noindent
对 $\mathbf{Q}$ 上的概形，可以构造陈特征。

\begin{lemma}
\label{derham-lemma-chern-character}
存在唯一的规则，它为每个拟紧拟分离概形 $X$（在 $\mathbf{Q}$ 上）
指定陈特征
$$
ch^{dR} : K_0(\textit{Vect}(X)) \longrightarrow
\prod\nolimits_{i \geq 0} H_{dR}^{2i}(X/\mathbf{Q})
$$
并具有下列性质：
\begin{enumerate}
\item $ch^{dR}$ 对所有 $X$ 都是环同态；
\item 如果 $f : X' \to X$ 是 $\mathbf{Q}$ 上拟紧拟分离概形之间的态射，
那么 $f^* \circ ch^{dR} =  ch^{dR} \circ f^*$；
\item 给定 $\mathcal{L} \in \Pic(X)$，有
$ch^{dR}([\mathcal{L}]) = \exp(c_1^{dR}(\mathcal{L}))$。
\end{enumerate}
\end{lemma}

\noindent
该构造很容易推广到 $\mathbf{Q}$ 上的所有概形，但为此需要稍微加强
Weil 上同调理论，第 \ref{weil-section-chern} 节中的讨论。

\begin{proof}
完全仿照引理 \ref{derham-lemma-chern-classes} 的证明，可以说明
$\mathbf{Q}$ 上拟紧拟分离概形的范畴连同函子
$A^*(X) = \bigoplus_{i \geq 0} H_{dR}^{2i}(X/\mathbf{Q})$
满足 Weil 上同调理论，第 \ref{weil-section-chern} 节的公理。
此外，在这种情形下，$A(X)$ 是 $\mathbf{Q}$-代数（对所有 $X$）。
因此本引理由 Weil 上同调理论，命题
\ref{weil-proposition-chern-character} 得出。
\end{proof}







\section{一种 Weil 上同调理论}
\label{derham-section-weil}

\noindent
设 $k$ 是特征 $0$ 的域。本节证明函子
$$
X \longmapsto H^*_{dR}(X/k)
$$
定义了 $k$ 上以 $k$ 为系数的 Weil 上同调理论，其含义见
Weil 上同调理论，定义 \ref{weil-definition-weil-cohomology-theory}。
我们将检验本章前面的构造给出数据 (D0)、(D1)、(D2')，
并满足 Weil 上同调理论，第 \ref{weil-section-c1} 节的公理 (A1)--(A9)。

\medskip\noindent
在本节余下部分中，固定域 $k$（特征 $0$），并置 $F = k$。
取下列数据：
\begin{enumerate}
\item[(D0)] 对 $1$ 维 $F$-向量空间 $F(1)$，取
$F(1) = F = k$.
\item[(D1)] 对函子 $H^*$，取把光滑射影概形 $X$（在 $k$ 上）映到
$H^*_{dR}(X/k)$ 的函子。函子性在第
\ref{derham-section-de-rham-cohomology} 节讨论，杯积在第
\ref{derham-section-cup-product} 节讨论。由引理
\ref{derham-lemma-cup-product-graded-commutative}，得到分次交换 $F$-代数。
\item[(D2')] 对映射 $c_1^H : \Pic(X) \to H^2(X)(1)$，使用第
\ref{derham-section-first-chern-class} 节引入的德拉姆第一陈类。
\end{enumerate}
我们将证明公理 (A1)--(A9) 成立。

\medskip\noindent
本段使用 Weil 上同调理论，引理
\ref{weil-lemma-check-over-extension}，把公理的检验约化到 $k$
为代数闭域的情形。把 $k'$ 记为 $k$ 的代数闭包，置 $F' = k'$。
完全按照上述方式，得到 $k'$ 上系数域为 $F'$ 的数据 (D0)、(D1)、(D2')。
由引理 \ref{derham-lemma-proper-smooth-de-Rham}，存在函子性同构
$$
H_{dR}^{2d}(X/k) \otimes_k k'
\longrightarrow
H_{dR}^{2d}(X_{k'}/k')
$$
，其中 $X$ 是 $k$ 上光滑射影概形。此外，诸图
$$
\xymatrix{
\Pic(X) \ar[r]_{c^{dR}_1} \ar[d] & H_{dR}^2(X/k) \ar[d] \\
\Pic(X_{k'}) \ar[r]^{c^{dR}_1} & H_{dR}^2(X_{k'}/k')
}
$$
由引理 \ref{derham-lemma-pullback-c1} 交换。这就完成了约化的证明。

\medskip\noindent
假设 $k$ 是特征零的代数闭域。我们将对上面给出的数据
(D0)、(D1)、(D2') 证明公理 (A1)--(A9)。

\medskip\noindent
公理 (A1)。这里需要检验
$H^*_{dR}(X \coprod Y/k) =  H^*_{dR}(X/k) \times H^*_{dR}(Y/k)$.
这是因为德拉姆上同调是取微分分次代数层（在扎里斯基拓扑中）的上同调构造的。

\medskip\noindent
公理 (A2)。它正是说取可逆模的第一陈类与拉回相容。
这由更一般的引理 \ref{derham-lemma-pullback-c1} 得出。

\medskip\noindent
公理 (A3)。这由更一般的命题
\ref{derham-proposition-projective-space-bundle-formula} 得出。

\medskip\noindent
公理 (A4)。这由更一般的引理
\ref{derham-lemma-log-complex-consequence} 得出。

\medskip\noindent
至此，使用 Weil 上同调理论，引理
\ref{weil-lemma-chern-classes} 和
\ref{weil-lemma-cycle-classes}，已经得到陈特征和循环类映射
$$
\gamma :
\CH^*(X)
\longrightarrow
\bigoplus\nolimits_{i \geq 0} H^{2i}_{dR}(X/k)
$$
；对光滑射影概形 $X$（在 $k$ 上），它们是分次环同态，并与态射
$f : X \to Y$（在 $k$ 上光滑射影概形之间）的拉回相容。

\medskip\noindent
公理 (A5)。有 $H_{dR}^*(\Spec(k)/k) = k = F$（在次数 $0$ 上）。
由引理 \ref{derham-lemma-kunneth-de-rham}，
两个光滑射影 $k$-概形的积满足 K\"unneth 公式（注意，该陈述中的导出张量积
没有问题，因为是在域 $k$ 上张量）。

\medskip\noindent
公理 (A7)。这由命题 \ref{derham-proposition-blowup-split} 得出。

\medskip\noindent
公理 (A8)。设 $X$ 是 $k$ 上光滑射影概形。由 Weil 上同调理论，
第 \ref{weil-section-c1} 节对该公理的解释可知，
$k' = H^0(X, \mathcal{O}_X)$ 是有限可分 $k$-代数。于是
$H_{dR}^*(\Spec(k')/k) = k'$ 位于次数 $0$，因为
$\Omega_{k'/k} = 0$。由引理 \ref{derham-lemma-bottom-part-degenerates}，
还有 $H_{dR}^0(X, \mathcal{O}_X) = k'$，从而得到该公理。

\medskip\noindent
公理 (A6)。设 $X$ 是 $k$ 上非空光滑射影概形，等维且维数为 $d$。
把 $\Delta : X \to X \times_{\Spec(k)} X$
记为 $X$ 在 $k$ 上的对角态射。需要证明存在 $k$-线性映射
$$
\lambda : H_{dR}^{2d}(X/k) \longrightarrow k
$$
使得 $(1 \otimes \lambda)\gamma([\Delta]) = 1$ 在
$H^0_{dR}(X/k)$ 中成立。写成
$$
\gamma = \gamma([\Delta]) = \gamma_0 + \ldots + \gamma_{2d}
$$
，其中 $\gamma_i \in H_{dR}^i(X/k) \otimes_k H_{dR}^{2d - i}(X/k)$
是 K\"unneth 分量。问题在于证明存在线性映射
$\lambda : H_{dR}^{2d}(X/k) \to k$，使得
$(1 \otimes \lambda)\gamma_0 = 1$ 在 $H^0_{dR}(X/k)$ 中成立。

\medskip\noindent
令 $X = \coprod X_i$ 为 $X$ 分解成连通（因而不可约）分支的分解。
相应地，有 $\Delta = \coprod \Delta_i$，其中
$\Delta_i \subset X_i \times X_i$。于是
$$
\gamma([\Delta]) = \sum \gamma([\Delta_i])
$$
而且，$\gamma([\Delta_i])$ 对应于
$\Delta_i \subset X_i \times X_i$ 的类，这是经由分解
$$
H^*_{dR}(X \times X) = \prod\nolimits_{i, j} H^*_{dR}(X_i \times X_j)
$$
实现的。
略去细节；一种证明方法是利用 $\CH^0(X \times X)$ 中存在幂等元
$e_{i, j}$，它对应于开闭子概形 $X_i \times X_j$；并利用 $\gamma$ 是环同态，
它把 $e_{i, j}$ 映到所示上同调乘积分解中相应的幂等元。
如果能找到 $\lambda_i : H_{dR}^{2d}(X_i/k) \to k$，使得
$(1 \otimes \lambda_i)\gamma([\Delta_i]) = 1$ 在
$H^0_{dR}(X_i/k)$ 中成立，那么取 $\lambda = \sum \lambda_i$
就能对 $X$ 解决问题。因此可以并且确实假设 $X$ 不可约。

\medskip\noindent
证明不可约 $X$ 的公理 (A6)。由于 $k$ 代数闭，有
$H^0_{dR}(X/k) = k$，因为 $H^0(X, \mathcal{O}_X) = k$；
这里使用了 $X$ 是代数闭域上的射影簇（例如见簇，引理
\ref{varieties-lemma-proper-geometrically-reduced-global-sections}）。
令 $x \in X$ 为任意闭点。考虑笛卡尔图
$$
\xymatrix{
x \ar[d] \ar[r] & X \ar[d]^\Delta \\
X \ar[r]^-{x \times \text{id}} & X \times_{\Spec(k)} X
}
$$
$\gamma$ 与拉回的相容性表明，$\gamma([\Delta])$ 映到
$\gamma([x])$（在 $H_{dR}^{2d}(X/k)$ 中）；换言之，有
$\gamma_0 = 1 \otimes \gamma([x])$。由此得到两点：
(a) 类 $\gamma([x])$ 与 $x$ 无关；(b) 只需证明类
$\gamma([x])$ 非零；因而 (c) 只需找到任意零维循环
$\alpha$（在 $X$ 上），使得 $\gamma(\alpha) \not = 0$。为此，选择有限态射
$$
f : X \longrightarrow \mathbf{P}^d_k
$$
关于这种态射的存在性，见相交理论，第
\ref{intersection-section-projection} 节，特别是引理
\ref{intersection-lemma-projection-generically-finite}。注意，$f$
是有限 syntomic 态射（由更多态射，引理
\ref{more-morphisms-lemma-lci-permanence}，它是局部完全交态射；
由代数，引理 \ref{algebra-lemma-CM-over-regular-flat}，它是平坦的）。
由命题 \ref{derham-proposition-Garel}，有迹映射
$$
\Theta_f :
f_*\Omega^\bullet_{X/k}
\longrightarrow
\Omega^\bullet_{\mathbf{P}^d_k/k}
$$
，它与典范映射
$$
\Omega^\bullet_{\mathbf{P}^d_k/k}
\longrightarrow
f_*\Omega^\bullet_{X/k}
$$
的复合是乘以 $f$ 的次数。因此得到映射
$$
\Theta : H_{dR}^{2d}(X/k) \to H_{dR}^{2d}(\mathbf{P}^d_k/k)
$$
使得 $\Theta \circ f^*$ 是乘以某个正整数。因此，如果能找到
$\mathbf{P}^d_k$ 上类非零的零维循环，就可由 $\gamma$ 与拉回的相容性得出结论。
这由引理 \ref{derham-lemma-de-rham-cohomology-projective-space} 成立，
从而完成公理 (A6) 的证明。

\medskip\noindent
下文将不加说明地使用以下事实。第一，由 Weil 上同调理论，备注
\ref{weil-remark-trace}，映射
$\lambda_X : H^{2d}_{dR}(X/k) \to k$ 是唯一的。第二，
在公理 (A6) 的证明中已经看到
$\lambda_X(\gamma([x])) = 1$，只要 $X$ 不可约；也就是说，循环类映射
$\gamma : \CH^d(X) \to H_{dR}^{2d}(X/k)$ 与 $\lambda_X$
的复合是次数映射。

\medskip\noindent
公理 (A9)。设 $Y \subset X$ 是非空光滑等维射影概形
$X$（在 $k$ 上，维数为 $d$）上的非空光滑除子。需要证明图
$$
\xymatrix{
H_{dR}^{2d - 2}(X/k)
\ar[rrr]_{c^{dR}_1(\mathcal{O}_X(Y)) \cap -} \ar[d]_{restriction} & & &
H_{dR}^{2d}(X) \ar[d]^{\lambda_X} \\
H_{dR}^{2d - 2}(Y/k) \ar[rrr]^-{\lambda_Y} & & & k
}
$$
交换，其中 $\lambda_X$ 和 $\lambda_Y$ 如公理 (A6) 所述。
上面已经看到，如果把 $X = \coprod X_i$ 分解为连通（等价地，不可约）分支，
那么相应地 $\lambda_X = \sum \lambda_{X_i}$。
类似地，如果把 $Y = \coprod Y_j$ 分解为连通（等价地，不可约）分支，
那么 $\lambda_Y = \sum \lambda_{Y_j}$。此外，此时有
$\mathcal{O}_X(Y) = \otimes_j \mathcal{O}_X(Y_j)$，因而
$$
c_1^{dR}(\mathcal{O}_X(Y)) = \sum\nolimits_j
c^{dR}_1(\mathcal{O}_X(Y_j))
$$
在 $H_{dR}^2(X/k)$ 中成立。直接追图表明，只需在 $X$ 和 $Y$
都不可约时证明该图交换。此时 $H_{dR}^{2d - 2}(Y/k)$ 是 $1$ 维的，
因为由 Weil 上同调理论，引理 \ref{weil-lemma-poincare-duality}，
$Y$ 满足庞加莱对偶。由公理 (A4)，限制映射（左竖直箭头）的核
包含在与 $c^{dR}_1(\mathcal{O}_X(Y))$ 作杯积的映射之核中。
这意味着，只需找到一个上同调类
$a \in H_{dR}^{2d - 2}(X)$，其在 $Y$ 上的限制非零，并且对 $a$
该图交换。取任意丰沛可逆模 $\mathcal{L}$，并置
$$
a = c^{dR}_1(\mathcal{L})^{d - 1}
$$
于是 $a|_Y = c^{dR}_1(\mathcal{L}|_Y)^{d - 1}$，因而
$$
\lambda_Y(a|_Y) = \deg(c_1(\mathcal{L}|_Y)^{d - 1} \cap [Y])
$$
，这里使用了上面对 $\lambda_Y$ 的描述。由 Chow 同调，引理
\ref{chow-lemma-degrees-and-numerical-intersections} 结合簇，引理
\ref{varieties-lemma-ample-positive}，这是一个正整数。类似地，得到
$$
\lambda_X(c^{dR}_1(\mathcal{O}_X(Y)) \cap a) =
\deg(c_1(\mathcal{O}_X(Y)) \cap c_1(\mathcal{L})^{d - 1} \cap [X])
$$
由于几乎由定义就知道 $c_1(\mathcal{O}_X(Y)) \cap [X] = [Y]$，
因而有零维循环的等式
$$
(Y \to X)_*\left(c_1(\mathcal{L}|_Y)^{d - 1} \cap [Y]\right) =
c_1(\mathcal{O}_X(Y)) \cap c_1(\mathcal{L})^{d - 1} \cap [X]
$$
（在 $X$ 上）。因此这些循环具有相同次数，证明完毕。

\begin{proposition}
\label{derham-proposition-de-rham-is-weil}
设 $k$ 是特征零的域。把光滑射影概形 $X$（在 $k$ 上）映到
$H_{dR}^*(X/k)$ 的函子，是 Weil 上同调理论，含义见
Weil 上同调理论，定义 \ref{weil-definition-weil-cohomology-theory}。
\end{proposition}

\begin{proof}
在上面的讨论中，我们说明了数据 (D0)、(D1)、(D2') 满足
Weil 上同调理论，第 \ref{weil-section-c1} 节的公理 (A1)--(A9)。
因此结论由 Weil 上同调理论，命题 \ref{weil-proposition-get-weil} 得出。

\medskip\noindent
请不要阅读下文。在这些断言的证明中，我们还使用了引理
\ref{derham-lemma-proper-smooth-de-Rham}、
\ref{derham-lemma-pullback-c1},
\ref{derham-lemma-log-complex-consequence},
\ref{derham-lemma-kunneth-de-rham},
\ref{derham-lemma-bottom-part-degenerates} 和
\ref{derham-lemma-de-rham-cohomology-projective-space},
命题
\ref{derham-proposition-projective-space-bundle-formula},
\ref{derham-proposition-blowup-split} 和
\ref{derham-proposition-Garel},
Weil 上同调理论，引理
\ref{weil-lemma-check-over-extension},
\ref{weil-lemma-chern-classes},
\ref{weil-lemma-cycle-classes} 和
\ref{weil-lemma-poincare-duality},
Weil 上同调理论，备注 \ref{weil-remark-trace}，簇，引理
\ref{varieties-lemma-proper-geometrically-reduced-global-sections} 和
\ref{varieties-lemma-ample-positive},
相交理论，第 \ref{intersection-section-projection} 节及引理
\ref{intersection-lemma-projection-generically-finite}，
更多态射，引理 \ref{more-morphisms-lemma-lci-permanence}，
代数，引理 \ref{algebra-lemma-CM-over-regular-flat}，以及 Chow 同调，引理
\ref{chow-lemma-degrees-and-numerical-intersections}.
\end{proof}

\begin{remark}
\label{derham-remark-hodge-cohomology-is-weil}
完全按照上述方式，可以证明霍奇上同调
$X \mapsto H_{Hodge}^*(X/k)$ 配备 $c_1^{Hodge}$ 后确定一个 Weil 上同调理论。
如果将来需要这一点，我们会在这里精确表述并证明。这导致以下有趣结论：
如果 Weil 上同调理论的 Betti 数与所选 Weil 上同调理论无关
（在我们的域 $k$（特征 $0$）上），那么霍奇到德拉姆谱序列在
$E_1$ 页退化！当然，霍奇到德拉姆谱序列的退化是已知的
（例如见 \cite{Deligne-Illusie} 中精彩的代数证明），但它绝非容易的结果！
这提示证明 Betti 数的无关性同样困难；据我们所知，这仍是一个开放问题。
相关问题见 Weil 上同调理论，备注
\ref{weil-remark-betti-numbers-in-some-sense}。
\end{remark}







\section{闭浸入的 Gysin 映射}
\label{derham-section-gysin}

\noindent
本节定义闭浸入的 Gysin 映射。

\begin{remark}
\label{derham-remark-gysin-equations}
设 $X \to S$ 是概形的态射。设
$f_1, \ldots, f_c \in \Gamma(X, \mathcal{O}_X)$。设 $Z \subset X$
是由 $f_1, \ldots, f_c$ 截出的闭子概形。下面研究 {\it Gysin 映射}
\begin{equation}
\label{derham-equation-gysin}
\gamma^p_{f_1, \ldots, f_c} :
\Omega^p_{Z/S}
\longrightarrow
\mathcal{H}_Z^c(\Omega^{p + c}_{X/S})
\end{equation}
，其定义如下。给定局部截面 $\omega$（属于 $\Omega^p_{Z/S}$），
它是某个截面 $\tilde \omega$（属于 $\Omega^p_{X/S}$）的限制，置
$$
\gamma^p_{f_1, \ldots, f_c}(\omega) =
c_{f_1, \ldots, f_c}(\tilde \omega|_Z) \wedge
\text{d}f_1 \wedge \ldots \wedge \text{d}f_c
$$
，其中 $c_{f_1, \ldots, f_c} : \Omega^p_{X/S} \otimes \mathcal{O}_Z \to
\mathcal{H}_Z^c(\Omega^p_{X/S})$ 是概形的导出范畴，备注
\ref{perfect-remark-supported-map-c-equations} 中构造的映射。
此定义是良定义的：给定 $\omega$，可以把所选的
$\tilde \omega$ 改变如下形式的元素：
$\sum f_i \omega'_i + \sum \text{d}(f_i) \wedge \omega''_i$
，而该构造把这些元素映到零。
\end{remark}

\begin{lemma}
\label{derham-lemma-gysin-differential}
Gysin 映射 (\ref{derham-equation-gysin}) 与 $\Omega^\bullet_{X/S}$ 和
$\Omega^\bullet_{Z/S}$ 上的德拉姆微分相容。
\end{lemma}

\begin{proof}
一旦正确解释该陈述，它就由一个几乎平凡的计算得出。首先回忆，
函子 $\mathcal{H}^c_Z$ 在 $\mathcal{O}_X$-模范畴上计算时，
与 $X$ 上阿贝尔层范畴中类似定义的函子一致；见上同调，引理
\ref{cohomology-lemma-sections-support-abelian-unbounded}。
因此，微分 $\text{d} : \Omega^p_{X/S} \to \Omega^{p + 1}_{X/S}$
诱导映射
$\mathcal{H}^c_Z(\Omega^p_{X/S}) \to \mathcal{H}^c_Z(\Omega^{p + 1}_{X/S})$.
此外，概形的导出范畴，备注
\ref{perfect-remark-support-c-equations} 中扩张交错 {\v C}ech 复形的构造
也适用于阿贝尔层范畴。映射
$$
\Coker\left(\bigoplus \mathcal{F}_{1 \ldots \hat i \ldots c} \to
\mathcal{F}_{1 \ldots c}\right)
\longrightarrow
i_*\mathcal{H}^c_Z(\mathcal{F})
$$
用于概形的导出范畴，备注
\ref{perfect-remark-supported-map-c-equations} 中
$c_{f_1, \ldots, f_c}$ 的构造；它在 $X$ 上所有阿贝尔层的范畴中
是良定义且函子性的。因此本引理由等式
$$
\text{d}\left(
\frac{\tilde \omega \wedge \text{d}f_1 \wedge \ldots \wedge
\text{d}f_c}{f_1 \ldots f_c}\right) =
\frac{\text{d}(\tilde \omega) \wedge
\text{d}f_1 \wedge \ldots \wedge \text{d}f_c}{f_1 \ldots f_c}
$$
得出，而该等式显然成立。
\end{proof}

\begin{lemma}
\label{derham-lemma-gysin-global}
设 $X \to S$ 是概形的态射。设 $Z \to X$ 是有限表示的闭浸入，
其余法层 $\mathcal{C}_{Z/X}$ 是秩 $c$ 的局部自由模。于是存在典范映射
$$
\gamma^p : \Omega^p_{Z/S} \to \mathcal{H}^c_Z(\Omega^{p + c}_{X/S})
$$
；局部上，它由备注 \ref{derham-remark-gysin-equations} 中的映射
$\gamma^p_{f_1, \ldots, f_c}$ 给出。
\end{lemma}

\begin{proof}
这些假设表明，给定 $x \in Z \subset X$，存在开邻域 $U$（包含 $x$），
使得 $Z$ 由 $c$ 个元素 $f_1, \ldots, f_c \in \mathcal{O}_X(U)$ 截出。
因此只需证明：给定 $f_1, \ldots, f_c$ 和 $g_1, \ldots, g_c$
（它们位于 $\mathcal{O}_X(U)$ 中并截出 $Z \cap U$），映射
$\gamma^p_{f_1, \ldots, f_c}$ 与
$\gamma^p_{g_1, \ldots, g_c}$ 相同。为此，缩小 $U$ 后可以假设
$g_j = \sum a_{ji} f_i$，其中
$a_{ji} \in \mathcal{O}_X(U)$。于是由概形的导出范畴，引理
\ref{perfect-lemma-supported-map-determinant}，有
$c_{f_1, \ldots, f_c} = \det(a_{ji}) c_{g_1, \ldots, g_c}$。
另一方面，有
$$
\text{d}(g_1) \wedge \ldots \wedge \text{d}(g_c) \equiv
\det(a_{ji}) \text{d}(f_1) \wedge \ldots \wedge \text{d}(f_c)
\bmod (f_1, \ldots, f_c)\Omega^c_{X/S}
$$
合并这些关系，直接计算即得所需等式。
\end{proof}

\begin{lemma}
\label{derham-lemma-gysin-differential-global}
设 $X \to S$ 和 $i : Z \to X$ 如引理 \ref{derham-lemma-gysin-global} 所述。
Gysin 映射 $\gamma^p$ 与 $\Omega^\bullet_{X/S}$ 和
$\Omega^\bullet_{Z/S}$ 上的德拉姆微分相容。
\end{lemma}

\begin{proof}
可以局部地检验，而此时结论由引理 \ref{derham-lemma-gysin-differential} 得出。
\end{proof}

\begin{lemma}
\label{derham-lemma-gysin-projection}
设 $X \to S$ 和 $i : Z \to X$ 如引理 \ref{derham-lemma-gysin-global} 所述。
给定 $\alpha \in H^q(X, \Omega^p_{X/S})$，有
$\gamma^p(\alpha|_Z) = i^{-1}\alpha \wedge \gamma^0(1)$，这是在
$H^q(Z, \mathcal{H}^c_Z(\Omega^{p + c}_{X/S}))$.
其中的记号见证明。
\end{lemma}

\begin{proof}
限制 $\alpha|_Z$ 是由霍奇上同调的函子性给出的
$H^q(Z, \Omega^p_{Z/S})$ 中的元素。对上同调应用函子性并使用
$\gamma^p : \Omega^p_{Z/S} \to \mathcal{H}^c_Z(\Omega^{p + c}_{X/S})$
，得到 $\gamma^p(\alpha|_Z)$，它属于
$H^q(Z, \mathcal{H}^c_Z(\Omega^{p + c}_{X/S}))$。
这解释了公式的左边。

\medskip\noindent
为解释右边，先沿环化空间的映射
$i : (Z, i^{-1}\mathcal{O}_X) \to (X, \mathcal{O}_X)$ 拉回，
得到元素 $i^{-1}\alpha \in H^q(Z, i^{-1}\Omega^p_{X/S})$。
令 $\gamma^0(1) \in H^0(Z, \mathcal{H}_Z^c(\Omega^c_{X/S}))$
为 $1 \in H^0(Z, \mathcal{O}_Z) = H^0(Z, \Omega^0_{Z/S})$
在 $\gamma^0$ 下的像。使用杯积，得到元素
$$
i^{-1}\alpha \cup \gamma^0(1)
\in
H^{q + c}(Z, 
i^{-1}\Omega^p_{X/S} \otimes_{i^{-1}\mathcal{O}_X}
\mathcal{H}^c_Z(\Omega^c_{X/S}))
$$
使用上同调，备注 \ref{cohomology-remark-support-cup-product}
以及楔积，得到典范映射
$$
i^{-1}\Omega^p_{X/S} \otimes_{i^{-1}\mathcal{O}_X}^\mathbf{L}
R\mathcal{H}_Z(\Omega^c_{X/S}) \to
R\mathcal{H}_Z(\Omega^p_{X/S} \otimes_{\mathcal{O}_X}^\mathbf{L}
\Omega^c_{X/S}) \to
R\mathcal{H}_Z(\Omega^{p + c}_{X/S})
$$
由概形的导出范畴，引理
\ref{perfect-lemma-supported-trivial-vanishing}，对象
$R\mathcal{H}_Z(\Omega^j_{X/S})$ 的上同调层在次数 $> c$ 时消失。
因此在次数 $c$ 的上同调层上得到映射
$$
i^{-1}\Omega^p_{X/S} \otimes_{i^{-1}\mathcal{O}_X}
\mathcal{H}^c_Z(\Omega^c_{X/S}) \longrightarrow
\mathcal{H}^c_Z(\Omega^{p + c}_{X/S})
$$
由上同调的函子性，表达式 $i^{-1}\alpha \wedge \gamma^0(1)$
是杯积 $i^{-1}\alpha \cup \gamma^0(1)$ 的像。

\medskip\noindent
按这种方式解释公式的含义后，由杯积的一般性质
（上同调，第 \ref{cohomology-section-cup-product} 节），
现在只需证明图
$$
\xymatrix{
i^{-1}\Omega^p_X \otimes \Omega^0_Z \ar[rr]_{\text{id} \otimes \gamma^0}
\ar[d] & &
i^{-1}\Omega^p_X \otimes \mathcal{H}^c_Z(\Omega^c_X) \ar[d]^\wedge \\
\Omega^p_Z \otimes \Omega^0_Z \ar[r]^\wedge &
\Omega^p_Z \ar[r]^{\gamma^p} &
\mathcal{H}^c_Z(\Omega^{p + c}_X)
}
$$
在 $Z$ 上层的范畴中交换（这里显然略有滥用记号）。
这归结为对映射 $\gamma^j_{f_1, \ldots, f_c}$ 的简单计算，我们将其略去；
事实上，选择这些映射正是为了使上述交换性成立，并使 $1$ 映到
$\frac{\text{d}f_1 \wedge \ldots \wedge \text{d}f_c}{f_1 \ldots f_c}$.
\end{proof}

\begin{lemma}
\label{derham-lemma-gysin-transverse}
设 $c \geq 0$ 是整数，并设
$$
\xymatrix{
Z' \ar[d]_h \ar[r] & X' \ar[d]_g \ar[r] & S' \ar[d] \\
Z \ar[r] & X \ar[r] & S
}
$$
为概形的交换图。假设：
\begin{enumerate}
\item $Z \to X$ and $Z' \to X'$
满足引理 \ref{derham-lemma-gysin-global} 的假设；
\item 图中的左方块是笛卡尔方块；
\item $h^*\mathcal{C}_{Z/X} \to \mathcal{C}_{Z'/X'}$
（态射，引理 \ref{morphisms-lemma-conormal-functorial}）是同构。
\end{enumerate}
那么图
$$
\xymatrix{
h^*\Omega^p_{Z/S} \ar[rr]_-{h^{-1}\gamma^p} \ar[d] & &
\mathcal{O}_{X'}|_{Z'} \otimes_{h^{-1}\mathcal{O}_X|_Z}
h^{-1}\mathcal{H}^c_Z(\Omega^{p + c}_{X/S}) \ar[d] \\
\Omega^p_{Z'/S'} \ar[rr]^{\gamma^p} & &
\mathcal{H}^c_{Z'}(\Omega^{p + c}_{X'/S'})
}
$$
交换。左竖直箭头来自微分模的函子性，右竖直箭头使用上同调，备注
\ref{cohomology-remark-support-functorial}。
\end{lemma}

\begin{proof}
更确切地说，考虑复合
\begin{align*}
\mathcal{O}_{X'}|_{Z'} \otimes_{h^{-1}\mathcal{O}_X|_Z}^\mathbf{L}
h^{-1}R\mathcal{H}_Z(\Omega^{p + c}_{X/S})
& \to
R\mathcal{H}_{Z'}(Lg^*\Omega^{p + c}_{X/S}) \\
& \to
R\mathcal{H}_{Z'}(g^*\Omega^{p + c}_{X/S}) \\
& \to
R\mathcal{H}_{Z'}(\Omega^{p + c}_{X'/S'})
\end{align*}
其中第一个箭头由上同调，备注
\ref{cohomology-remark-support-functorial} 给出，最后一个箭头由微分的函子性给出。
由概形的导出范畴，引理
\ref{perfect-lemma-supported-trivial-vanishing}，上同调层在次数 $> c$
时消失，因此该复合诱导右竖直箭头。交换性可以局部检验。
因此可以假设 $Z$ 由
$f_1, \ldots, f_c \in \Gamma(X, \mathcal{O}_X)$ 截出。
那么 $Z'$ 由 $f'_i = g^\sharp(f_i)$ 截出。映射
$c_{f_1, \ldots, f_c}$ 和 $c_{f'_1, \ldots, f'_c}$
嵌入交换图
$$
\xymatrix{
h^*i^*\Omega^p_{X/S} \ar[rr]_-{h^{-1}c_{f_1, \ldots, f_c}} \ar[d] & &
\mathcal{O}_{X'}|_{Z'} \otimes_{h^{-1}\mathcal{O}_X|_Z}
h^{-1}\mathcal{H}^c_Z(\Omega^p_{X/S}) \ar[d] \\
(i')^*\Omega^p_{X'/S'} \ar[rr]^{c_{f'_1, \ldots, f'_c}} & &
\mathcal{H}^c_{Z'}(\Omega^p_{X'/S'})
}
$$
见概形的导出范畴，备注
\ref{perfect-remark-supported-functorial}。回忆，给定一个
$p$-形式 $\omega$（在 $Z$ 上），我们如下定义 $\gamma^p(\omega)$：
局部地在 $X$ 和 $Z$ 上，选择一个 $p$-形式 $\tilde \omega$
（在 $X$ 上）来提升 $\omega$，并取
$\gamma^p(\omega) =
c_{f_1, \ldots, f_c}(\tilde \omega) \wedge
\text{d}f_1 \wedge \ldots \wedge \text{d}f_c$.
由于形式
$\text{d}f_1 \wedge \ldots \wedge \text{d}f_c$ 拉回到
$\text{d}f'_1 \wedge \ldots \wedge \text{d}f'_c$，结论成立。
\end{proof}

\begin{remark}
\label{derham-remark-how-to-use}
设 $X \to S$、$i : Z \to X$ 和 $c \geq 0$ 如引理
\ref{derham-lemma-gysin-global} 所述。设 $p \geq 0$，并假设
$\mathcal{H}^i_Z(\Omega^{p + c}_{X/S}) = 0$ 对
$i = 0, \ldots, c - 1$ 成立。如果 $X \to S$ 光滑且
$Z \to X$ 是 Koszul 正则浸入，则此消失成立；见概形的导出范畴，引理
\ref{perfect-lemma-supported-vanishing}。于是得到映射
$$
\gamma^{p, q} :
H^q(Z, \Omega^p_{Z/S})
\longrightarrow
H^{q + c}(X, \Omega^{p + c}_{X/S})
$$
：先使用
$\gamma^p : \Omega^p_{Z/S} \to \mathcal{H}^c_Z(\Omega^{p + c}_{X/S})$
映入
$$
H^q(Z, \mathcal{H}^c_Z(\Omega^{p + c}_{X/S})) =
H^q(Z, R\mathcal{H}_Z(\Omega^{p + c}_{X/S})[c]) =
H^q(X, i_*R\mathcal{H}_Z(\Omega^{p + c}_{X/S})[c])
$$
，再使用伴随映射
$i_*R\mathcal{H}_Z(\Omega^{p + c}_{X/S}) \to \Omega^{p + c}_{X/S}$
继续映到所需的霍奇上同调模。
\end{remark}

\begin{lemma}
\label{derham-lemma-gysin-differential-hodge}
设 $X \to S$ 和 $i : Z \to X$ 如引理 \ref{derham-lemma-gysin-global} 所述。
假设 $X \to S$ 光滑且 $Z \to X$ 为 Koszul 正则浸入。
Gysin 映射 $\gamma^{p, q}$ 与 $\Omega^\bullet_{X/S}$ 和
$\Omega^\bullet_{Z/S}$ 上的德拉姆微分相容。
\end{lemma}

\begin{proof}
这立即由引理 \ref{derham-lemma-gysin-differential-global} 得出。
\end{proof}

\begin{lemma}
\label{derham-lemma-gysin-projection-global}
设 $X \to S$、$i : Z \to X$ 和 $c \geq 0$ 如引理
\ref{derham-lemma-gysin-global} 所述。假设 $X \to S$ 光滑且
$Z \to X$ 为 Koszul 正则浸入。给定
$\alpha \in H^q(X, \Omega^p_{X/S})$，有
$\gamma^{p, q}(\alpha|_Z) = \alpha \cup \gamma^{0, 0}(1)$，这是在
$H^{q + c}(X, \Omega^{p + c}_{X/S})$ 中成立，其中
$\gamma^{a, b}$ 如备注 \ref{derham-remark-how-to-use} 所述。
\end{lemma}

\begin{proof}
本引理由引理 \ref{derham-lemma-gysin-projection} 和上同调，引理
\ref{cohomology-lemma-support-cup-product} 得出。
建议读者跳过下面更详细的讨论。

\medskip\noindent
下文将不加说明地使用
$R\mathcal{H}_Z(\Omega^j_{X/S}) = \mathcal{H}^c_Z(\Omega^j_{X/S})[-c]$
对所有 $j$ 成立，正如备注 \ref{derham-remark-how-to-use} 所指出的。
还将不加说明地使用等同
$H^{q + c}_Z(X, \Omega^j_{X/S}) = H^{q + c}(Z, R\mathcal{H}_Z(\Omega^j_{X/S}) =
H^q(Z, \mathcal{H}^c_Z(\Omega^j_{X/S}))$；第一个等同见上同调，引理
\ref{cohomology-lemma-local-to-global-sections-with-support}。
在这些等同下：
\begin{enumerate}
\item $\gamma^0(1) \in H^c_Z(X, \Omega^c_{X/S})$ 映到 $\gamma^{0, 0}(1)$，
后者位于 $H^c(X, \Omega^c_{X/S})$ 中；
\item 等式的右边 $i^{-1}\alpha \wedge \gamma^0(1)$（在引理
\ref{derham-lemma-gysin-projection} 中）是杯积在楔积下的像；该杯积来自
上同调，备注 \ref{cohomology-remark-support-cup-product-global}，
作用于元素 $\alpha$ 和 $\gamma^0(1)$。换言之，引理
\ref{derham-lemma-gysin-projection} 的证明中的构造与上同调，备注
\ref{cohomology-remark-support-cup-product-global} 中的构造相符；
\item 由上同调，引理 \ref{cohomology-lemma-support-cup-product}，
它映到 $\alpha \cup \gamma^{0, 0}(1)$（在
$H^{q + c}(X, \Omega^p_{X/S} \otimes \Omega^c_{X/S})$ 中）；
\item 等式的左边 $\gamma^p(\alpha|_Z)$（在引理
\ref{derham-lemma-gysin-projection} 中）映到
$\gamma^{p, q}(\alpha|_Z)$。
\end{enumerate}
证明完毕。
\end{proof}

\begin{lemma}
\label{derham-lemma-gysin-transverse-global}
设 $c \geq 0$，并设
$$
\xymatrix{
Z' \ar[d]_h \ar[r] & X' \ar[d]_g \ar[r] & S' \ar[d] \\
Z \ar[r] & X \ar[r] & S
}
$$
满足引理 \ref{derham-lemma-gysin-transverse} 的假设；此外还假设
$X \to S$ 和 $X' \to S'$ 光滑，且 $Z \to X$ 和 $Z' \to X'$
是 Koszul 正则浸入。那么图
$$
\xymatrix{
H^q(Z, \Omega^p_{Z/S}) \ar[rr]_-{\gamma^{p, q}} \ar[d] & &
H^{q + c}(X, \Omega^{p + c}_{X/S}) \ar[d] \\
H^q(Z', \Omega^p_{Z'/S'}) \ar[rr]^{\gamma^{p, q}} & &
H^{q + c}(X', \Omega^{p + c}_{X'/S'})
}
$$
交换，其中 $\gamma^{p, q}$ 如备注 \ref{derham-remark-how-to-use} 所述。
\end{lemma}

\begin{proof}
这由引理 \ref{derham-lemma-gysin-transverse} 与上同调，引理
\ref{cohomology-lemma-support-functorial} 结合得出。
\end{proof}

\begin{lemma}
\label{derham-lemma-class-of-a-point}
设 $k$ 是一个域。设 $X$ 是 $k$ 上维数为 $d$ 的不可约光滑固有概形。
设 $Z \subset X$ 是仅由一个 $k$-有理点 $x$ 组成的既约闭子概形。那么
$1 \in k = H^0(Z, \mathcal{O}_Z) = H^0(Z, \Omega^0_{Z/k})$
在备注 \ref{derham-remark-how-to-use} 的映射
$H^0(Z, \Omega^0_{Z/k}) \to H^d(X, \Omega^d_{X/k})$ 下的像非零。
\end{lemma}

\begin{proof}
映射 $\gamma^0 : \mathcal{O}_Z \to
\mathcal{H}^d_Z(\Omega^d_{X/k}) = R\mathcal{H}_Z(\Omega^d_{X/k})[d]$
伴随于映射
$$
g^0 : i_*\mathcal{O}_Z \longrightarrow \Omega^d_{X/k}[d]
$$
（在 $D(\mathcal{O}_X)$ 中）。回忆，
$\Omega^d_{X/k} = \omega_X$ 是 $X/k$ 的对偶层；见概形对偶，引理
\ref{duality-lemma-duality-proper-over-field}。因此，本引理陈述中映射的
$k$-线性对偶是映射
$$
H^0(X, \mathcal{O}_X) \to \Ext^d_X(i_*\mathcal{O}_Z, \omega_X)
$$
，它把 $1$ 映到 $g^0$。因此只需证明 $g^0$ 非零。
这可以在任意邻域 $U$ 中完成，该邻域包含点 $x$。选择 $U$，使得存在
$f_1, \ldots, f_d \in \mathcal{O}_X(U)$，它们只在 $x$ 处消失，
并生成极大理想 $\mathfrak m_x \subset \mathcal{O}_{X, x}$。
可以假设 $U = \Spec(R)$ 是仿射的。回顾 $\gamma^0$ 的构造，
可见我们的扩张由下式给出：
$$
k \to
(R \to \bigoplus\nolimits_{i_0} R_{f_{i_0}} \to
\bigoplus\nolimits_{i_0 < i_1} R_{f_{i_0}f_{i_1}} \to
\ldots \to R_{f_1\ldots f_r})[d] \to R[d]
$$
其中 $1$ 在第一个映射下映到 $1/f_1 \ldots f_c$。
它非零，因为在此情形下 $1/f_1 \ldots f_c$ 是局部上同调群
$H^d_{(f_1, \ldots, f_d)}(R)$ 的非零元素，
\end{proof}






\section{相对庞加莱对偶}
\label{derham-section-relative-poincare-duality}

\noindent
本节证明基底上光滑固有概形的相对德拉姆上同调所满足的庞加莱对偶。
我们强烈建议读者先阅读第 \ref{derham-section-poincare-duality} 节。

\begin{situation}
\label{derham-situation-relative-duality}
这里 $S$ 是拟紧拟分离概形，而 $f : X \to S$ 是光滑固有的概形态射，
其所有纤维均非空且等维，维数为 $n$。
\end{situation}

\begin{lemma}
\label{derham-lemma-relative-bottom-part-degenerates}
在情形 \ref{derham-situation-relative-duality} 中，正向像
$f_*\mathcal{O}_X$ 是有限平展 $\mathcal{O}_S$-代数；并且在 $S$ 上局部地，
$Rf_*\mathcal{O}_X = f_*\mathcal{O}_X \oplus P$ 在 $D(\mathcal{O}_S)$ 中成立，
其中 $P$ 是 Tor 振幅落在 $[1, \infty)$ 中的完美对象。映射
$\text{d} : f_*\mathcal{O}_X \to f_*\Omega_{X/S}$ 为零。
\end{lemma}

\begin{proof}
陈述的第一部分由概形的导出范畴，引理
\ref{perfect-lemma-proper-flat-geom-red} 得出。置
$S' = \underline{\Spec}_S(f_*\mathcal{O}_X)$，便得到分解
$X \to S' \to S$（这是 Stein 分解；见态射进阶，第
\ref{more-morphisms-section-stein-factorization} 节，不过这里并不需要这一点）。
例如由态射，引理 \ref{morphisms-lemma-triangle-differentials} 和
\ref{morphisms-lemma-etale-at-point} 可见
$\Omega_{X/S} = \Omega_{X/S'}$。这当然蕴含
$\text{d} : f_*\mathcal{O}_X \to f_*\Omega_{X/S}$ 为零。
\end{proof}

\begin{lemma}
\label{derham-lemma-relative-duality-hodge}
在情形 \ref{derham-situation-relative-duality} 中，存在 $\mathcal{O}_S$-模映射
$$
t : Rf_*\Omega^n_{X/S}[n] \longrightarrow \mathcal{O}_S
$$
它在预复合以 $H^0(X, \mathcal{O}_X)$ 中某个单位的乘法这个意义下唯一，
并具有如下性质：对每个 $p$，配对
$$
Rf_*\Omega^p_{X/S}
\otimes_{\mathcal{O}_S}^\mathbf{L}
Rf_*\Omega^{n - p}_{X/S}[n]
\longrightarrow
\mathcal{O}_S
$$
由相对杯积再复合 $t$ 给出，并且是 $S$ 上完美复形之间的完美配对。
\end{lemma}

\begin{proof}
首先假设 $S$ 是诺特概形。令 $\omega^\bullet_{X/S}$ 为概形的对偶性，
备注 \ref{duality-remark-relative-dualizing-complex} 所述的
$X$ 在 $S$ 上的相对对偶复形，并令
$Rf_*\omega_{X/S}^\bullet \to \mathcal{O}_S$ 为其迹映射。由概形的对偶性，
引理 \ref{duality-lemma-smooth-proper}（此处使用了 $S$ 的诺特性），存在同构
$\omega^\bullet_{X/S} \cong \Omega^n_{X/S}[n]$；利用此同构便得到 $t$。
由引理 \ref{derham-lemma-proper-smooth-de-Rham}，复形 $Rf_*\Omega^p_{X/S}$ 是完美的。
由于 $\Omega^p_{X/S}$ 局部自由，并且
$\Omega^p_{X/S} \otimes_{\mathcal{O}_X} \Omega^{n - p}_{X/S} \to
\Omega^n_{X/S}$ 给出同构 $\Omega^p_{X/S} \cong
\SheafHom_{\mathcal{O}_X}(\Omega^{n - p}_{X/S}, \Omega^n_{X/S})$，
由概形的对偶性，备注
\ref{duality-remark-relative-dualizing-complex-relative-cup-product} 可知，
相对杯积诱导的配对是完美的。

\medskip\noindent
一般情形下存在性的证明。
由绝对诺特逼近，可以找到笛卡尔图
$$
\xymatrix{
X \ar[r] \ar[d]^f & X_0 \ar[d]^{f_0} \\
S \ar[r]^g & S_0
}
$$
其中 $S_0$ 是诺特概形，$f_0$ 是光滑固有态射，且其所有纤维均非空并且
等维，维数为 $n$。见极限，引理 \ref{limits-lemma-descend-finite-presentation}、
\ref{limits-lemma-descend-smooth}、\ref{limits-lemma-descend-dimension-d}、
\ref{limits-lemma-descend-surjective} 和
\ref{limits-lemma-eventually-proper}。选择陈述中那样的映射
$$
t_0 : Rf_{0, *}\Omega^n_{X_0/S_0}[n] \longrightarrow \mathcal{O}_{S_0}
$$
。由上同调与基变换，复形 $Rf_*(\Omega^p_{X/S})$ 等于
$Lg^*(Rf_{0, *}\Omega^p_{X_0/S_0})$，这对每个 $p \geq 0$ 均成立；见引理
\ref{derham-lemma-proper-smooth-de-Rham}。特别地，$t_0$ 沿 $g$ 的拉回给出
本引理陈述中的映射 $t$。具体地，由上同调，引理
\ref{cohomology-lemma-base-change-cup-product}，杯积映射
$$
Rf_*\Omega^p_{X/S}
\otimes_{\mathcal{O}_S}^\mathbf{L}
Rf_*\Omega^{n - p}_{X/S}[n]
\longrightarrow
Rf_*\Omega^n_{X/S}
$$
是沿 $g$ 拉回 $f_0$ 的相应杯积映射。因此，$t_0$ 给出完美配对这一事实
拉回后说明 $t$ 在 $S$ 上具有同一性质。

\medskip\noindent
映射 $t$ 的唯一性。选择可区别三角
$f_*\mathcal{O}_X \to Rf_*\mathcal{O}_X \to P \to f_*\mathcal{O}_X[1]$.
。由引理 \ref{derham-lemma-relative-bottom-part-degenerates}，对象 $P$ 是
Tor 振幅落在 $[1, \infty)$ 中的完美对象，并且该三角在 $S$ 上局部分裂。
因此 $R\SheafHom_{\mathcal{O}_X}(P, \mathcal{O}_X)$ 是 Tor 振幅落在
$(-\infty, -1]$ 中的完美对象。于是上述对偶性表明，在 $S$ 上局部地有
$$
Rf_*\Omega^n_{X/S}[n] \cong
R\SheafHom_{\mathcal{O}_S}(f_*\mathcal{O}_X, \mathcal{O}_S)
\oplus R\SheafHom_{\mathcal{O}_X}(P, \mathcal{O}_X)
$$
这说明 $R^nf_*\Omega^n_{X/S}$ 是有限局部自由的，并且得到完美的
$\mathcal{O}_S$-双线性配对
$$
f_*\mathcal{O}_X \times R^nf_*\Omega^n_{X/S} \longrightarrow \mathcal{O}_S
$$
；这里使用了 $t$。这蕴含任何 $\mathcal{O}_S$-线性映射
$t' : R^nf_*\Omega^n_{X/S} \to \mathcal{O}_S$ 都具有形式
$t' = t \circ g$，其中某个
$g \in \Gamma(S, f_*\mathcal{O}_X) = \Gamma(X, \mathcal{O}_X)$。
要使 $t'$ 仍确定完美配对，$g$ 必须是单位。证明完毕。
\end{proof}

\begin{lemma}
\label{derham-lemma-relative-top-part-degenerates}
在情形 \ref{derham-situation-relative-duality} 中，映射
$\text{d} : R^nf_*\Omega^{n - 1}_{X/S} \to R^nf_*\Omega^n_{X/S}$
为零。
\end{lemma}

\noindent
正如我们在引理 \ref{derham-lemma-top-part-degenerates} 的证明中提到的，
本引理并非引理
\ref{derham-lemma-relative-duality-hodge} 和
\ref{derham-lemma-relative-bottom-part-degenerates} 的简单推论。

\begin{proof}[$S$ 既约时的证明]
假设 $S$ 既约。注意
$\text{d} : R^nf_*\Omega^{n - 1}_{X/S} \to R^nf_*\Omega^n_{X/S}$
是（拟凝聚）$\mathcal{O}_S$-模之间的 $\mathcal{O}_S$-线性映射。
$\mathcal{O}_S$-模 $R^nf_*\Omega^n_{X/S}$ 是有限局部自由的
（由引理
\ref{derham-lemma-relative-duality-hodge} 和
\ref{derham-lemma-relative-bottom-part-degenerates}，它是有限局部自由
$\mathcal{O}_S$-模 $f_*\mathcal{O}_X$ 的对偶）。由于 $S$ 既约，
只需证明 $\text{d}$ 在每个一般点 $\eta \in S$ 处的茎为零；
这可通过考察仿射开集上的截面、利用 $\text{d}$ 的目标局部自由以及代数，
引理 \ref{algebra-lemma-reduced-ring-sub-product-fields} 的第 (2) 部分得出。
由于 $S$ 既约，有 $\mathcal{O}_{S, \eta} = \kappa(\eta)$；见代数，引理
\ref{algebra-lemma-minimal-prime-reduced-ring}。因此 $\text{d}_\eta$ 等同于映射
$$
\text{d} :
H^n(X_\eta, \Omega^{n - 1}_{X_\eta/\kappa(\eta)})
\longrightarrow
H^n(X_\eta, \Omega^n_{X_\eta/\kappa(\eta)})
$$
，它由引理 \ref{derham-lemma-top-part-degenerates} 为零。
\end{proof}

\begin{proof}[一般情形下的证明]
注意该问题在 $S$ 上对平坦拓扑是局部的：若 $S' \to S$ 是满的平坦概形态射，
并且该映射拉回到 $S'$ 后为零，则原映射为零。此外，由平坦基变换（概形的上同调，
引理 \ref{coherent-lemma-flat-base-change-cohomology}），该映射的形成与沿平坦态射的
基变换交换。

\medskip\noindent
考虑态射进阶，定理
\ref{more-morphisms-theorem-stein-factorization-general} 所述的 Stein 分解
$X \to S' \to S$。由引理 \ref{derham-lemma-relative-bottom-part-degenerates}，态射
$\pi : S' \to S$ 有限平展。态射 $f : X \to S'$ 固有（由该定理），
光滑（由态射进阶，引理
\ref{more-morphisms-lemma-smooth-etale-permanence}），并且由 Stein 分解定理，
其纤维几何连通。在引理 \ref{derham-lemma-relative-bottom-part-degenerates} 的证明中，
我们已看到 $\Omega_{X/S} = \Omega_{X/S'}$，因为 $S' \to S$ 平展。
因此 $\Omega^\bullet_{X/S} = \Omega^\bullet_{X/S'}$。我们有
$$
R^qf_*\Omega^p_{X/S} = \pi_*R^qf'_*\Omega^p_{X/S'}
$$
，对所有 $p, q$ 都成立。这由 Leray 谱序列（上同调，引理
\ref{cohomology-lemma-relative-Leray}）、$\pi$ 有限因而仿射，以及概形的上同调，
引理 \ref{coherent-lemma-relative-affine-vanishing} 得出（当然还使用了
$R^qf'_*\Omega^p_{X'/S}$ 拟凝聚）。因此，本引理中的映射是将 $\pi_*$ 作用于
$\text{d} : R^nf'_*\Omega^{n - 1}_{X/S'} \to R^nf'_*\Omega^n_{X/S'}$ 所得的映射。
换言之，为证明本引理，可以把 $f : X \to S$ 替换为 $f' : X \to S'$，
从而归约到下一段所讨论的情形。

\medskip\noindent
假设 $f$ 的纤维几何连通，并且 $f_*\mathcal{O}_X = \mathcal{O}_S$。
对每个 $s \in S$，可以选择平展邻域 $(S', s') \to (S, s)$，
使得该基变换 $X' \to S'$ 作为 $S$ 的基变换具有截面；见态射进阶，引理
\ref{more-morphisms-lemma-etale-nbhd-dominates-smooth}。
由证明开头的说明，这又归约到下一段所讨论的情形。

\medskip\noindent
假设 $f$ 的纤维几何连通，$f_*\mathcal{O}_X = \mathcal{O}_S$，并且有截面
$s : S \to X$，它是 $f$ 的截面。我们可以并且将会假设 $S = \Spec(A)$ 仿射。映射
$s^* : R\Gamma(X, \mathcal{O}_X) \to R\Gamma(S, \mathcal{O}_S) = A$
是映射 $A \to R\Gamma(X, \mathcal{O}_X)$ 的分裂。因此可写成
$$
R\Gamma(X, \mathcal{O}_X) = A \oplus P
$$
，其中 $P$ 是 $s^*$ 的“核”。由引理
\ref{derham-lemma-relative-bottom-part-degenerates}，对象 $P$（属于 $D(A)$）是
Tor 振幅落在 $[1, n]$ 中的完美对象。与引理
\ref{derham-lemma-relative-duality-hodge} 的证明相同，可见
$H^n(X, \Omega^n_{X/S})$ 是局部自由 $A$-模，其秩为 $1$
（事实上它与 $A$ 对偶，因而是秩为 $1$ 的自由模；我们很快会选择一个生成元，
但既不想核对它是否为同一个生成元，也无需如此）。

\medskip\noindent
用 $Z \subset X$ 表示 $s$ 的像；由概形，引理
\ref{schemes-lemma-section-immersion}，它是 $X$ 的闭子概形。注意，由除子，引理
\ref{divisors-lemma-regular-quasi-regular-immersion}，$Z \to X$ 是正则闭浸入
（从而更是 Koszul 正则闭浸入）；这由除子，引理
\ref{divisors-lemma-section-smooth-regular-immersion} 得出。当然，
$Z \to X$ 的余维数为 $n$。因此，由备注 \ref{derham-remark-how-to-use}，
可以考虑映射
$$
\gamma^{0, 0} : H^0(Z, \Omega^0_{Z/S}) \longrightarrow H^n(X, \Omega^n_{X/S})
$$
，并置 $\xi = \gamma^{0, 0}(1) \in H^n(X, \Omega^n_{X/S})$。

\medskip\noindent
我们断言 $\xi$ 是一个基元素。事实上，由于最高次数满足基变换（例如见极限，
引理 \ref{limits-lemma-higher-direct-images-zero-above-dimension-fibre}），可见
$H^n(X, \Omega^n_{X/S}) \otimes_A k = H^n(X_k, \Omega^n_{X_k/k})$
，这对任意环映射 $A \to k$ 均成立。由 $\xi$ 的构造与基变换相容（见引理
\ref{derham-lemma-gysin-transverse-global}），可见 $\xi$ 在
$H^n(X_k, \Omega^n_{X_k/k})$ 中的像由引理 \ref{derham-lemma-class-of-a-point} 非零，
只要 $k$ 是域。
因此 $\xi$ 是可逆模的处处不消失截面，从而是生成元。

\medskip\noindent
取 $\theta \in H^n(X, \Omega^{n - 1}_{X/S})$。我们必须证明
$\text{d}(\theta)$ 在 $H^n(X, \Omega^n_{X/S})$ 中为零。
可以写成 $\text{d}(\theta) = a \xi$，其中某个 $a \in A$，
因为 $\xi$ 是基元素。于是必须证明 $a = 0$。

\medskip\noindent
考虑闭浸入
$$
\Delta : X \to X \times_S X
$$
。它也是某个光滑态射（即任一投影）的截面，因而也是余维数为 $n$ 的正则且
Koszul 正则浸入。因此可以考虑备注 \ref{derham-remark-how-to-use} 中的映射
$$
\gamma^{p, q} :
H^q(X, \Omega^p_{X/S})
\longrightarrow
H^{q + n}(X \times_S X, \Omega^{p + n}_{X \times_S X/S})
$$
。考虑像
$$
\gamma^{n - 1, n}(\theta) \in
H^{2n}(X \times_S X, \Omega^{2n - 1}_{X \times_S X})
$$
。由引理 \ref{derham-lemma-de-rham-complex-product} 有
$$
\Omega^{2n - 1}_{X \times_S X} =
\Omega^{n - 1}_{X/S} \boxtimes \Omega^n_{X/S} \oplus
\Omega^n_{X/S} \boxtimes \Omega^{n - 1}_{X/S}
$$
由 K\"unneth 公式（可用概形的导出范畴，引理 \ref{perfect-lemma-kunneth}，
或概形的导出范畴，引理 \ref{perfect-lemma-kunneth-single-sheaf}）可见
$$
H^{2n}(X \times_S X, \Omega^{n - 1}_{X/S} \boxtimes \Omega^n_{X/S}) =
H^n(X, \Omega^{n - 1}_{X/S}) \otimes_A H^n(X, \Omega^n_{X/S})
$$
以及
$$
H^{2n}(X \times_S X, \Omega^n_{X/S} \boxtimes \Omega^{n - 1}_{X/S}) =
H^n(X, \Omega^n_{X/S}) \otimes_A H^n(X, \Omega^{n - 1}_{X/S})
$$
。事实上，由于我们考察的是最高次数，不会有更高的 Tor 群介入。再结合
$\xi$ 是生成元这一事实，可以写成
$$
\gamma^{n - 1, n}(\theta) = \theta_1 \otimes \xi + \xi \otimes \theta_2
$$
，其中 $\theta_1, \theta_2 \in H^n(X, \Omega^{n - 1}_{X/S})$。
完全相同的论证给出
$$
\gamma^{n, n}(\xi) = b \xi \otimes \xi
$$
，它位于
$H^{2n}(X \times_S X, \Omega^{2n}_{X \times_S X/S}) =
H^n(X, \Omega^n_{X/S}) \otimes_A H^n(X, \Omega^n_{X/S})$
中，其中某个 $b \in H^0(S, \mathcal{O}_S)$。

\medskip\noindent
{\bf 断言：} $\theta_1 = \theta$、$\theta_2 = \theta$ 且 $b = 1$。
先说明该断言蕴含所需结论 $a = 0$。事实上，由引理
\ref{derham-lemma-gysin-differential-hodge} 有
$$
\gamma^{n, n}(\text{d}(\theta)) = \text{d}(\gamma^{n - 1, n}(\theta))
$$
由上面的选择，这给出
$$
a \xi \otimes \xi =
\gamma^{n, n}(a\xi) =
\text{d}(\theta \otimes \xi + \xi \otimes \theta) =
a \xi \otimes \xi + (-1)^n a \xi \otimes \xi
$$
最右边的等式源于如下事实：映射
$\text{d} : \Omega^{2n - 1}_{X \otimes_S X/S} \to \Omega^{2n}_{X \times_S X/S}$
由引理 \ref{derham-lemma-de-rham-complex-product} 是微分
$\text{d} \boxtimes 1 : \Omega^{n - 1}_{X/S} \boxtimes \Omega^n_{X/S}
\to \Omega^n_{X/S} \boxtimes \Omega^n_{X/S}$
与微分
$(-1)^n 1 \boxtimes \text{d} : \Omega^n_{X/S} \boxtimes \Omega^{n - 1}_{X/S}
\to \Omega^n_{X/S} \boxtimes \Omega^n_{X/S}$ 之和。更多说明见第
\ref{derham-section-kunneth} 节和概形的导出范畴，第
\ref{perfect-section-kunneth-complexes} 节。由于 $\xi \otimes \xi$ 是秩为
$1$ 的自由 $A$-模
$H^n(X, \Omega^n_{X/S}) \otimes_A H^n(X, \Omega^n_{X/S})$
的一组基，故有
$$
a = a + (-1)^n a \Rightarrow a = 0
$$
，正是所需结论。

\medskip\noindent
在证明的其余部分，我们证明上述断言。记
$\eta = \gamma^{0, 0}(1) \in H^n(X \times_S X, \Omega^n_{X \times_S X/S})$。
由于 $\Omega^n_{X \times_S X/S} =
\bigoplus_{p + p' = n} \Omega^p_{X/S} \boxtimes \Omega^{p'}_{X/S}$，可写成
$$
\eta = \eta_0 + \eta_1 + \ldots + \eta_n
$$
，其中 $\eta_p$ 属于
$H^n(X \times_S X, \Omega^p_{X/S} \boxtimes \Omega^{n - p}_{X/S})$.
当 $p = 0$ 时，可写成
\begin{align*}
H^n(X \times_S X, \mathcal{O}_X \boxtimes \Omega^n_{X/S})
& =
H^n(R\Gamma(X, \mathcal{O}_X) \otimes_A^\mathbf{L}
R\Gamma(X, \Omega^n_{X/S})) \\
& =
A \otimes_A H^n(X, \Omega^n_{X/S}) \oplus
H^n(P \otimes_A^\mathbf{L} R\Gamma(X, \Omega^n_{X/S}))
\end{align*}
；这里使用了先前给出的分解 $R\Gamma(X, \mathcal{O}_X) = A \oplus P$。
考虑态射 $(s, \text{id}) : X \to X \times_S X$。于是概形论意义下有
$(s, \text{id})^{-1}(\Delta) = Z$。因此，由引理
\ref{derham-lemma-gysin-transverse-global} 可见 $(s, \text{id})^*\eta = \xi$。这意味着
$$
\xi = (s, \text{id})^*\eta = (s^* \otimes \text{id})(\eta_0)
$$
这恰好意味着 $\eta_0$ 在上述直和分解中的第一个分量是 $\xi$。换言之，可以写成
$$
\eta_0 = 1 \otimes \xi + \eta'_0
$$
，其中 $\eta'_0 \in H^n(P \otimes_A^\mathbf{L} R\Gamma(X, \Omega^n_{X/S}))$。
完全相同地，当 $p = n$ 时，可以写成
\begin{align*}
H^n(X \times_S X, \Omega^n_{X/S} \boxtimes \mathcal{O}_X)
& =
H^n(R\Gamma(X, \Omega^n_{X/S}) \otimes_A^\mathbf{L}
R\Gamma(X, \mathcal{O}_X)) \\
& =
H^n(X, \Omega^n_{X/S}) \otimes_A A \oplus
H^n(R\Gamma(X, \Omega^n_{X/S}) \otimes_A^\mathbf{L} P)
\end{align*}
并且可以写成
$$
\eta_n = \xi \otimes 1 + \eta'_n
$$
，其中 $\eta'_n \in H^n(R\Gamma(X, \Omega^n_{X/S}) \otimes_A^\mathbf{L} P)$。

\medskip\noindent
注意，$\text{pr}_1^*\theta = \theta \otimes 1$ 和
$\text{pr}_2^*\theta = 1 \otimes \theta$ 是 $X \times_S X$ 上的霍奇上同调类，
它们拉回为 $\theta$，这里的拉回沿 $\Delta$ 进行。因此，由引理
\ref{derham-lemma-gysin-projection-global} 有
$$
\theta_1 \otimes \xi + \xi \otimes \theta_2 =
\gamma^{n - 1, n}(\theta) =
(\theta \otimes 1) \cup \eta =
(1 \otimes \theta) \cup \eta
$$
，这是在 $X \times_S X$ 相对于 $S$ 的霍奇上同调环中成立。按照
$X \times_S X/S$ 的微分模上的直和分解，得到
$$
\theta_1 \otimes \xi =
(\theta \otimes 1) \cup \eta_0
\quad\text{且}\quad
\xi \otimes \theta_2 =
(1 \otimes \theta) \cup \eta_n
$$
考察上面所得公式 $\eta_0 = 1 \otimes \xi + \eta'_0$，可见要证明
$\theta_1 = \theta$，只需证明
$$
(\theta \otimes 1) \cup \eta'_0 = 0
$$
为此，注意与 $\theta \otimes 1$ 作杯积由如下映射在上同调上的作用给出：
$$
(P \otimes_A^\mathbf{L} R\Gamma(X, \Omega^n_{X/S}))[-n]
\xrightarrow{\theta \otimes 1}
R\Gamma(X, \Omega^{n - 1}_{X/S}) \otimes_A^\mathbf{L}
R\Gamma(X, \Omega^n_{X/S})
$$
，这是在导出范畴中；见上同调，备注
\ref{cohomology-remark-cup-with-element-map-total-cohomology}。该映射是以下两个映射的
导出张量积：
$$
\theta : P[-n] \to R\Gamma(X, \Omega^{n - 1}_{X/S})
\quad\text{且}\quad
1 : R\Gamma(X, \Omega^n_{X/S}) \to R\Gamma(X, \Omega^n_{X/S})
$$
；见概形的导出范畴，备注
\ref{perfect-remark-annoying-compatibility}。然而，其中第一个映射在 $D(A)$ 中为零，
因为它从 Tor 振幅落在 $[n + 1, 2n]$ 中的完美复形映到一个仅在次数
$0, 1, \ldots, n$ 具有上同调的复形；见代数进阶，引理
\ref{more-algebra-lemma-splitting-unique}。类似论证表明
$(1 \otimes \theta) \cup \eta'_n$ 也为零。最后，完全相同地得到
$$
b \xi \otimes \xi = \gamma^{n, n}(\xi) = (\xi \otimes 1) \cup \eta_0
$$
，并像前面一样，通过以同样方式证明
$(\xi \otimes 1) \cup \eta'_0 = 0$ 而得到结论。证明完毕。
\end{proof}

\begin{proposition}
\label{derham-proposition-relative-poincare-duality}
设 $S$ 是拟紧拟分离概形。设 $f : X \to S$ 是光滑固有的概形态射，
其所有纤维均非空且等维，维数为 $n$。存在 $\mathcal{O}_S$-模映射
$$
t : R^{2n}f_*\Omega^\bullet_{X/S} \longrightarrow \mathcal{O}_S
$$
它在预复合以 $H^0(X, \mathcal{O}_X)$ 中某个单位的乘法这个意义下唯一，
并具有如下性质：配对
$$
Rf_*\Omega^\bullet_{X/S}
\otimes_{\mathcal{O}_S}^\mathbf{L}
Rf_*\Omega^\bullet_{X/S}[2n]
\longrightarrow
\mathcal{O}_S, \quad
(\xi, \xi') \longmapsto t(\xi \cup \xi')
$$
是 $S$ 上完美复形之间的完美配对。
\end{proposition}

\begin{proof}
本证明与命题 \ref{derham-proposition-poincare-duality} 的证明完全相同。

\medskip\noindent
由相对霍奇--德拉姆谱序列
$$
E_1^{p, q} = R^qf_*\Omega^p_{X/S} \Rightarrow R^{p + q}f_*\Omega^\bullet_{X/S}
$$
（第 \ref{derham-section-hodge-to-de-rham} 节）、$\Omega^i_{X/S}$ 对 $i > n$ 的消失、
例如极限，引理 \ref{limits-lemma-higher-direct-images-zero-above-dimension-fibre}
中的消失，以及引理 \ref{derham-lemma-relative-bottom-part-degenerates} 和
\ref{derham-lemma-relative-top-part-degenerates} 的结论，可见
$R^0f_*\Omega_{X/S} = R^0f_*\mathcal{O}_X$ 且
$R^nf_*\Omega^n_{X/S} = R^{2n}f_*\Omega^\bullet_{X/S}$。
更确切地说，这些等同来自复形映射
$$
\Omega^\bullet_{X/S} \to \mathcal{O}_X[0]
\quad\text{且}\quad
\Omega^n_{X/S}[-n] \to \Omega^\bullet_{X/S}
$$
选择 $t : R^{2n}f_*\Omega_{X/S} \to \mathcal{O}_S$，使它经由此等同对应于
引理 \ref{derham-lemma-relative-duality-hodge} 中的一个 $t$。

\medskip\noindent
简记 $\Omega^\bullet = \Omega^\bullet_{X/S}$。考虑映射
(\ref{derham-equation-wedge})，它在当前情形中写作
$$
\wedge :
\text{Tot}(\Omega^\bullet \otimes_{f^{-1}\mathcal{O}_S} \Omega^\bullet)
\longrightarrow
\Omega^\bullet
$$
对每个整数 $p = 0, 1, \ldots, n$，由于次数的原因，该映射在子复形
$\text{Tot}(\sigma_{> p} \Omega^\bullet \otimes_{f^{-1}\mathcal{O}_S}
\sigma_{\geq n - p} \Omega^\bullet)$ 上为零。因此，$\wedge$ 在子复形
$\text{Tot}(\Omega^\bullet \otimes_{f^{-1}\mathcal{O}_S}
\geq_{n - p}\Omega^\bullet)$ 上的限制分解经过一个复形映射
$$
\gamma_p :
\text{Tot}(\sigma_{\leq p} \Omega^\bullet \otimes_{f^{-1}\mathcal{O}_S}
\sigma_{\geq n - p} \Omega^\bullet)
\longrightarrow
\Omega^\bullet
$$
采用第 \ref{derham-section-cup-product} 节中的相同步骤，得到相对杯积
$$
Rf_*\sigma_{\leq p} \Omega^\bullet
\otimes_{\mathcal{O}_S}^\mathbf{L}
Rf_*\sigma_{\geq n - p}\Omega^\bullet
\longrightarrow
Rf_*\Omega^\bullet
$$
我们将对 $p$ 归纳证明：这些杯积经由 $t$ 在
$Rf_*\sigma_{\leq p} \Omega^\bullet$ 与
$Rf_*\sigma_{\geq n - p}\Omega^\bullet[2n]$ 之间诱导完美配对。当 $p = n$ 时，
这正是命题的断言。

\medskip\noindent
基础情形是 $p = 0$。此时有
$$
Rf_*\sigma_{\leq p}\Omega^\bullet = Rf_*\mathcal{O}_X
\quad\text{且}\quad
Rf_*\sigma_{\geq n - p}\Omega^\bullet[2n] = Rf_*(\Omega^n[-n])[2n] =
Rf_*\Omega^n[n]
$$
此时所得恰好是引理 \ref{derham-lemma-relative-duality-hodge} 中
$Rf_*\mathcal{O}_X$ 与 $Rf_*\Omega^n[n]$ 之间的配对，故结论成立。

\medskip\noindent
归纳步骤。假设已知结论对 $p$ 成立。考虑可区别三角
$$
\Omega^{p + 1}[-p - 1] \to
\sigma_{\leq p + 1}\Omega^\bullet \to
\sigma_{\leq p}\Omega^\bullet \to
\Omega^{p + 1}[-p]
$$
以及可区别三角
$$
\sigma_{\geq n - p}\Omega^\bullet \to
\sigma_{\geq n - p - 1}\Omega^\bullet \to
\Omega^{n - p - 1}[-n + p + 1] \to
(\sigma_{\geq n - p}\Omega^\bullet)[1]
$$
注意，两者都是 $f^{-1}\mathcal{O}_S$-模层复形的同伦范畴中的可区别三角；
特别地，映射 $\sigma_{\leq p}\Omega^\bullet \to \Omega^{p + 1}[-p]$ 和
$\Omega^{n - p - 1}[-d + p + 1] \to (\sigma_{\geq n - p}\Omega^\bullet)[1]$
由实际的复形映射给出，具体使用微分 $\Omega^p \to \Omega^{p + 1}$ 和微分
$\Omega^{n - p - 1} \to \Omega^{n - p}$。考虑对这些可区别三角作用
$Rf_*$ 所得的相应可区别三角
$$
\xymatrix{
Rf_*\sigma_{\leq p}\Omega^\bullet \ar[d]_a \\
Rf_*\Omega^{p + 1}[-p - 1] \ar[d]_b \\
Rf_*\sigma_{\leq p + 1}\Omega^\bullet \ar[d]_c \\
Rf_*\sigma_{\leq p}\Omega^\bullet \ar[d]_d \\
Rf_*\Omega^{p + 1}[-p - 1]
}
\quad\quad
\xymatrix{
Rf_*\sigma_{\geq n - p}\Omega^\bullet \\
Rf_*\Omega^{n - p - 1}[-n + p + 1] \ar[u]_{a'} \\
Rf_*\sigma_{\geq n - p - 1}\Omega^\bullet \ar[u]_{b'} \\
Rf_*\sigma_{\geq n - p}\Omega^\bullet \ar[u]_{c'} \\
Rf_*\Omega^{n - p - 1}[-n + p + 1] \ar[u]_{d'}
}
$$
下文将证明各对 $(a, a')$、$(b, b')$、$(c, c')$ 和 $(d, d')$
与给定配对相容。这意味着我们得到一个从左侧可区别三角到如下可区别三角的映射：
对右侧可区别三角作用 $R\SheafHom(-, \mathcal{O}_S)$ 所得的可区别三角。
由归纳假设和引理 \ref{derham-lemma-duality-hodge}，上面在第一、第二、第四和第五行的
复形之间构造的配对是完美的；也就是说，对偶后它们确定同构。由导出范畴，
引理 \ref{derived-lemma-third-isomorphism-triangle}，中间一行复形之间的配对
也是完美的，正如所需。

\medskip\noindent
设 $e : K \to K'$ 和 $e' : M' \to M$ 是 $D(\mathcal{O}_S)$ 中对象的映射，
并设 $K \otimes_{\mathcal{O}_S}^\mathbf{L} M \to \mathcal{O}_S$ 和
$K' \otimes_{\mathcal{O}_S}^\mathbf{L} M' \to \mathcal{O}_S$ 是配对。
如果下图交换，就称这些配对相容：
$$
\xymatrix{
K' \otimes_{\mathcal{O}_S}^\mathbf{L} M' \ar[d] &
K \otimes_{\mathcal{O}_S}^\mathbf{L} M'
\ar[l]^{e \otimes 1} \ar[d]^{1 \otimes e'} \\
\mathcal{O}_S &
K \otimes_{\mathcal{O}_S}^\mathbf{L} M \ar[l]
}
$$
事实上，这意味着下图交换：
$$
\xymatrix{
K \ar[r] \ar[d]_e & R\SheafHom(M, \mathcal{O}_S) \ar[d]^{R\SheafHom(e', -)} \\
K' \ar[r] & R\SheafHom(M', \mathcal{O}_S)
}
$$
，所以足以满足我们的需要。

\medskip\noindent
先对 $(c, c')$ 证明这一点。这里只需注意有交换图
$$
\xymatrix{
\text{Tot}(\sigma_{\leq p} \Omega^\bullet \otimes_{f^{-1}\mathcal{O}_S}
\sigma_{\geq n - p} \Omega^\bullet) \ar[d]_{\gamma_p} &
\text{Tot}(\sigma_{\leq p + 1} \Omega^\bullet \otimes_{f^{-1}\mathcal{O}_S}
\sigma_{\geq n - p} \Omega^\bullet) \ar[l] \ar[d] \\
\Omega^\bullet &
\text{Tot}(\sigma_{\leq p + 1} \Omega^\bullet \otimes_{f^{-1}\mathcal{O}_S}
\sigma_{\geq n - p - 1} \Omega^\bullet) \ar[l]_-{\gamma_{p + 1}}
}
$$
由杯积的函子性，便得到所需图的交换性。

\medskip\noindent
类似地，对 $(b, b')$ 使用交换图
$$
\xymatrix{
\text{Tot}(\sigma_{\leq p + 1} \Omega^\bullet \otimes_{f^{-1}\mathcal{O}_S}
\sigma_{\geq n - p - 1} \Omega^\bullet) \ar[d]_{\gamma_{p + 1}} &
\text{Tot}(\Omega^{p + 1}[-p - 1] \otimes_{f^{-1}\mathcal{O}_S}
\sigma_{\geq n - p - 1} \Omega^\bullet) \ar[l] \ar[d] \\
\Omega^\bullet &
\Omega^{p + 1}[-p - 1]
\otimes_{f^{-1}\mathcal{O}_S}
\Omega^{n - p - 1}[-n + p + 1] \ar[l]_-\wedge
}
$$

\medskip\noindent
对 $(d, d')$ 和 $(a, a')$，使用交换图
$$
\xymatrix{
\Omega^{p + 1}[-p] \otimes_{f^{-1}\mathcal{O}_S}
\Omega^{n - p - 1}[-n + p] \ar[d] &
\text{Tot}(\sigma_{\leq p}\Omega^\bullet \otimes_{f^{-1}\mathcal{O}_S}
\Omega^{n - p - 1}[-n + p]) \ar[l] \ar[d] \\
\Omega^\bullet &
\text{Tot}(\sigma_{\leq p}\Omega^\bullet \otimes_{f^{-1}\mathcal{O}_S}
\sigma_{\geq n - p}\Omega^\bullet) \ar[l]
}
$$

\medskip\noindent
关于 $t$ 在预复合以 $H^0(X, \mathcal{O}_X)$ 中某个单位的乘法这个意义下
唯一的论证，从略。
\end{proof}
